% Mechanical Waves — Physics Upper Sixth — Cours signé OnBuch+
% Official MINESEC syllabus. Compiler avec : tectonic PHY02-mechanical-waves.tex
\newcommand{\DOCMATIERE}{Physics}
\newcommand{\DOCNIVEAU}{Upper Sixth}
\newcommand{\DOCMODULE}{Upper Sixth — Physics}
\newcommand{\DOCLECON}{2}
\newcommand{\DOCTITRE}{Mechanical Waves}
% OnBuch+ — préambule « cours signés » ENRICHI (compilé avec tectonic / XeLaTeX)
% Système visuel « Pop » d'OnBuch+ : fond crème (jamais blanc), bordures encre
% épaisses, ombres « dures » décalées, titres Archivo Black en capitales.
% Version MATHÉMATIQUES : graphes (pgfplots/tikz), boîtes pédagogiques
% (définition, propriété, méthode, attention, le savais-tu, exercice résolu,
% exercice, TP, bilan), page de garde et signature OnBuch+ ancrée en bas.
%
% Macros attendues dans main.tex AVANT \input{preamble} :
%   \DOCMATIERE  \DOCNIVEAU  \DOCMODULE  \DOCLECON (numéro)  \DOCTITRE
\documentclass[11pt]{article}

\usepackage{fontspec}
\usepackage[english]{babel}
\usepackage[a4paper,margin=2cm,top=3.4cm,bottom=2.8cm,headheight=1.4cm,footskip=1.3cm]{geometry}
\usepackage{amsmath,amssymb,amsfonts}
\usepackage{mathtools} % \xmapsto, \coloneqq... souvent utilisés par les modèles
\usepackage{cancel} % simplifications barrées (bilans de forces)
% \abs{z} (module, valeur absolue) et \norm{u} (norme), utilisés par les modèles
\providecommand{\abs}{}\let\abs\relax\DeclarePairedDelimiter{\abs}{\lvert}{\rvert}
\providecommand{\norm}{}\let\norm\relax\DeclarePairedDelimiter{\norm}{\lVert}{\rVert}
\usepackage[table]{xcolor} % [table] : fournit aussi \rowcolors (alternance de lignes)
\usepackage{pagecolor}
\usepackage{tikz}
\usetikzlibrary{arrows.meta,positioning,calc,shapes.geometric,shapes.misc,shapes.arrows,decorations.pathmorphing,decorations.pathreplacing,decorations.markings,patterns,shadows,mindmap,trees,fit,backgrounds,matrix,intersections,angles,quotes}
\usepackage{pgfplots}
\usepackage[version=4]{mhchem} % équations nucléaires \ce{^{235}_{92}U}
\usepackage[european,siunitx]{circuitikz} % schémas électriques
\pgfplotsset{compat=1.18}
\usepgfplotslibrary{fillbetween,groupplots} % aires entre courbes (calcul intégral)
\usepackage{enumitem}
\usepackage{fancyhdr}
% [most] charge toutes les bibliothèques tcolorbox (listings, minted, raster,
% documentation...) et gonfle inutilement la mémoire TeX sur un document long
% (~300 boîtes) : on ne charge que ce qui est réellement utilisé.
\usepackage[skins,breakable]{tcolorbox}
\usepackage{graphicx}
\usepackage{colortbl}
\usepackage{booktabs}
\usepackage{tabularx}
\usepackage{diagbox} % case d'en-tête diagonale des tableaux à double entrée
% Colonnes extensibles centrée / gauche / droite (C, L, R), souvent utilisées par les modèles
\newcolumntype{C}{>{\centering\arraybackslash}X}
\newcolumntype{L}{>{\raggedright\arraybackslash}X}
\newcolumntype{R}{>{\raggedleft\arraybackslash}X}
\usepackage{array}
\usepackage{multicol}
\usepackage{siunitx}
% parse-numbers=false : accepte aussi \SI{2,5 \times 10^{-3}}{...}
\sisetup{locale=UK,output-decimal-marker={.},group-minimum-digits=4,parse-numbers=false}
\DeclareSIUnit\ohm{\ensuremath{\symup{\Omega}}} % U+2126 absent de la police
\DeclareSIUnit\division{div} % réglages d'oscilloscope (ms/div, V/div)
\DeclareSIUnit\tr{tr} % tours (vitesse de rotation, ex. \SI{1500}{\tr\per\minute})
\usepackage{qrcode}
% \degree : commande fréquemment utilisée par les modèles pour un angle ou une
% température en dehors de \SI{}{} (non fournie par les paquets chargés).
\providecommand{\degree}{^{\circ}}
% \qed (fin de démonstration), utilisé par les modèles sans amsthm
\providecommand{\qed}{\hfill$\square$}
% \barre : double barre des valeurs interdites dans un tableau de variations (tkz-tab)
\providecommand{\barre}{\ensuremath{\|}}
% environnement proof (amsthm non chargé) utilisé par les modèles
\ifdefined\proof\else\newenvironment{proof}[1][Proof]{\par\noindent\textit{#1.}\ }{\qed\par}\fi
\usepackage{lastpage}
\usepackage{hyperref}
\hypersetup{hidelinks}

% ============================================================
% Palette « Pop » OnBuch+ — mêmes hex que la classe P (plus_ui.dart)
% ============================================================
\definecolor{poporange}{HTML}{F59321}
\definecolor{popdark}{HTML}{E07A0C}
\definecolor{popcream}{HTML}{FFF6E8}
\definecolor{popink}{HTML}{1C1714}
\definecolor{poppurple}{HTML}{7A5AE0}
\definecolor{popgreen}{HTML}{2E9E6A}
\definecolor{poppink}{HTML}{E0457B}
\definecolor{popblue}{HTML}{2F7DE1}
\definecolor{popgold}{HTML}{C98A12}
\definecolor{popmuted}{HTML}{8A7F76}
\definecolor{popline}{HTML}{EADFD2}
\definecolor{popgray}{HTML}{8A7F76}
\colorlet{popgrey}{popgray}
% Alias tolérants (le modèle nomme parfois les couleurs différemment)
\colorlet{popwhite}{white}
\colorlet{popblack}{popink}
\colorlet{popred}{poppink}
\colorlet{popyellow}{popgold}
% Teintes claires pour fonds de tableaux / zones
\colorlet{popblueL}{popblue!10!white}
\colorlet{poporangeL}{poporange!14!white}
\colorlet{popgreenL}{popgreen!12!white}
\colorlet{poppurpleL}{poppurple!12!white}
\colorlet{poppinkL}{poppink!12!white}
\colorlet{popgoldL}{popgold!14!white}

\pagecolor{popcream}

% ============================================================
% Typographie
% ============================================================
\defaultfontfeatures{Ligatures=TeX}
\setmainfont{PlusJakartaSans-Medium.ttf}[Path=../../../fonts/,BoldFont=PlusJakartaSans-Bold.ttf]
% Maths en sans-serif (Fira Math) avec chiffres et lettres droites de Plus
% Jakarta, pour que les formules se fondent dans le texte.
\usepackage[math-style=ISO,bold-style=ISO]{unicode-math}
\setmathfont{FiraMath-Regular.otf}[Path=../../../fonts/]
\setmathfont{PlusJakartaSans-Medium.ttf}[Path=../../../fonts/,range={up/num,up/{Latin,latin},\mathup}]
% (icomma retiré : décimales avec point en anglais)
% Caractères mathématiques Unicode absents des polices de texte (Plus Jakarta,
% Archivo Black) : rendus par leur équivalent mathématique, en texte comme en
% formule — sinon ils s'affichent en carré vide (ex. « Arithmétique dans ℤ »).
\usepackage{newunicodechar}
\newunicodechar{ℝ}{\ensuremath{\mathbb{R}}}
\newunicodechar{ℤ}{\ensuremath{\mathbb{Z}}}
\newunicodechar{ℂ}{\ensuremath{\mathbb{C}}}
\newunicodechar{ℕ}{\ensuremath{\mathbb{N}}}
\newunicodechar{ℚ}{\ensuremath{\mathbb{Q}}}
\newunicodechar{𝔻}{\ensuremath{\mathbb{D}}}
\newunicodechar{α}{\ensuremath{\alpha}}
\newunicodechar{β}{\ensuremath{\beta}}
\newunicodechar{γ}{\ensuremath{\gamma}}
\newunicodechar{δ}{\ensuremath{\delta}}
\newunicodechar{ε}{\ensuremath{\varepsilon}}
\newunicodechar{θ}{\ensuremath{\theta}}
\newunicodechar{λ}{\ensuremath{\lambda}}
\newunicodechar{μ}{\ensuremath{\mu}}
\newunicodechar{σ}{\ensuremath{\sigma}}
\newunicodechar{τ}{\ensuremath{\tau}}
\newunicodechar{φ}{\ensuremath{\varphi}}
\newunicodechar{ψ}{\ensuremath{\psi}}
\newunicodechar{ω}{\ensuremath{\omega}}
\newunicodechar{Σ}{\ensuremath{\Sigma}}
% majuscules produites par \MakeUppercase dans les titres (α-aminés -> Α)
\newunicodechar{Α}{\ensuremath{\alpha}}
\newunicodechar{Β}{\ensuremath{\beta}}
\newunicodechar{Γ}{\ensuremath{\Gamma}}
\newunicodechar{Δ}{\ensuremath{\Delta}}
\newunicodechar{Ε}{\ensuremath{\varepsilon}}
\newunicodechar{Θ}{\ensuremath{\Theta}}
\newunicodechar{Λ}{\ensuremath{\Lambda}}
\newunicodechar{Μ}{\ensuremath{\mu}}
\newunicodechar{Τ}{\ensuremath{\tau}}
\newunicodechar{Φ}{\ensuremath{\Phi}}
\newunicodechar{Ψ}{\ensuremath{\Psi}}
\newunicodechar{Ω}{\ensuremath{\Omega}}
\newunicodechar{↦}{\ensuremath{\mapsto}}
\newunicodechar{∈}{\ensuremath{\in}}
\newunicodechar{∉}{\ensuremath{\notin}}
\newunicodechar{∧}{\ensuremath{\wedge}}
\newunicodechar{⇔}{\ensuremath{\Leftrightarrow}}
\newunicodechar{⟺}{\ensuremath{\Longleftrightarrow}}
\newunicodechar{⇒}{\ensuremath{\Rightarrow}}
\newunicodechar{⟹}{\ensuremath{\Longrightarrow}}
\newunicodechar{∘}{\ensuremath{\circ}}
\newunicodechar{∪}{\ensuremath{\cup}}
\newunicodechar{∩}{\ensuremath{\cap}}
\newunicodechar{≡}{\ensuremath{\equiv}}
\newunicodechar{⊂}{\ensuremath{\subset}}
\newunicodechar{⊥}{\ensuremath{\perp}}
\newunicodechar{∃}{\ensuremath{\exists}}
\newunicodechar{∀}{\ensuremath{\forall}}
\newunicodechar{∛}{\ensuremath{\sqrt[3]{\,}}}
\newunicodechar{∜}{\ensuremath{\sqrt[4]{\,}}}
\newunicodechar{⃗}{\ensuremath{{}^{\to}}}
\newunicodechar{ⁿ}{\ifmmode^{n}\else\textsuperscript{\ensuremath{n}}\fi}
\newunicodechar{ˣ}{\ifmmode^{x}\else\textsuperscript{\ensuremath{x}}\fi}
\newunicodechar{ᵇ}{\ifmmode^{b}\else\textsuperscript{\ensuremath{b}}\fi}
\newunicodechar{ʳ}{\ifmmode^{r}\else\textsuperscript{\ensuremath{r}}\fi}
\newunicodechar{ᵘ}{\ifmmode^{u}\else\textsuperscript{\ensuremath{u}}\fi}
\newunicodechar{ʸ}{\ifmmode^{y}\else\textsuperscript{\ensuremath{y}}\fi}
\newunicodechar{ᵃ}{\ifmmode^{a}\else\textsuperscript{\ensuremath{a}}\fi}
\newunicodechar{ᵉ}{\ifmmode^{e}\else\textsuperscript{\ensuremath{e}}\fi}
\newunicodechar{ᵏ}{\ifmmode^{k}\else\textsuperscript{\ensuremath{k}}\fi}
\newunicodechar{ᵖ}{\ifmmode^{p}\else\textsuperscript{\ensuremath{p}}\fi}
\newunicodechar{ᴺ}{\ifmmode^{N}\else\textsuperscript{\ensuremath{N}}\fi}
\newunicodechar{ᶜ}{\ifmmode^{c}\else\textsuperscript{\ensuremath{c}}\fi}
\newunicodechar{ʰ}{\ifmmode^{h}\else\textsuperscript{\ensuremath{h}}\fi}
\newunicodechar{ᵠ}{\ifmmode^{q}\else\textsuperscript{\ensuremath{q}}\fi}
\newunicodechar{ₙ}{\ifmmode_{n}\else\textsubscript{\ensuremath{n}}\fi}
\newunicodechar{ₐ}{\ifmmode_{a}\else\textsubscript{\ensuremath{a}}\fi}
\newunicodechar{ᵢ}{\ifmmode_{i}\else\textsubscript{\ensuremath{i}}\fi}
\newunicodechar{ᵣ}{\ifmmode_{r}\else\textsubscript{\ensuremath{r}}\fi}
\newunicodechar{ₖ}{\ifmmode_{k}\else\textsubscript{\ensuremath{k}}\fi}
\newunicodechar{ᵦ}{\ifmmode_{\beta}\else\textsubscript{\ensuremath{\beta}}\fi}
\newunicodechar{ₓ}{\ifmmode_{x}\else\textsubscript{\ensuremath{x}}\fi}
\newfontfamily\popdisplay{ArchivoBlack-Regular.ttf}[Path=../../../fonts/]
\newfontfamily\poplogo{SpaceGrotesk-Bold.ttf}[Path=../../../fonts/]
\newfontfamily\popsemi{PlusJakartaSans-SemiBold.ttf}[Path=../../../fonts/]
\newfontfamily\popxbold{PlusJakartaSans-ExtraBold.ttf}[Path=../../../fonts/]
\newfontfamily\popxboldls{PlusJakartaSans-ExtraBold.ttf}[Path=../../../fonts/,LetterSpace=3.0]

\setlist[enumerate]{leftmargin=1.7em,itemsep=5pt,topsep=4pt}
\setlist[itemize]{leftmargin=1.7em,itemsep=4pt,topsep=4pt,label={\color{poporange}\textbullet}}
\setlength{\parindent}{0pt}
\setlength{\parskip}{5pt}
\renewcommand{\arraystretch}{1.25}

% pgfplots : style Pop par défaut
\pgfplotsset{
  every axis/.append style={axis line style={popink,line width=1pt},tick style={popink},
    grid style={popline,line width=0.5pt},label style={font=\small},tick label style={font=\footnotesize}},
  popcurve/.style={poporange,line width=1.8pt,no markers,smooth},
}
% Même style disponible dans un \draw TikZ simple (hors pgfplots)
\tikzset{popcurve/.style={poporange,line width=1.8pt}}
\tikzset{popgrid/.style={popline,very thin}}
\colorlet{popcurve}{poporange} % le nom du style est parfois utilisé comme couleur

% ============================================================
% Pill / squiggle
% ============================================================
\newcommand{\poppill}[3]{%
  \tikz[baseline=-0.55ex]\node[draw=popink,line width=1.2pt,rounded corners=2.6pt,%
    fill=#2,inner xsep=6.5pt,inner ysep=3pt]{%
    \popxboldls\fontsize{7}{7}\selectfont\color{#3}\MakeUppercase{#1}};%
}
\newcommand{\popsquiggle}[1]{%
  \tikz[baseline=-0.4ex]{
    \draw[#1,line width=2.6pt,line cap=round]
      plot [smooth] coordinates {(0,0.09) (0.34,0.01) (0.68,0.21) (1.02,0.01) (1.36,0.10)};
  }%
}
% Mot-clé surligné (marqueur orange sous le texte)
\newcommand{\cle}[1]{{\popxbold\color{popdark}#1}}

% ============================================================
% En-tête / pied
% ============================================================
\pagestyle{fancy}
\fancyhf{}
\renewcommand{\headrulewidth}{0pt}
\renewcommand{\footrulewidth}{0pt}
\lhead{\raisebox{-0.15ex}{\poplogo\fontsize{13}{13}\selectfont\color{popink}OnBuch\color{poporange}+}}
\rhead{\poppill{\DOCMATIERE\ \textbullet\ \DOCNIVEAU\ \textbullet\ Chapter \DOCLECON}{white}{popink}}
\lfoot{\popxbold\fontsize{7.6}{7.6}\selectfont\color{popmuted}\MakeUppercase{Signed course OnBuch+}\ \textbullet\ {\popsemi\DOCTITRE}}
\rfoot{\tikz[baseline=-0.6ex]\node[rounded corners=8.5pt,fill=popink,minimum height=17pt,minimum width=17pt,inner xsep=5pt,inner sep=0]{\popxbold\fontsize{8}{8}\selectfont\color{popcream}\thepage\,/\,\pageref*{LastPage}};}

% ============================================================
% Titres
% ============================================================
\newcommand{\chaptitre}[2]{%
  \begin{tcolorbox}[enhanced,colback=poporange,colframe=popink,boxrule=2.6pt,arc=11pt,
      drop shadow=popink,left=15pt,right=15pt,top=12pt,bottom=11pt,before skip=3pt,after skip=18pt]
    \poppill{#2}{popink}{popcream}\par\vspace{9pt}
    {\raggedright\hyphenpenalty=10000\exhyphenpenalty=10000\popdisplay\fontsize{22}{24}\selectfont\color{popcream}\MakeUppercase{#1}\par}
  \end{tcolorbox}
}
% Section numérotée : pastille orange avec le numéro + titre Archivo
\newcounter{popsec}
\newcommand{\coursec}[1]{%
  \stepcounter{popsec}\setcounter{popsub}{0}%
  \par\vspace{16pt}\noindent
  \begin{minipage}{\linewidth}
  \tikz[baseline=-0.7ex]\node[draw=popink,line width=1.6pt,fill=poporange,rounded corners=4pt,minimum size=22pt,inner sep=0]{\popdisplay\fontsize{12}{12}\selectfont\color{popcream}\thepopsec};%
  \hspace{8pt}{\popdisplay\fontsize{15}{17}\selectfont\color{popink}\MakeUppercase{#1}}\par
  \vspace{4pt}\noindent{\color{poporange}\rule{36pt}{3.6pt}}
  \end{minipage}\par\nopagebreak\vspace{8pt}
}
\newcounter{popsub}
\newcommand{\courssub}[1]{%
  \stepcounter{popsub}%
  \par\vspace{9pt}\noindent{\popxbold\fontsize{11.5}{13}\selectfont\color{popink}{\color{poporange}\thepopsec.\thepopsub}\hspace{6pt}#1}\par\nopagebreak\vspace{4pt}
}

% ============================================================
% Boîtes pédagogiques — même recette : fond blanc, bordure épaisse colorée,
% ombre dure encre, badge plein. Titre optionnel [#1].
% ============================================================
\tcbset{popbase/.style={breakable,enhanced,colback=white,boxrule=2.3pt,arc=9pt,
  shadow={3pt}{-3pt}{0pt}{popink},left=14pt,right=14pt,top=11pt,bottom=11pt,before skip=11pt,after skip=15pt,
  fontupper=\popsemi\fontsize{10.6}{14.5}\selectfont\color{popink},
  fontlower=\popsemi\fontsize{10.6}{14.5}\selectfont\color{popink}}}
% Coche dessinée (le glyphe ✓ n'existe pas dans les polices du document)
\newcommand{\popcheck}[1]{\tikz[baseline=-0.5ex]\draw[#1,line width=1.6pt,line cap=round,line join=round](-0.13,0.0)--(-0.04,-0.09)--(0.13,0.1);}
\providecommand{\coche}{\popcheck{popgreen}}
\providecommand{\resultat}[1]{\par\smallskip\noindent\fcolorbox{popgreen}{popgreenL}{\parbox{\dimexpr\linewidth-2\fboxsep-2\fboxrule\relax}{\textbf{Result.} #1}}\par\smallskip}
\newcommand{\popboxhead}[3]{\poppill{#1}{#2}{white}\ifx\relax#3\relax\else\hspace{6pt}{\popxbold\fontsize{10.6}{12}\selectfont\color{popink}#3}\fi\par\vspace{7pt}}

\newtcolorbox{aretenir}[1][]{popbase,colframe=popgreen,before upper={\popboxhead{Key points}{popgreen}{#1}}}
\newtcolorbox{exemplebox}[1][]{popbase,colframe=poporange,before upper={\popboxhead{Exemple}{poporange}{#1}}}
\newtcolorbox{definition}[1][]{popbase,colframe=popblue,before upper={\popboxhead{Definition}{popblue}{#1}}}
\newtcolorbox{propriete}[1][]{popbase,colframe=poppurple,before upper={\popboxhead{Property}{poppurple}{#1}}}
\newtcolorbox{methode}[1][]{popbase,colframe=popgold,before upper={\popboxhead{Method}{popgold}{#1}}}
\newtcolorbox{attention}[1][]{popbase,colframe=poppink,before upper={\popboxhead{Warning}{poppink}{#1}}}
\newtcolorbox{experience}[1][]{popbase,colframe=popblue,colback=popblueL,before upper={\popboxhead{Experiment}{popblue}{#1}}}
% « Le savais-tu ? » : carte violette pleine (contexte, culture, Cameroun)
\newtcolorbox{savaistu}[1][]{popbase,colback=poppurple,colframe=popink,
  fontupper=\popsemi\fontsize{10.4}{14.2}\selectfont\color{white},
  before upper={\poppill{Did you know?}{white}{poppurple}\ifx\relax#1\relax\else\hspace{6pt}{\popxbold\color{white}#1}\fi\par\vspace{7pt}}}
% Exercice résolu : énoncé en haut, \tcblower puis la solution sur fond crème
\newcounter{popexo}\newcounter{popexr}
\newtcolorbox{exoresolu}[1][]{popbase,colframe=poporange,colbacklower=popcream,
  segmentation style={popink,line width=1.2pt,dashed},
  before upper={\stepcounter{popexr}\popboxhead{Worked example \thepopexr}{poporange}{#1}},
  before lower={\poppill{Solution}{popink}{popcream}\par\vspace{6pt}}}
% Exercice (non résolu en place) — la correction est dans la partie Corrigés
\newtcolorbox{exercice}[1][]{popbase,colframe=popink,
  before upper={\stepcounter{popexo}\popboxhead{Exercise \thepopexo}{popink}{#1}}}
% Corrigé
\newtcolorbox{corrige}[1][]{popbase,colframe=popgreen,colback=popgreenL,
  before upper={\popboxhead{Solution}{popgreen}{#1}}}
% Objectifs / prérequis / bilan
\newtcolorbox{objectifs}{popbase,colframe=popink,colback=poporangeL,
  before upper={\popboxhead{Chapter objectives}{popink}{}}}
\newtcolorbox{prerequis}{popbase,colframe=popink,colback=popblueL,
  before upper={\popboxhead{Prerequisites}{popblue}{}}}
\newtcolorbox{bilan}{popbase,colframe=popink,colback=white,boxrule=2.8pt,
  before upper={\popboxhead{Fiche bilan}{poporange}{L'essentiel à connaître pour l'examen}}}
% Figure encadrée avec légende
\newtcolorbox{popfigure}[1][]{enhanced,breakable=false,colback=white,colframe=popline,boxrule=1.4pt,arc=8pt,
  left=10pt,right=10pt,top=10pt,bottom=8pt,before skip=10pt,after skip=12pt,halign=center,
  lower separated=false,halign lower=center,
  fontlower=\popsemi\fontsize{9}{11.5}\selectfont\color{popmuted},
  #1}
\newcommand{\legende}[1]{\par\vspace{6pt}{\popsemi\fontsize{9}{11.5}\selectfont\color{popmuted}#1}}

% ============================================================
% Signature OnBuch+
% ============================================================
\newcommand{\poplogoinline}[1]{{\poplogo\fontsize{#1}{#1}\selectfont\color{popink}OnBuch\color{poporange}+}}
\newcommand{\signatureonbuch}{%
  \begin{tcolorbox}[enhanced,colback=popink,colframe=popink,boxrule=2.2pt,arc=10pt,
      drop shadow=poporange,left=14pt,right=14pt,top=10pt,bottom=10pt,before skip=0pt,after skip=0pt]
    \tikz[baseline=-0.7ex]\node[circle,fill=popgreen,draw=popcream,line width=1.3pt,minimum size=22pt,inner sep=0]{\popcheck{white}};%
    \hspace{9pt}\begin{minipage}[c]{0.62\linewidth}
      {\popxbold\fontsize{12}{13}\selectfont\color{popcream}Signed\ \textbullet\ {\poplogo OnBuch\color{poporange}+}}\par\vspace{2pt}
      {\raggedright\popsemi\fontsize{8}{10.5}\selectfont\color{popline}\MakeUppercase{OnBuch+ teaching team}\\\MakeUppercase{Follows the official MINESEC syllabus}\par}
    \end{minipage}\hfill
    {\poplogo\fontsize{22}{22}\selectfont\color{popcream}OnBuch\color{poporange}+}
  \end{tcolorbox}%
}

% ============================================================
% Page de garde
% #1 titre · #2 matière · #3 niveau · #4 module · #5 n° leçon · #6 durée ·
% #7 description / objectifs (texte ou itemize)
% ============================================================
\newcommand{\pagegarde}[7]{%
  \thispagestyle{empty}
  \newgeometry{a4paper,margin=1.7cm,top=1.6cm,bottom=1.4cm}%
  \begin{tikzpicture}[remember picture,overlay]
    \fill[poporange!18] ([xshift=-3.2cm,yshift=-3.6cm]current page.north east) circle (4.6cm);
    \fill[poppurple!14] ([xshift=2.4cm,yshift=7.6cm]current page.south west) circle (3.8cm);
  \end{tikzpicture}%
  {\poplogo\fontsize{30}{30}\selectfont\color{popink}OnBuch\color{poporange}+}\hfill
  \raisebox{0.6ex}{\poppill{Signed course \textbullet\ Premium}{popink}{popcream}}\par
  \vspace{1pt}\popsquiggle{poppurple}\par
  \vspace{8pt}
  \begin{tcolorbox}[enhanced,colback=poporange,colframe=popink,boxrule=2.8pt,arc=13pt,
      drop shadow=popink,left=18pt,right=18pt,top=12pt,bottom=13pt,after skip=10pt]
    \poppill{#2 \textbullet\ #3}{popink}{popcream}\hspace{5pt}\poppill{#4}{white}{popink}\par\vspace{8pt}
    {\popdisplay\fontsize{14}{14}\selectfont\color{popink}CHAPTER #5}\par\vspace{4pt}
    {\raggedright\hyphenpenalty=10000\exhyphenpenalty=10000\popdisplay\fontsize{25}{27}\selectfont\color{popcream}\MakeUppercase{#1}\par}
    \vspace{8pt}
    {\popxbold\fontsize{9.5}{9.5}\selectfont\color{popink}\MakeUppercase{Suggested time: #6}}
  \end{tcolorbox}
  \begin{tcolorbox}[enhanced,colback=white,colframe=popink,boxrule=2.2pt,arc=10pt,
      drop shadow=popink,left=16pt,right=16pt,top=10pt,bottom=10pt,after skip=8pt]
    \poppill{In this chapter}{poppurple}{white}\par\vspace{5pt}
    \popsemi\fontsize{9.3}{12.6}\selectfont\color{popink}#7
  \end{tcolorbox}
  \noindent\begin{minipage}[t]{0.487\linewidth}
    \begin{tcolorbox}[enhanced,colback=popgreen,colframe=popink,boxrule=2.2pt,arc=10pt,
        drop shadow=popink,left=11pt,right=11pt,top=10pt,bottom=10pt,height=3.1cm,valign=center]
      \tcbox[colback=white,colframe=popink,boxrule=1.3pt,arc=5pt,boxsep=0pt,nobeforeafter,
        left=3pt,right=3pt,top=3pt,bottom=3pt,valign=center]{\qrcode[height=1.9cm]{https://play.google.com/store/apps/details?id=com.luvvix.onbuch&hl=en}}%
      \hspace{8pt}\begin{minipage}[c]{0.5\linewidth}
        {\popdisplay\fontsize{11}{12.5}\selectfont\color{white}\MakeUppercase{Download OnBuch}}\par\vspace{4pt}
        {\raggedright\popsemi\fontsize{8.4}{11}\selectfont\color{white}Free \textbullet\ Google Play\\AI tutor, past papers, results\par}
      \end{minipage}
    \end{tcolorbox}
  \end{minipage}\hfill
  \begin{minipage}[t]{0.487\linewidth}
    \begin{tcolorbox}[enhanced,colback=poppurple,colframe=popink,boxrule=2.2pt,arc=10pt,
        drop shadow=popink,left=11pt,right=11pt,top=10pt,bottom=10pt,height=3.1cm,valign=center]
      \tikz[baseline=-0.7ex]\node[circle,fill=white,draw=popink,line width=1.3pt,minimum size=1.6cm,inner sep=0]{\popdisplay\fontsize{24}{24}\selectfont\color{poppurple}+};%
      \hspace{8pt}\begin{minipage}[c]{0.6\linewidth}
        {\popdisplay\fontsize{11}{12.5}\selectfont\color{white}\MakeUppercase{Go OnBuch+}}\par\vspace{4pt}
        {\raggedright\popsemi\fontsize{8.4}{11}\selectfont\color{white}All signed courses, unlimited AI tutor, PDF sheets — OnBuch+ tab in the app\par}
      \end{minipage}
    \end{tcolorbox}
  \end{minipage}
  \vspace{8pt}
  \popcontenu
  \vfill
  \signatureonbuch
  \restoregeometry
  \newpage
}

% Rangée « Ce cours contient » (page de garde)
\newcommand{\poptile}[3]{% #1 couleur #2 gros texte #3 légende
  \begin{tcolorbox}[enhanced,colback=#1,colframe=popink,boxrule=2pt,arc=9pt,shadow={2.5pt}{-2.5pt}{0pt}{popink},
      left=6pt,right=6pt,top=9pt,bottom=9pt,halign=center,valign=center,height=1.75cm,width=\linewidth,nobeforeafter]
    {\popdisplay\fontsize{12.5}{13}\selectfont\color{white}\MakeUppercase{#2}}\par\vspace{3pt}
    {\centering\popsemi\fontsize{7.8}{9.5}\selectfont\color{white}#3\par}
  \end{tcolorbox}}
\newcommand{\popcontenu}{%
  \par\noindent{\popxboldls\fontsize{7.5}{7.5}\selectfont\color{popmuted}THIS SIGNED COURSE CONTAINS}\par\vspace{7pt}
  \noindent\begin{minipage}[t]{0.235\linewidth}\poptile{popblue}{Course}{complete, illustrated, step by step}\end{minipage}\hfill
  \begin{minipage}[t]{0.235\linewidth}\poptile{poporange}{Methods}{and worked examples}\end{minipage}\hfill
  \begin{minipage}[t]{0.235\linewidth}\poptile{poppink}{Exercises}{exam style + solutions}\end{minipage}\hfill
  \begin{minipage}[t]{0.235\linewidth}\poptile{popgreen}{Summary sheet}{the essentials to remember}\end{minipage}\par}

% Sommaire
\newtcolorbox{sommaire}{enhanced,colback=white,colframe=popink,boxrule=2.2pt,arc=10pt,
  drop shadow=popink,left=18pt,right=18pt,top=13pt,bottom=13pt,before skip=6pt,after skip=20pt,
  before upper={\poppill{Contents}{poppurple}{white}\par\vspace{9pt}\popsemi\fontsize{10.6}{14.5}\selectfont\color{popink}}}

% Signature de fin de cours — poussée tout en bas de la dernière page
\newcommand{\findecours}{%
  \par\vfill\nopagebreak
  \begin{center}\popsquiggle{poporange}\end{center}\vspace{4pt}
  \signatureonbuch
}
\AtBeginDocument{\def\llbracket{\mathopen{[\mkern-3.5mu[}}\def\rrbracket{\mathclose{]\mkern-3.5mu]}}} % intervalles d'entiers (glyphe ⟦ absent de la police)
% Glyphes absents de Fira Math : redéfinis à partir de glyphes disponibles
\AtBeginDocument{%
  \def\setminus{\mathbin{\backslash}}\let\smallsetminus\setminus
  \def\vdots{\mathord{\vbox{\baselineskip3.2pt\lineskiplimit0pt\kern4pt\hbox{.}\hbox{.}\hbox{.}}}}%
  \def\ddots{\mathinner{\mkern1mu\raise6.5pt\hbox{.}\mkern2mu\raise3.6pt\hbox{.}\mkern2mu\raise0.7pt\hbox{.}\mkern1mu}}%
  \def\triangle{\mathord{\scalebox{0.9}{$\symup{\Delta}$}}}%
  \def\nabla{\mathord{\raisebox{\depth}{\scalebox{1}[-1]{$\symup{\Delta}$}}}}%
  \def\checkmark{\popcheck{popdark}}%
  \def\star{\mathbin{\ast}}\def\bigstar{\mathord{\ast}}%
  \def\diamond{\mathbin{\scalebox{0.6}{$\Diamond$}}}%
  \def\bigcup{\mathop{\mathchoice{\raisebox{-0.3em}{\scalebox{1.7}{$\cup$}}}{\scalebox{1.25}{$\cup$}}{\cup}{\cup}}\displaylimits}%
  \def\bigcap{\mathop{\mathchoice{\raisebox{-0.3em}{\scalebox{1.7}{$\cap$}}}{\scalebox{1.25}{$\cap$}}{\cap}{\cap}}\displaylimits}%
  \def\lesseqqgtr{\mathrel{\lessgtr}}\def\bigoplus{\mathop{\oplus}\displaylimits}%
}
\providecommand{\ssi}{if and only if} % abréviation fréquente
\AtBeginDocument{\providecommand{\wideparen}{\overparen}} % arc de cercle

%% FIGFIX-BEGIN v5 — correctifs globaux des figures (docs/onbuch-generation/tools/figfix)
\usepackage{adjustbox}\usepackage{etoolbox}\tcbuselibrary{raster}
% toute figure TikZ/pgfplots plus large que la ligne est réduite à la largeur disponible
\BeforeBeginEnvironment{tikzpicture}{\begin{adjustbox}{max width=\linewidth}}
\AfterEndEnvironment{tikzpicture}{\end{adjustbox}}
% légendes pgfplots lisibles par-dessus les courbes
\pgfplotsset{every axis/.append style={legend style={fill=white,fill opacity=0.92,text opacity=1,draw=black!15,inner sep=3pt}}}
% étiquettes d'arêtes (syntaxe edge["..."]) avec fond blanc
\tikzset{every edge quotes/.append style={fill=white,inner sep=1.2pt}}
%% FIGFIX-END
\begin{document}

\pagegarde{Mechanical Waves}{Physics}{Upper Sixth}{Upper Sixth — Physics}{2}{6 periods}{This chapter introduces mechanical waves: their nature, classification and mathematical description as progressive waves. It then explores the great wave phenomena — reflection, refraction, diffraction, interference, the Doppler effect and stationary waves — which lie at the heart of the GCE Advanced Level Physics examination.
\vspace{4pt}
\begin{multicols}{2}\small
\begin{itemize}
\item By the end of this chapter I can define a wave and classify waves according to their mode of propagation (transverse/longitudinal) and their need for a material medium, giving specific examples.
\item By the end of this chapter I can state and use the characteristics of a wave: period, frequency, wavelength, amplitude and speed, with v = fλ.
\item By the end of this chapter I can write and interpret the equation of a progressive wave y(x,t) and extract amplitude, wavelength, period and wave speed from it.
\item By the end of this chapter I can draw and interpret the graphical representations of a wave: y against x at fixed t, and y against t at fixed x.
\item By the end of this chapter I can describe and explain reflection, refraction, diffraction and interference of waves, including the single-slit diffraction pattern.
\item By the end of this chapter I can analyse double-slit and multiple-slit interference and use the fringe spacing to measure a wavelength.
\end{itemize}\end{multicols}}

\begin{sommaire}
\begin{multicols}{2}
\begin{enumerate}
\item Nature, Classification and Characteristics of Waves
\item The Progressive Wave: Equation and Graphical Representation
\item Wave Phenomena I: Reflection, Refraction, Diffraction and the Single Slit
\item Wave Phenomena II: Interference — Double and Multiple Slits, Measurement of Wavelength
\item Doppler Effect and Stationary Waves
\item Integration activity
\item Methods and worked examples
\item Exercises
\item Solutions to the exercises
\item Summary sheet
\end{enumerate}
\end{multicols}
\end{sommaire}

\begin{objectifs}
\begin{itemize}
\item By the end of this chapter I can define a wave and classify waves according to their mode of propagation (transverse/longitudinal) and their need for a material medium, giving specific examples.
\item By the end of this chapter I can state and use the characteristics of a wave: period, frequency, wavelength, amplitude and speed, with v = fλ.
\item By the end of this chapter I can write and interpret the equation of a progressive wave y(x,t) and extract amplitude, wavelength, period and wave speed from it.
\item By the end of this chapter I can draw and interpret the graphical representations of a wave: y against x at fixed t, and y against t at fixed x.
\item By the end of this chapter I can describe and explain reflection, refraction, diffraction and interference of waves, including the single-slit diffraction pattern.
\item By the end of this chapter I can analyse double-slit and multiple-slit interference and use the fringe spacing to measure a wavelength.
\item By the end of this chapter I can apply the Doppler effect formulae for sound in air for a moving source/stationary observer and a moving observer/stationary source.
\item By the end of this chapter I can explain the formation of stationary waves, identify nodes and antinodes, and relate string length to wavelength and frequency.
\end{itemize}
\end{objectifs}

\begin{prerequis}
\begin{itemize}
\item Simple harmonic motion: amplitude, period, frequency, angular frequency ω = 2πf (Lower Sixth oscillations)
\item Trigonometric functions sine and cosine, phase angles in radians
\item Graph plotting and interpretation of gradients and intercepts
\item Basic properties of sound and light from everyday experience and Ordinary Level physics
\item Algebraic manipulation of formulae and use of SI units
\end{itemize}
\end{prerequis}

\newpage

\coursec{Nature, Classification and Characteristics of Waves}

During a football match at the Ahmadou Ahidjo Stadium in Yaoundé, a spectator at the far end of the stands sees the striker kick the ball a split second before hearing the sound of the kick. In the village, a fisherman on Lake Nyos watches circular ripples spread from where his paddle struck the water, while a church bell heard from behind a hill still reaches the faithful though its source is hidden. And every day, the siren of a passing ambulance in Douala sounds higher-pitched as it approaches and lower as it moves away. All these everyday observations share one explanation: the physics of mechanical waves. This chapter gives you the tools to describe, represent and predict all of them.

\courssub{What Is a Wave?}

\begin{definition}[Wave]
A \cle{wave} is a propagating disturbance that transfers \textbf{energy} (and information) from one point to another \textbf{without net transfer of matter}. The particles of the medium oscillate about their equilibrium positions; they do not travel with the wave.
\end{definition}

\begin{popfigure}
\begin{tikzpicture}[scale=1.1]
% Medium particles at equilibrium
\foreach \x in {0,1,2,3,4,5,6,7,8,9} {
    \fill[popblue!60] (\x,0) circle (0.12);
    \draw[popmuted,->,thick] (\x,0) -- (\x,0.8) node[above,font=\tiny] {$y$};
    \draw[popmuted,->,thick] (\x,0) -- (\x+0.8,0) node[right,font=\tiny] {$x$};
}
% Wave profile
\draw[popink,thick,domain=0:9,smooth,variable=\x] plot ({\x},{0.6*sin(360*\x/2)});
\draw[poporange,thick,<->] (0,-1.2) -- (2,-1.2) node[midway,below] {$\lambda$};
\draw[poporange,thick,<->] (2,0) -- (2,0.6) node[midway,right] {$A$};
\node[popdark] at (4.5,-1.8) {Particles oscillate vertically (transverse)};
\node[popdark] at (4.5,-2.2) {Wave travels horizontally $\rightarrow$};
\end{tikzpicture}
\legende{A transverse wave on a rope: each particle moves up and down while the wave pattern moves to the right.}
\end{popfigure}

\courssub{Classification of Waves}

Waves are classified according to two independent criteria: the \textbf{need for a material medium} and the \textbf{direction of particle vibration} relative to the direction of propagation.

\begin{popfigure}
\begin{tikzpicture}[scale=0.9, every node/.style={font=\small}]
% Root
\node[draw,rounded corners,fill=popblueL,thick] (root) at (0,0) {Waves};
% First level: by medium
\node[draw,rounded corners,fill=popgreenL,thick, align=center] (mech) at (-5,-2){Mechanical Waves\\ (need a medium)};
\node[draw,rounded corners,fill=poppink,thick, align=center] (em) at (5,-2){Electromagnetic Waves\\ (no medium needed)};
\draw[thick,->] (root) -- (mech);
\draw[thick,->] (root) -- (em);
% Second level for mechanical: by mode
\node[draw,rounded corners,fill=popgreenL!80!white, align=center] (trans) at (-8,-4){Transverse\\ (e.g. rope, water ripples)};
\node[draw,rounded corners,fill=popgreenL!80!white, align=center] (long) at (-2,-4){Longitudinal\\ (e.g. sound in air, slinky)};
\draw[thick,->] (mech) -- (trans);
\draw[thick,->] (mech) -- (long);
% Second level for EM: all transverse
\node[draw,rounded corners,fill=poppink!80!white, align=center] (emtrans) at (5,-4){Transverse only\\ (light, radio, X-rays, $\gamma$)};
\draw[thick,->] (em) -- (emtrans);
% Examples
\node[align=center,font=\tiny] at (-8,-5) {Rope wave\\ Water surface\\ Seismic S-waves};
\node[align=center,font=\tiny] at (-2,-5) {Sound in air/gas/liquid\\ Seismic P-waves\\ Slinky spring};
\node[align=center,font=\tiny] at (5,-5) {Light, radio, microwave\\ Infrared, UV, X-rays, $\gamma$};
\end{tikzpicture}
\legende{Classification tree of waves. Mechanical waves require a material medium; electromagnetic waves do not. Mechanical waves can be transverse or longitudinal; electromagnetic waves are exclusively transverse.}
\end{popfigure}

\courssub{Mechanical Waves vs Electromagnetic Waves}

\begin{itemize}
\item \textbf{Mechanical waves} require a material medium (solid, liquid or gas) to propagate. The disturbance is a mechanical deformation (compression, shear, displacement) that travels through the medium. Examples: sound waves, waves on a stretched rope, water surface waves, seismic waves.
\item \textbf{Electromagnetic waves} do \emph{not} require a medium; they propagate through vacuum at speed $c = 3.00\times10^8\ \mathrm{m\,s^{-1}}$. They consist of oscillating electric and magnetic fields perpendicular to each other and to the direction of propagation. Examples: visible light, radio waves, microwaves, X-rays, $\gamma$-rays.
\end{itemize}

\begin{attention}[GCE Trap: Sound in Vacuum]
A classic GCE question: ``Can sound travel on the Moon?'' \textbf{No.} The Moon has no atmosphere (vacuum), so there is no medium for sound waves. Light, however, travels perfectly well in vacuum — that is how sunlight reaches Earth.
\end{attention}

\courssub{Transverse and Longitudinal Mechanical Waves}

\begin{propriete}[Transverse Wave]
In a \cle{transverse wave}, the displacement of the medium's particles is \textbf{perpendicular} to the direction of wave propagation.
\begin{itemize}
\item Examples: wave on a stretched rope, ripples on water surface, seismic S-waves (shear waves).
\item Can propagate in solids and on liquid surfaces (surface waves), but \textbf{not in bulk fluids} (gases/liquids) because fluids cannot sustain shear stress.
\end{itemize}
\end{propriete}

\begin{propriete}[Longitudinal Wave]
In a \cle{longitudinal wave}, the displacement of the medium's particles is \textbf{parallel} to the direction of wave propagation. The wave consists of alternating \cle{compressions} (regions of high pressure/density) and \cle{rarefactions} (regions of low pressure/density).
\begin{itemize}
\item Examples: sound waves in air, gases, liquids and solids; seismic P-waves (primary/pressure waves); waves along a slinky spring.
\item Can propagate in solids, liquids and gases.
\end{itemize}
\end{propriete}

\begin{popfigure}
\begin{tikzpicture}[scale=1.1]
% Transverse wave on a rope
\begin{scope}[yshift=3cm]
\draw[popblue,thick] (-0.5,0) -- (8.5,0); % equilibrium line
\draw[popink,thick,domain=0:8,smooth,variable=\x] plot ({\x},{0.5*sin(360*\x/2)});
\draw[poporange,<->,thick] (0,0) -- (0,0.5) node[midway,left] {$A$};
\draw[poporange,<->,thick] (0,-0.8) -- (2,-0.8) node[midway,below] {$\lambda$};
\draw[popdark,->,thick] (4,1.2) -- (4,0.5) node[above] {vibration $\uparrow\downarrow$};
\draw[popdark,->,thick] (4,-1.2) -- (6,-1.2) node[below] {propagation $\rightarrow$};
\node[popdark] at (4,-1.6) {Transverse wave on a rope};
\end{scope}

% Longitudinal wave on a slinky
\begin{scope}[yshift=-3cm]
% Equilibrium positions of coils
\foreach \x in {0,1,2,3,4,5,6,7,8} {
    \draw[popmuted,thin] (\x,-0.4) -- (\x,0.4);
}
% Compressions and rarefactions represented by displaced vertical lines
\foreach \x in {0,0.5,1,1.5,2,2.5,3,3.5,4,4.5,5,5.5,6,6.5,7,7.5,8} {
    \pgfmathsetmacro{\disp}{0.3*sin(360*\x/2)};
    \draw[popblue,thick] (\x+\disp,-0.3) -- (\x+\disp,0.3);
}
% Labels
\node[popdark,align=center] at (1,0.8) {Compression};
\node[popdark,align=center] at (3,-0.8) {Rarefaction};
\draw[poporange,<->,thick] (0,-1.2) -- (2,-1.2) node[midway,below] {$\lambda$};
\draw[popdark,->,thick] (4,1.2) -- (4,0.5) node[above] {vibration $\leftrightarrow$};
\draw[popdark,->,thick] (4,-1.5) -- (6,-1.5) node[below] {propagation $\rightarrow$};
\node[popdark] at (4,-2) {Longitudinal wave on a slinky};
\end{scope}
\end{tikzpicture}
\legende{Comparison of transverse (top) and longitudinal (bottom) mechanical waves. In the longitudinal wave, the horizontal displacement of coils creates compressions (close together) and rarefactions (far apart).}
\end{popfigure}

\begin{experience}[Observing Transverse and Longitudinal Waves]
\textbf{Apparatus:} Long stretched rope (or Slinky spring), ripple tank with motor, signal generator, strobe light.

\textbf{Procedure:}
\begin{enumerate}
\item \textbf{Rope (transverse):} Fix one end of a rope to a vibrator driven by a signal generator ($f \approx 2$--$5\ \mathrm{Hz}$). Observe the travelling sine wave. Increase frequency $\rightarrow$ wavelength decreases.
\item \textbf{Slinky (longitudinal):} Stretch a slinky on a bench. Give a sharp longitudinal push at one end. Observe a compression pulse travelling. For a continuous wave, drive one end longitudinally with a low-frequency vibrator.
\item \textbf{Ripple tank (transverse water waves):} Generate straight waves with a straight dipper. Use a strobe light to ``freeze'' the pattern and measure $\lambda$ directly on the screen.
\end{enumerate}
\textbf{Observation:} In the rope and ripple tank, particles move perpendicular to propagation. In the slinky, coils move back and forth parallel to propagation.
\end{experience}

\courssub{Characteristics of a Periodic Wave}

A \cle{periodic wave} repeats its pattern in space and time. Five key quantities describe it:

\begin{center}
\begin{tabularx}{\linewidth}{|l|X|l|}\hline
\rowcolor{popblueL}
\textbf{Quantity} & \textbf{Definition} & \textbf{SI Unit} \\ \hline
Amplitude $A$ & Maximum displacement of a particle from its equilibrium position. & metre ($\mathrm{m}$) \\ \hline
Period $T$ & Time for one complete oscillation of a particle (or for the wave to travel one wavelength). & second ($\mathrm{s}$) \\ \hline
Frequency $f$ & Number of oscillations per unit time: $f = \dfrac{1}{T}$. & hertz ($\mathrm{Hz}$) \\ \hline
Wavelength $\lambda$ & Distance between two consecutive points in phase (e.g. crest to crest, compression to compression). & metre ($\mathrm{m}$) \\ \hline
Wave speed $v$ & Speed at which the wave pattern (energy) travels through the medium. & metre per second ($\mathrm{m\,s^{-1}}$) \\ \hline
\end{tabularx}
\end{center}

\begin{popfigure}
\begin{tikzpicture}[scale=1.2]
% Sine wave with labels
\draw[popblue,thick] (-0.5,0) -- (8.5,0) node[right] {$x$};
\draw[popink,thick,domain=0:8,smooth,variable=\x] plot ({\x},{0.6*sin(360*\x/2)});
% Amplitude
\draw[poporange,<->,thick] (1,0) -- (1,0.6) node[midway,right] {$A$};
% Wavelength
\draw[poporange,<->,thick] (0,-0.9) -- (2,-0.9) node[midway,below] {$\lambda$};
\draw[poporange,<->,thick] (2,-0.9) -- (4,-0.9) node[midway,below] {$\lambda$};
% Crest and trough
\node[popdark,above] at (1,0.6) {Crest};
\node[popdark,below] at (3,-0.6) {Trough};
% Points in phase
\fill[popink] (1,0.6) circle (0.06);
\fill[popink] (3,0.6) circle (0.06);
\fill[popink] (5,0.6) circle (0.06);
\draw[popmuted,dashed] (1,0.6) -- (1,-1.2);
\draw[popmuted,dashed] (3,0.6) -- (3,-1.2);
\draw[popmuted,dashed] (5,0.6) -- (5,-1.2);
\node[popmuted,font=\tiny] at (1,-1.3) {in phase};
\node[popmuted,font=\tiny] at (3,-1.3) {in phase};
\node[popmuted,font=\tiny] at (5,-1.3) {in phase};
\end{tikzpicture}
\legende{Labelled sinusoidal wave showing amplitude $A$, wavelength $\lambda$, crests and troughs. Points separated by integer multiples of $\lambda$ are \emph{in phase}.}
\end{popfigure}

\courssub{The Wave Equation: $v = f\lambda$}

\begin{propriete}[Wave Equation]
For any periodic wave, the wave speed $v$, frequency $f$ and wavelength $\lambda$ are related by
\[
v = f\lambda = \frac{\lambda}{T}.
\]
\textbf{Derivation:} In one period $T$, the wave travels a distance equal to one wavelength $\lambda$. By definition of speed, $v = \frac{\text{distance}}{\text{time}} = \frac{\lambda}{T} = f\lambda$.
This equation holds for \textbf{all} waves (mechanical and electromagnetic) provided $v$ is the speed in the given medium.
\end{propriete}

\begin{exemplebox}[Using the Wave Equation]
A sound wave in air at $20^\circ\mathrm{C}$ has frequency $f = 440\ \mathrm{Hz}$ (concert A). The speed of sound in air at this temperature is $v \approx 343\ \mathrm{m\,s^{-1}}$. Find the wavelength.
\[
\lambda = \frac{v}{f} = \frac{343}{440} = 0.780\ \mathrm{m}.
\]
If the same frequency is produced in water where $v \approx 1480\ \mathrm{m\,s^{-1}}$, the wavelength becomes $\lambda = 1480/440 = 3.36\ \mathrm{m}$. \textbf{Key point:} frequency is determined by the source and does not change when a wave passes from one medium to another; $v$ and $\lambda$ change together.
\end{exemplebox}

\courssub{Factors Affecting Wave Speed}

The speed of a mechanical wave depends \textbf{only on the properties of the medium}, not on the frequency or amplitude (for small amplitudes, linear regime).

\begin{propriete}[Speed of a Transverse Wave on a Stretched String]
For a string of mass per unit length $\mu$ (linear density) under tension $T$,
\[
v = \sqrt{\frac{T}{\mu}}.
\]
\begin{itemize}
\item Increasing tension $T$ $\rightarrow$ increases $v$.
\item Increasing mass per unit length $\mu$ (thicker/heavier string) $\rightarrow$ decreases $v$.
\end{itemize}
\end{propriete}

\begin{propriete}[Speed of Sound in a Fluid]
For a fluid (liquid or gas) of density $\rho$ and bulk modulus $K$,
\[
v = \sqrt{\frac{K}{\rho}}.
\]
For an ideal gas, $v = \sqrt{\frac{\gamma RT}{M}}$, where $\gamma = C_p/C_v$, $R$ is the molar gas constant, $T$ is absolute temperature, $M$ is molar mass.
\begin{itemize}
\item In air at $0^\circ\mathrm{C}$: $v \approx 331\ \mathrm{m\,s^{-1}}$.
\item In air at $20^\circ\mathrm{C}$: $v \approx 343\ \mathrm{m\,s^{-1}}$.
\item Approximate formula: $v \approx 331 + 0.6\,\theta$ (with $\theta$ in $^\circ\mathrm{C}$).
\item Sound travels faster in liquids ($\sim 1500\ \mathrm{m\,s^{-1}}$ in water) and solids ($\sim 5000\ \mathrm{m\,s^{-1}}$ in steel) because they are much less compressible (larger $K$).
\end{itemize}
\end{propriete}

\begin{attention}[Wave Speed vs Particle Speed]
\textbf{Common mistake:} Confusing the \textbf{wave speed} $v$ (speed of energy propagation) with the \textbf{particle speed} $u = \dfrac{\partial y}{\partial t}$ (speed of oscillation of a medium particle).
\begin{itemize}
\item Wave speed $v$ is constant for a given medium and wave type (e.g. $340\ \mathrm{m\,s^{-1}}$ for sound in air).
\item Particle speed $u$ varies sinusoidally: $u_{\max} = 2\pi f A = \omega A$. For a sound wave of $f=1000\ \mathrm{Hz}$, $A=10^{-6}\ \mathrm{m}$, $u_{\max} \approx 0.006\ \mathrm{m\,s^{-1}}$, while $v = 340\ \mathrm{m\,s^{-1}}$.
\end{itemize}
In exams, read carefully: ``speed of the wave'' $\neq$ ``speed of the particles''.
\end{attention}

\courssub{Measuring Wavelength and Period in the Laboratory}

\begin{methode}[Measuring $\lambda$ and $T$ from a Ripple Tank or Rope]
\begin{enumerate}
\item \textbf{Set up a steady periodic wave} (constant frequency $f$ from a signal generator).
\item \textbf{Freeze the pattern} using a strobe light synchronized with the generator (ripple tank) or take a photograph (rope).
\item \textbf{Measure $\lambda$}: On the frozen image, measure the distance between two consecutive crests (or compressions). Use a scale on the screen or a ruler on the photograph. $\lambda = \dfrac{\text{measured distance on image}}{\text{magnification factor}}$.
\item \textbf{Measure $T$}: $T = 1/f$ where $f$ is the known frequency of the signal generator. Alternatively, use a high-speed camera to record the motion of one particle and measure the time between two successive maxima.
\item \textbf{Calculate $v$}: $v = f\lambda$ (or $v = \lambda/T$).
\end{enumerate}
\end{methode}

\begin{exoresolu}[Wave on a Stretched String]
A string of length $L = 2.00\ \mathrm{m}$ and mass $m = 5.00\ \mathrm{g}$ is stretched under a tension $T = 80.0\ \mathrm{N}$. A vibrator at one end generates a transverse wave of frequency $f = 120\ \mathrm{Hz}$. Find (a) the wave speed, (b) the wavelength, (c) the maximum transverse speed of a particle if the amplitude is $A = 4.00\ \mathrm{mm}$.

\tcblower
\textbf{Solution:}
\begin{enumerate}
\item Linear density $\mu = \dfrac{m}{L} = \dfrac{5.00\times10^{-3}\ \mathrm{kg}}{2.00\ \mathrm{m}} = 2.50\times10^{-3}\ \mathrm{kg\,m^{-1}}$.
\item Wave speed $v = \sqrt{\dfrac{T}{\mu}} = \sqrt{\dfrac{80.0}{2.50\times10^{-3}}} = \sqrt{32000} = 179\ \mathrm{m\,s^{-1}}$.
\item Wavelength $\lambda = \dfrac{v}{f} = \dfrac{179}{120} = 1.49\ \mathrm{m}$.
\item Angular frequency $\omega = 2\pi f = 240\pi\ \mathrm{rad\,s^{-1}}$.
  Maximum particle speed $u_{\max} = \omega A = 240\pi \times 4.00\times10^{-3} = 3.02\ \mathrm{m\,s^{-1}}$.
  \textbf{Note:} $u_{\max} \ll v$ ($3.02$ vs $179\ \mathrm{m\,s^{-1}}$), illustrating the warning above.
\end{enumerate}
\end{exoresolu}

\courssub{Key Points to Remember}

\begin{aretenir}[Summary of Section 1]
\begin{itemize}
\item A wave transfers energy without net transfer of matter.
\item Mechanical waves need a medium; electromagnetic waves do not.
\item Transverse waves: particle vibration $\perp$ propagation (rope, water ripples, light).
\item Longitudinal waves: particle vibration $\parallel$ propagation; compressions and rarefactions (sound, slinky).
\item Characteristics: $A$, $T$, $f=1/T$, $\lambda$, $v$.
\item Universal wave equation: $v = f\lambda = \lambda/T$.
\item Wave speed depends on medium properties only (tension and $\mu$ for string; $K$ and $\rho$ for fluids).
\item Frequency is source-dependent; $v$ and $\lambda$ change with medium.
\item \textbf{Don't confuse} wave speed $v$ with particle oscillation speed $u$.
\end{itemize}
\end{aretenir}

\begin{exercice}[Practice Problems]
\begin{itemize}
\item A fisherman on Lake Nyos observes that wave crests pass his boat every $2.5\ \mathrm{s}$. He estimates the distance between crests as $6.0\ \mathrm{m}$. What is the speed of the water waves?
\item A transverse wave on a string is described by $y(x,t) = 0.03\sin(2\pi(50t - 0.2x))$ with $y,x$ in metres, $t$ in seconds. Determine (a) amplitude, (b) frequency, (c) wavelength, (d) wave speed, (e) direction of propagation.
\item The speed of sound in air at $27^\circ\mathrm{C}$ is approximately $347\ \mathrm{m\,s^{-1}}$. A tuning fork of frequency $512\ \mathrm{Hz}$ is struck. Find the wavelength of the sound wave in air at this temperature.
\item A steel wire of length $1.50\ \mathrm{m}$ and mass $12.0\ \mathrm{g}$ is under tension $100\ \mathrm{N}$. Calculate the speed of transverse waves on the wire. If the wire is driven at $400\ \mathrm{Hz}$, what is the wavelength?
\item Explain why sound cannot be heard on the Moon, but light from the Sun reaches the Moon and Earth.
\end{itemize}
\end{exercice}

\begin{corrige}[Exercise 1]
$T = 2.5\ \mathrm{s}$, $\lambda = 6.0\ \mathrm{m}$. $v = \lambda/T = 6.0/2.5 = 2.4\ \mathrm{m\,s^{-1}}$.
\end{corrige}

\begin{corrige}[Exercise 2]
Given $y = 0.03\sin(2\pi(50t - 0.2x)) = 0.03\sin(100\pi t - 0.4\pi x)$.
(a) $A = 0.03\ \mathrm{m} = 3.0\ \mathrm{cm}$.
(b) $\omega = 100\pi \Rightarrow f = \omega/(2\pi) = 50\ \mathrm{Hz}$.
(c) $k = 0.4\pi \Rightarrow \lambda = 2\pi/k = 2\pi/(0.4\pi) = 5.0\ \mathrm{m}$.
(d) $v = f\lambda = 50 \times 5.0 = 250\ \mathrm{m\,s^{-1}}$ (or $v = \omega/k = 100\pi/(0.4\pi) = 250\ \mathrm{m\,s^{-1}}$).
(e) The phase is $(50t - 0.2x)$; as $t$ increases, $x$ must increase to keep phase constant $\Rightarrow$ wave travels in $+x$ direction.
\end{corrige}

\begin{corrige}[Exercise 3]
$\lambda = v/f = 347/512 = 0.678\ \mathrm{m}$.
\end{corrige}

\begin{corrige}[Exercise 4]
$\mu = 0.012/1.50 = 8.00\times10^{-3}\ \mathrm{kg\,m^{-1}}$.
$v = \sqrt{100/(8.00\times10^{-3})} = \sqrt{12500} = 112\ \mathrm{m\,s^{-1}}$.
$\lambda = v/f = 112/400 = 0.280\ \mathrm{m}$.
\end{corrige}

\begin{corrige}[Exercise 5]
Sound is a mechanical longitudinal wave requiring a material medium (air, water, solid) to propagate. The Moon has no atmosphere (vacuum), so there is no medium for sound waves to travel through. Light is an electromagnetic wave; it does not require a medium and propagates perfectly through vacuum. Hence sunlight reaches the Moon and Earth, but sound cannot travel on the Moon.
\end{corrige}

\coursec{The Progressive Wave: Equation and Graphical Representation}

In the previous section we established what a mechanical wave is: a disturbance that propagates through a material medium, transferring energy without a net transport of matter. We also classified waves (transverse vs longitudinal, one-dimensional vs two- or three-dimensional) and introduced the fundamental characteristics: amplitude $A$, period $T$, frequency $f$, wavelength $\lambda$ and wave speed $v$. Now we go a step further: we want a \emph{mathematical description} that tells us the displacement of any particle of the medium at any position $x$ and any time $t$. This is the equation of a \emph{progressive (travelling) wave}.

\courssub{Definition of a Progressive Wave and the Harmonic Progressive Wave}

\begin{definition}[Progressive (Travelling) Wave]
A \cle{progressive wave} is a wave whose spatial profile (shape) moves through the medium with a constant speed $v$ without changing its form (in a non-dispersive, lossless medium). Every particle of the medium executes the same vibratory motion as the source, but with a time delay that increases with distance from the source.
\end{definition}

If the source vibrates with simple harmonic motion (SHM), the wave is called a \cle{harmonic progressive wave} (or sinusoidal progressive wave). This is the most important case because, by Fourier's theorem, any periodic wave can be decomposed into a sum of harmonic waves.

\courssub{Derivation of the Wave Equation: The Delay Principle}

Consider a one-dimensional wave travelling along the $+x$ direction. The source at $x=0$ oscillates according to
\[
y(0,t) = A \sin(\omega t) \quad \text{or} \quad y(0,t) = A \cos(\omega t)
\]
where $\omega = 2\pi f = \frac{2\pi}{T}$ is the angular frequency. A particle situated at position $x$ receives the disturbance after a time delay $\tau = \frac{x}{v}$ (the time needed for the wave to travel the distance $x$ at speed $v$). Hence the displacement of the particle at $x$ at time $t$ is exactly the displacement the source had at the earlier time $t - \tau$:
\[
y(x,t) = y(0,\, t - \tfrac{x}{v}) = A \sin\bigl[\omega(t - \tfrac{x}{v})\bigr].
\]

Expanding the argument:
\[
\omega(t - \tfrac{x}{v}) = \omega t - \frac{\omega}{v}x = \omega t - \frac{2\pi f}{f\lambda}x = \omega t - \frac{2\pi}{\lambda}x.
\]
We define the \cle{wave number} (or propagation constant) $k = \frac{2\pi}{\lambda}$ (SI unit: rad\,m$^{-1}$). The wave speed satisfies $v = \frac{\omega}{k} = f\lambda = \frac{\lambda}{T}$.

\begin{propriete}[Equation of a One-Dimensional Harmonic Progressive Wave]
For a wave travelling in the \textbf{$+x$ direction}:
\[
\boxed{y(x,t) = A \sin(\omega t - k x)} \qquad \text{or} \qquad y(x,t) = A \cos(\omega t - k x)
\]
where
\begin{itemize}
  \item $A$ = amplitude (maximum displacement), SI unit: metre (m);
  \item $\omega = 2\pi f = \frac{2\pi}{T}$ = angular frequency, SI unit: rad\,s$^{-1}$;
  \item $k = \frac{2\pi}{\lambda}$ = wave number, SI unit: rad\,m$^{-1}$;
  \item $v = \frac{\omega}{k} = f\lambda = \frac{\lambda}{T}$ = wave speed, SI unit: m\,s$^{-1}$.
\end{itemize}
For a wave travelling in the \textbf{$-x$ direction}, the sign in front of $kx$ changes:
\[
y(x,t) = A \sin(\omega t + k x) \quad \text{or} \quad y(x,t) = A \cos(\omega t + k x).
\]
\end{propriete}

\begin{aretenir}[Key Relations]
\[
\omega = 2\pi f = \frac{2\pi}{T}, \qquad k = \frac{2\pi}{\lambda}, \qquad v = \frac{\omega}{k} = f\lambda = \frac{\lambda}{T}.
\]
The argument $\phi(x,t) = \omega t \mp k x$ is called the \cle{phase} of the wave at $(x,t)$.
\end{aretenir}

\courssub{Phase and Phase Difference}

The phase of a particle at position $x$ and time $t$ for a $+x$ travelling wave is $\phi(x,t) = \omega t - k x$. Two particles at positions $x_1$ and $x_2$ at the \emph{same instant} $t$ have a phase difference
\[
\Delta \phi = \phi(x_2,t) - \phi(x_1,t) = -k(x_2 - x_1) = -\frac{2\pi}{\lambda}\,\Delta x.
\]
The magnitude of the phase difference is therefore
\[
|\Delta \phi| = \frac{2\pi}{\lambda}\,|\Delta x|.
\]

\begin{propriete}[Phase Difference and Spatial Separation]
\begin{itemize}
  \item Two points are \textbf{in phase} ($\Delta \phi = 2n\pi,\ n\in\mathbb{Z}$) if their separation is an integer multiple of the wavelength: $\Delta x = n\lambda$.
  \item Two points are \textbf{in opposition of phase} ($\Delta \phi = (2n+1)\pi$) if $\Delta x = (n+\frac12)\lambda$.
\end{itemize}
\end{propriete}

\begin{popfigure}
\begin{tikzpicture}[scale=1.1, every node/.style={font=\small}]
% Axes
\draw[->] (-0.5,0) -- (7.5,0) node[right] {$x$};
\draw[->] (0,-1.8) -- (0,1.8) node[above] {$y$};
% Sine wave: period = 2 (lambda), so k = 2pi/lambda = pi rad/unit -> 180 deg/unit
\draw[popblue, thick, domain=0:7, samples=200] plot (\x, {1.5*sin(180*\x)});
% Wavelength markers
\draw[poporange, <->] (0,1.7) -- (2,1.7) node[midway, above] {$\lambda$};
\draw[poporange, <->] (2,1.7) -- (4,1.7) node[midway, above] {$\lambda$};
% Amplitude
\draw[popgreen, <->] (0,0) -- (0,1.5) node[midway, left] {$A$};
% Points at x=0.5 and x=2.5
\fill[popdark] (0.5, {1.5*sin(180*0.5)}) circle (2pt) node[above right] {$x_1$};
\fill[popdark] (2.5, {1.5*sin(180*2.5)}) circle (2pt) node[above right] {$x_2$};
% Phase difference annotation
\draw[poppurple, dashed] (0.5, -1.5) -- (0.5, {1.5*sin(180*0.5)});
\draw[poppurple, dashed] (2.5, -1.5) -- (2.5, {1.5*sin(180*2.5)});
\draw[poppurple, <->] (0.5, -1.6) -- (2.5, -1.6) node[midway, below] {$\Delta x$};
\node[poppurple] at (1.5, -1.2) {$\Delta\phi = \frac{2\pi}{\lambda}\Delta x$};
% Tick marks
\foreach \x in {0,2,4,6} \draw (\x, -0.05) -- (\x, 0.05) node[below] {$\x$};
\foreach \y in {-1.5, -1, -0.5, 0.5, 1, 1.5} \draw (-0.05, \y) -- (0.05, \y) node[left] {$\y$};
\end{tikzpicture}
\legende{Snapshot of a harmonic progressive wave at a fixed instant $t$. The phase difference between two points separated by $\Delta x$ is $\Delta\phi = \frac{2\pi}{\lambda}\Delta x$.}
\end{popfigure}

\courssub{Graphical Representations of a Progressive Wave}

A progressive wave is a function of two variables ($x$ and $t$). To visualise it on paper we must freeze one variable. This gives two complementary graphs.

\begin{enumerate}
  \item \textbf{Snapshot (Wave Profile):} $y$ versus $x$ at a \emph{fixed time} $t = t_0$. This is a photograph of the wave at one instant. The horizontal distance between two consecutive crests (or troughs) is the \textbf{wavelength $\lambda$}. The vertical maximum is the \textbf{amplitude $A$}.
  \item \textbf{Vibration Graph (Time History):} $y$ versus $t$ at a \emph{fixed position} $x = x_0$. This shows how a single particle moves with time. The horizontal distance between two consecutive peaks is the \textbf{period $T$}. The vertical maximum is again the \textbf{amplitude $A$}.
\end{enumerate}

\begin{popfigure}
\begin{tikzpicture}
\begin{axis}[
  width=0.9\linewidth, height=5cm,
  axis lines=middle,
  xlabel={$x$}, ylabel={$y$},
  xmin=0, xmax=4*pi,
  ymin=-2, ymax=2,
  xtick={0, pi, 2*pi, 3*pi, 4*pi},
  xticklabels={$0$, $\pi$, $2\pi$, $3\pi$, $4\pi$},
  ytick={-1.5, 0, 1.5},
  yticklabels={$-A$, $0$, $A$},
  samples=200,
  domain=0:4*pi,
  clip=false
]
\addplot[popblue, thick] {1.5*sin(deg(x))};
% Wavelength annotation
\draw[poporange, <->] (pi, 1.7) -- (3*pi, 1.7) node[midway, above] {$\lambda$};
\draw[poporange, <->] (0, 1.7) -- (pi, 1.7) node[midway, above] {$\lambda$};
% Amplitude annotation
\draw[popgreen, <->] (pi/2, 0) -- (pi/2, 1.5) node[midway, right] {$A$};
\end{axis}
\end{tikzpicture}
\legende{Snapshot $y(x)$ at fixed $t$: the spatial period is $\lambda$.}
\end{popfigure}

\begin{popfigure}
\begin{tikzpicture}
\begin{axis}[
  width=0.9\linewidth, height=5cm,
  axis lines=middle,
  xlabel={$t$}, ylabel={$y$},
  xmin=0, xmax=4*pi,
  ymin=-2, ymax=2,
  xtick={0, pi, 2*pi, 3*pi, 4*pi},
  xticklabels={$0$, $\pi$, $2\pi$, $3\pi$, $4\pi$},
  ytick={-1.5, 0, 1.5},
  yticklabels={$-A$, $0$, $A$},
  samples=200,
  domain=0:4*pi,
  clip=false
]
\addplot[poporange, thick] {1.5*sin(deg(x))};
% Period annotation
\draw[popblue, <->] (0, 1.7) -- (2*pi, 1.7) node[midway, above] {$T$};
\draw[popblue, <->] (2*pi, 1.7) -- (4*pi, 1.7) node[midway, above] {$T$};
% Amplitude annotation
\draw[popgreen, <->] (pi/2, 0) -- (pi/2, 1.5) node[midway, right] {$A$};
\end{axis}
\end{tikzpicture}
\legende{Vibration graph $y(t)$ at fixed $x$: the temporal period is $T$.}
\end{popfigure}

\courssub{Determining the Direction of Propagation from the Two Graphs}

This is a classic GCE examination skill. Given a snapshot $y(x)$ at time $t=0$ and the vibration graph $y(t)$ of a particular particle (usually at $x=0$), you can deduce the direction of travel by comparing the immediate future of the particle with the spatial profile to its left and right.

\begin{methode}[Finding the Direction of Propagation (Graphical Method)]
\begin{enumerate}
  \item On the snapshot $y(x,0)$, note the displacement of the particle at $x=0$ and the displacements of its neighbours at small $x>0$ (right) and $x<0$ (left).
  \item On the vibration graph $y(0,t)$, observe the displacement of the particle at $x=0$ for small $t>0$ (just after the snapshot instant). Is it moving upward (increasing $y$) or downward (decreasing $y$)?
  \item \textbf{Comparison:} For a wave travelling in the $+x$ direction, the profile moves right. The particle at $x=0$ at time $t>0$ will have the displacement that was at $x<0$ at $t=0$. Therefore, the vibration graph for $t>0$ must match the snapshot for $x<0$. For a wave travelling in the $-x$ direction, the vibration for $t>0$ matches the snapshot for $x>0$.
  \item \textbf{In practice:} If the particle at $x=0$ moves \emph{towards} the displacement of its \emph{left} neighbour ($x<0$), the wave travels $+x$. If it moves \emph{towards} the displacement of its \emph{right} neighbour ($x>0$), the wave travels $-x$.
\end{enumerate}
\end{methode}

\begin{exemplebox}[Direction of Propagation -- Corrected Analysis]
A snapshot at $t=0$ shows $y(x,0) = A\sin(kx)$ (with $k>0$). The particle at $x=0$ has $y=0$. For small $x>0$, $\sin(kx)>0$ so the right neighbour has positive displacement. For small $x<0$, $\sin(kx)<0$ so the left neighbour has negative displacement.
The vibration graph of the particle at $x=0$ shows $y(0,t) = A\sin(\omega t)$. Just after $t=0$, $y$ becomes \emph{positive} (moves upward).
The particle at $x=0$ moves toward the displacement of its \emph{right} neighbour (both positive). According to the method, this means the wave travels in the \textbf{$-x$ direction}.
Indeed, the full equation is $y(x,t) = A\sin(\omega t + kx)$.
\end{exemplebox}

\courssub{Worked Examples: Extracting Parameters and Writing the Equation}

\begin{exoresolu}[From Graphs to Equation]
The figure below shows (a) a snapshot of a progressive wave at $t=0$ and (b) the vibration graph of the particle at $x=0$.

\begin{popfigure}
\begin{tikzpicture}
\begin{scope}
\begin{axis}[
  width=0.45\linewidth, height=4cm,
  axis lines=middle,
  xlabel={$x\,\mathrm{(m)}$}, ylabel={$y\,\mathrm{(cm)}$},
  xmin=0, xmax=4,
  ymin=-3.5, ymax=3.5,
  xtick={0,1,2,3,4},
  ytick={-3,-2,-1,0,1,2,3},
  samples=200, domain=0:4,
  clip=false
]
\addplot[popblue, thick] {3*sin(deg(pi*x))};
\draw[poporange, <->] (0, 3.2) -- (2, 3.2) node[midway, above] {$\lambda=2\,\mathrm{m}$};
\draw[popgreen, <->] (0.5, 0) -- (0.5, 3) node[midway, right] {$A=3\,\mathrm{cm}$};
\end{axis}
\node at (0.225\linewidth, -2.8) {(a) Snapshot at $t=0$};
\end{scope}
\begin{scope}[xshift=0.55\linewidth]
\begin{axis}[
  width=0.45\linewidth, height=4cm,
  axis lines=middle,
  xlabel={$t\,\mathrm{(s)}$}, ylabel={$y\,\mathrm{(cm)}$},
  xmin=0, xmax=0.04,
  ymin=-3.5, ymax=3.5,
  xtick={0,0.01,0.02,0.03,0.04},
  ytick={-3,-2,-1,0,1,2,3},
  samples=200, domain=0:0.04,
  clip=false
]
\addplot[poporange, thick] {3*sin(deg(2*pi*50*x))};
\draw[popblue, <->] (0, 3.2) -- (0.02, 3.2) node[midway, above] {$T=0.02\,\mathrm{s}$};
\draw[popgreen, <->] (0.005, 0) -- (0.005, 3) node[midway, right] {$A=3\,\mathrm{cm}$};
\end{axis}
\node at (0.225\linewidth, -2.8) {(b) Vibration at $x=0$};
\end{scope}
\end{tikzpicture}
\legende{Graphs for the worked example.}
\end{popfigure}

\textbf{Determine:} (i) amplitude $A$, wavelength $\lambda$, period $T$, frequency $f$, wave speed $v$; (ii) the wave equation $y(x,t)$; (iii) the direction of propagation.

\tcblower
\textbf{Solution.}
\begin{enumerate}
  \item \textbf{Reading the graphs:}
  \begin{itemize}
    \item From (a): $A = 3\,\mathrm{cm} = 0.03\,\mathrm{m}$; $\lambda = 2\,\mathrm{m}$.
    \item From (b): $T = 0.02\,\mathrm{s}$; hence $f = \frac{1}{T} = 50\,\mathrm{Hz}$.
    \item Wave speed: $v = f\lambda = 50 \times 2 = 100\,\mathrm{m\,s^{-1}}$.
  \end{itemize}
  \item \textbf{Angular frequency and wave number:}
  \[
  \omega = 2\pi f = 100\pi\,\mathrm{rad\,s^{-1}}, \qquad k = \frac{2\pi}{\lambda} = \pi\,\mathrm{rad\,m^{-1}}.
  \]
  \item \textbf{Direction of propagation:}
  At $t=0$, the snapshot (a) gives $y(x,0) = 3\sin(\pi x)$ (in cm). At $x=0$, $y=0$. For small $x>0$, $\sin(\pi x)>0$ so the right neighbour has \emph{positive} displacement. For small $x<0$, the left neighbour has \emph{negative} displacement.
  The vibration graph (b) at $x=0$ shows $y(0,t) = 3\sin(100\pi t)$. Just after $t=0$, $y$ becomes \emph{positive} (since $\sin(100\pi t) > 0$ for small $t>0$).
  The particle at $x=0$ moves \emph{towards} the displacement of its \emph{right} neighbour (both become positive). Therefore the wave travels in the \textbf{$-x$ direction}.
  \item \textbf{Wave equation:}
  For $-x$ direction: $y(x,t) = A \sin(\omega t + k x)$.
  \[
  y(x,t) = 0.03 \sin(100\pi t + \pi x) \quad \text{(SI units: m, s)}.
  \]
  In cm and m mixed (often accepted in GCE if consistent):
  \[
  y(x,t) = 3 \sin(100\pi t + \pi x) \quad \text{with $y$ in cm, $x$ in m, $t$ in s}.
  \]
\end{enumerate}
\end{exoresolu}

\begin{exoresolu}[Phase Difference Calculation]
Two points on a progressive wave are separated by $\Delta x = 0.5\,\mathrm{m}$. The wave has frequency $f = 400\,\mathrm{Hz}$ and speed $v = 320\,\mathrm{m\,s^{-1}}$. Calculate the phase difference between the two points.

\tcblower
\textbf{Solution.}
\[
\lambda = \frac{v}{f} = \frac{320}{400} = 0.8\,\mathrm{m}.
\]
\[
\Delta \phi = \frac{2\pi}{\lambda}\,\Delta x = \frac{2\pi}{0.8} \times 0.5 = \frac{\pi}{0.8} = 1.25\pi = \frac{5\pi}{4}\,\mathrm{rad}.
\]
(Equivalently $225^\circ$.) The points are neither in phase nor in opposition of phase.
\end{exoresolu}

\courssub{Common Mistakes and Examination Traps}

\begin{attention}[Mixing up $y(x)$ and $y(t)$]
The most frequent error is reading \textbf{$\lambda$ on a $y(t)$ graph} or \textbf{$T$ on a $y(x)$ graph}. Always check the horizontal axis label:
\begin{itemize}
  \item Horizontal axis = $x$ (metres) $\rightarrow$ spatial period = $\lambda$.
  \item Horizontal axis = $t$ (seconds) $\rightarrow$ temporal period = $T$.
\end{itemize}
\end{attention}

\begin{attention}[Phase Difference Formula]
Phase difference is $\Delta \phi = \frac{2\pi}{\lambda}\Delta x$ (spatial separation over wavelength). \textbf{Never} use $\frac{2\pi}{T}\Delta x$ or $\frac{2\pi}{T}\Delta t$ for spatial phase difference. For a time difference $\Delta t$ at the \emph{same point}, the phase change is $\Delta \phi = \omega \Delta t = \frac{2\pi}{T}\Delta t$.
\end{attention}

\begin{attention}[Sign in the Wave Equation]
For a wave travelling in the $+x$ direction the phase is $\omega t - kx$ (minus sign). For $-x$ direction it is $\omega t + kx$ (plus sign). A common mistake is to write the opposite. Remember: a point at larger $x$ experiences the disturbance \emph{later}, so its phase lags behind $\rightarrow$ subtract $kx$ for $+x$ travel.
\end{attention}

\begin{attention}[Units Consistency]
When writing the final equation $y(x,t) = A\sin(\omega t \mp kx)$, ensure that if $x$ is in metres, $k$ is in rad\,m$^{-1}$; if $t$ is in seconds, $\omega$ is in rad\,s$^{-1}$. The argument of sine must be dimensionless (radians). Do not mix cm and m without conversion.
\end{attention}

\begin{attention}[Direction of Propagation from Graphs]
Do not rely on the oversimplified rule "particle moves towards right neighbour $\rightarrow$ +x direction". The correct rule is: \textbf{compare the vibration for $t>0$ with the snapshot for $x<0$ (left) and $x>0$ (right)}. If it matches the left side, wave goes $+x$; if it matches the right side, wave goes $-x$. Always verify with the constant-phase argument $\omega t \mp kx = \text{const}$.
\end{attention}

\courssub{Summary of Key Formulas}

\begin{aretenir}[Progressive Wave Cheat Sheet]
\begin{align*}
\text{Wave equation (+x)} &\quad y(x,t) = A \sin(\omega t - k x) \\
\text{Wave equation (-x)} &\quad y(x,t) = A \sin(\omega t + k x) \\
\text{Angular frequency} &\quad \omega = 2\pi f = \frac{2\pi}{T} \\
\text{Wave number} &\quad k = \frac{2\pi}{\lambda} \\
\text{Wave speed} &\quad v = \frac{\omega}{k} = f\lambda = \frac{\lambda}{T} \\
\text{Phase difference (spatial)} &\quad \Delta \phi = \frac{2\pi}{\lambda}\,\Delta x \\
\text{In phase} &\quad \Delta x = n\lambda \quad (n \in \mathbb{Z}) \\
\text{Opposition of phase} &\quad \Delta x = (n+\tfrac12)\lambda
\end{align*}
\end{aretenir}

\begin{exercice}[Practice Problems]
\begin{enumerate}
  \item A progressive wave is described by $y(x,t) = 0.02 \sin(200\pi t - 0.5\pi x)$ (SI units). Find $A$, $f$, $T$, $\lambda$, $v$ and the direction of propagation.
  \item The snapshot of a wave at $t=0$ is $y(x,0) = 4\cos(2x)$ (cm, m). The particle at $x=0$ has a vibration graph $y(0,t) = 4\cos(10t)$ (cm, s). Write the full wave equation $y(x,t)$ and state the direction of travel.
  \item Two points on a wave of wavelength $1.2\,\mathrm{m}$ are separated by $0.3\,\mathrm{m}$. What is their phase difference in radians and degrees? Are they in phase, in opposition, or neither?
  \item A wave travels at $1500\,\mathrm{m\,s^{-1}}$ in water. If the frequency is $50\,\mathrm{kHz}$, what is the wavelength? What is the phase difference between two points $0.5\,\mathrm{m}$ apart?
\end{enumerate}
\end{exercice}

\begin{corrige}[Exercise 1]
$A = 0.02\,\mathrm{m}$; $\omega = 200\pi \Rightarrow f = 100\,\mathrm{Hz},\; T = 0.01\,\mathrm{s}$; $k = 0.5\pi \Rightarrow \lambda = 4\,\mathrm{m}$; $v = \omega/k = 400\,\mathrm{m\,s^{-1}}$; sign is minus $\rightarrow$ $+x$ direction.
\end{corrige}

\begin{corrige}[Exercise 2]
Snapshot: $y(x,0)=4\cos(2x) \Rightarrow k=2\,\mathrm{rad\,m^{-1}}$. Vibration: $y(0,t)=4\cos(10t) \Rightarrow \omega=10\,\mathrm{rad\,s^{-1}}$.
At $x=0$, $y=4$ at $t=0$ (maximum). For small $x>0$, $\cos(2x) < 1$ so $y$ decreases (downward from max). For small $x<0$, $\cos(2x) < 1$ also (symmetric), so left neighbour also decreases.
For small $t>0$, $\cos(10t) < 1$, so particle at $x=0$ moves downward.
The vibration for $t>0$ (downward) matches the snapshot for $x<0$ (downward) \textbf{and} for $x>0$ (downward) because the snapshot is symmetric at a maximum. We need to look at the \emph{rate} of change or use the constant-phase method.
Phase: $\phi = 10t \mp 2x$. For $+x$: $\phi = 10t - 2x$. Constant phase $\Rightarrow 10t - 2x = \text{const} \Rightarrow dx/dt = 5 > 0$ ($+x$). For $-x$: $dx/dt = -5$.
Check: $y(x,t)=4\cos(10t - 2x)$ gives $y(x,0)=4\cos(-2x)=4\cos(2x)$ and $y(0,t)=4\cos(10t)$. Matches both graphs. The wave travels in the \textbf{$+x$ direction}.
Equation: $y(x,t)=4\cos(10t - 2x)$ (cm, m, s).
\end{corrige}

\begin{corrige}[Exercise 3]
$\Delta\phi = \frac{2\pi}{1.2}\times 0.3 = \frac{\pi}{2}\,\mathrm{rad} = 90^\circ$. Neither in phase nor opposition.
\end{corrige}

\begin{corrige}[Exercise 4]
$\lambda = v/f = 1500 / 50000 = 0.03\,\mathrm{m} = 3\,\mathrm{cm}$. $\Delta\phi = \frac{2\pi}{0.03}\times 0.5 = \frac{100\pi}{3}\,\mathrm{rad} \equiv \frac{4\pi}{3}\,\mathrm{rad}$ (mod $2\pi$) $= 240^\circ$.
\end{corrige}

This completes our rigorous treatment of the progressive wave equation and its graphical representations. In the next section we will explore what happens when waves encounter boundaries and obstacles: reflection, refraction, diffraction and the single-slit pattern.

\coursec{Wave Phenomena I: Reflection, Refraction, Diffraction and the Single Slit}

In the previous section we learned to describe a progressive wave mathematically, $y(x,t)=A\sin(\omega t - kx)$, and to represent it graphically both in space (the wave profile) and in time (the vibration of one point). We now ask a very natural question: what happens to such a wave when it meets something --- a wall, a boundary between two media, an obstacle, or a hole? Stand at the edge of a swimming pool and watch the ripples from a dropped stone: they bounce off the pool wall, they bend where the water becomes shallow, and they spread out after passing through a gap between two rocks. These everyday observations --- \cle{reflection}, \cle{refraction} and \cle{diffraction} --- are the subject of this section. They are exactly the behaviours examiners love, because they test whether you truly understand what a wave \emph{is}.

\courssub{Production of Waves, Wavefronts and Huygens' Principle}

Every wave begins at a \cle{source}: a vibrating blade in a ripple tank, a loudspeaker cone, a stone striking water. The disturbance then travels outward. To describe \emph{where} the disturbance has reached, we use the idea of a wavefront.

\begin{definition}[Wavefront and ray]
A \cle{wavefront} is a surface (in two dimensions, a line) joining all points of the medium that are vibrating \emph{in phase} at the same instant. A \cle{ray} is a line drawn perpendicular to the wavefronts, showing the direction of propagation of the wave energy.
\end{definition}

For a point source in a uniform medium, wavefronts are concentric circles (spheres in 3D) and rays are radial. Far from the source, a small portion of a circular wavefront looks almost straight: we then speak of \cle{plane wavefronts}, with parallel rays.

How does a wavefront advance? The Dutch physicist Christiaan Huygens proposed a beautifully simple picture.

\begin{propriete}[Huygens' principle]
Every point on a wavefront acts as a \cle{secondary source} of small wavelets, which spread out in the forward direction at the wave speed. The new wavefront, a short time later, is the envelope (the common tangent surface) of all these wavelets.
\end{propriete}

Huygens' principle is qualitative at this level, but it is powerful: it explains in one stroke why waves reflect with equal angles, why they bend on entering a new medium, and why they spread after passing through a gap. We shall use it repeatedly.

\courssub{Reflection of Waves}

\begin{experience}[Observing reflection in a ripple tank]
Generate straight (plane) wavefronts with a straight vibrating dipper and place a flat barrier in the tank at an angle. The incident wavefronts strike the barrier and bounce off as a new family of straight wavefronts. Measuring the directions of travel shows a striking regularity: the wave comes in and goes out at equal angles. The same experiment works with circular pulses from a point dipper reflecting off a straight barrier --- the reflected pulses appear to come from an ``image source'' behind the barrier.
\end{experience}

\begin{propriete}[Laws of reflection]
When a wave reflects from a surface:
\begin{enumerate}
\item the incident ray, the normal to the surface at the point of incidence, and the reflected ray all lie in the same plane;
\item the \cle{angle of incidence} $i$ equals the \cle{angle of reflection} $r$: $i = r$, both measured from the normal.
\end{enumerate}
Reflection changes neither the speed $v$, nor the frequency $f$, nor the wavelength $\lambda$ of the wave, because the wave stays in the same medium.
\end{propriete}

\begin{propriete}[Phase change on reflection]
When a mechanical wave reflects from a \textbf{rigid (fixed) boundary} (e.g., a pulse on a string tied to a wall, or a water wave hitting a vertical barrier), the reflected wave undergoes a \cle{phase reversal} (a phase change of $\pi$ rad or $\lambda/2$): a crest returns as a trough. When it reflects from a \textbf{free boundary} (e.g., a string attached to a frictionless ring, or the free surface of water), there is \textbf{no phase change}: a crest returns as a crest.
\end{propriete}

\begin{popfigure}
\begin{center}
\begin{tikzpicture}[scale=1.0, >=stealth]
% barrier
\fill[popmuted!40] (3,-0.35) rectangle (7,0);
\draw[thick, popink] (3,0) -- (7,0);
\node[popink] at (6.6,-0.2) {\scriptsize barrier};
% normal
\draw[dashed, popdark] (5,0) -- (5,2.6);
\node[popdark] at (5.15,2.5) {\scriptsize normal};
% incident ray
\draw[->, very thick, popblue] (2.6,2.4) -- (5,0);
% reflected ray
\draw[->, very thick, poporange] (5,0) -- (7.4,2.4);
% incident wavefronts (perpendicular to incident ray)
\draw[popblue] (2.9,2.05) -- (3.55,2.7);
\draw[popblue] (3.35,1.6) -- (4.0,2.25);
\draw[popblue] (3.8,1.15) -- (4.45,1.8);
% reflected wavefronts
\draw[poporange] (5.55,1.8) -- (6.2,1.15);
\draw[poporange] (6.0,2.25) -- (6.65,1.6);
\draw[poporange] (6.45,2.7) -- (7.1,2.05);
% angles
\draw[popdark] (4.35,0) arc (180:124:0.65);
\node[popdark] at (4.55,0.42) {\scriptsize $i$};
\draw[popdark] (5.65,0) arc (0:56:0.65);
\node[popdark] at (5.45,0.42) {\scriptsize $r$};
\node[popblue] at (2.5,2.6) {\scriptsize incident};
\node[poporange] at (7.6,2.6) {\scriptsize reflected};
\end{tikzpicture}
\end{center}
\legende{Plane wavefronts reflecting on a barrier: the rays are perpendicular to the wavefronts, and $i = r$.}
\end{popfigure}

Sound obeys the same law: an \cle{echo} is simply reflected sound. If you clap your hands facing a distant wall in Buea and hear the echo after a time $t$, the sound has travelled twice the distance $d$ to the wall, so $2d = vt$. This is the principle of \cle{echo sounding} and \cle{sonar}: a ship sends an ultrasound pulse downwards and times the returning echo to measure the depth of the sea bed.

\begin{exoresolu}[Depth of the sea bed by echo sounding]
A ship's sonar emits an ultrasound pulse and receives the echo from the sea bed after $\SI{0.40}{s}$. The speed of sound in sea water is $\SI{1500}{m.s^{-1}}$. Find the depth of the sea bed.
\tcblower
The pulse makes a round trip of length $2d$ in time $t = \SI{0.40}{s}$:
\[
2d = vt \quad\Longrightarrow\quad d = \frac{vt}{2} = \frac{1500 \times 0.40}{2} = \SI{3.0e2}{m}.
\]
The sea bed is $\SI{3.0e2}{m}$ below the ship. A common error is to forget the factor of 2: the time measured is for the \emph{go and return} journey.
\end{exoresolu}

\courssub{Refraction of Waves}

Roll a ripple-tank wave from a deep region into a shallow region: the wave slows down. Since the frequency is fixed by the source, the relation $v = f\lambda$ forces the wavelength to shrink. And if the wavefronts meet the boundary at an angle, they \emph{bend}. This is \cle{refraction}.

\begin{definition}[Refraction]
\cle{Refraction} is the change in direction of propagation of a wave when it passes from one medium into another in which its speed is different.
\end{definition}

\begin{propriete}[What changes and what stays constant in refraction]
When a wave crosses a boundary between two media:
\begin{itemize}
\item the \cle{frequency} $f$ is \textbf{unchanged} (it is imposed by the source; the boundary cannot create or destroy vibrations);
\item the speed $v$ changes (a property of the medium);
\item the wavelength changes accordingly: $\lambda = v/f$.
\end{itemize}
Consequently, for two media: $\dfrac{\lambda_1}{\lambda_2} = \dfrac{v_1}{v_2}$.
\end{propriete}

\begin{propriete}[Snell's law for waves]
The angles of incidence $i$ and refraction $r$ (measured from the normal) obey
\[
\frac{\sin i}{\sin r} = \frac{v_1}{v_2} = \frac{\lambda_1}{\lambda_2}.
\]
If $v_2 < v_1$ (wave slows down), then $r < i$: the ray bends \textbf{towards} the normal. If $v_2 > v_1$, the ray bends \textbf{away from} the normal.
\end{propriete}

Why do the wavefronts bend? Apply Huygens' idea: the end of a wavefront that enters the slow medium first is retarded while the other end is still travelling fast, so the whole wavefront swings round, exactly like a line of marchers wheeling when one end of the line steps onto muddy ground.

\begin{popfigure}
\begin{center}
\begin{tikzpicture}[scale=1.0, >=stealth]
% media
\fill[popblueL] (0,0) rectangle (8,3.2);
\fill[popgreenL] (0,-2.4) rectangle (8,0);
\draw[thick, popink] (0,0) -- (8,0);
\node[popink] at (0.9,2.9) {\scriptsize medium 1 (fast)};
\node[popink] at (0.9,-2.1) {\scriptsize medium 2 (slow)};
% normal
\draw[dashed, popdark] (4,3.0) -- (4,-2.3);
\node[popdark] at (4.2,2.8) {\scriptsize normal};
% incident ray
\draw[->, very thick, popblue] (2.2,3.0) -- (4,0);
% refracted ray (bent towards normal)
\draw[->, very thick, poporange] (4,0) -- (4.95,-2.3);
% incident wavefronts: spacing lambda1 along ray direction
\draw[popblue] (2.35,2.55) -- (3.35,3.05);
\draw[popblue] (2.75,2.0) -- (3.75,2.5);
\draw[popblue] (3.15,1.45) -- (4.15,1.95);
\draw[popblue] (3.55,0.9) -- (4.55,1.4);
% refracted wavefronts: shorter spacing lambda2
\draw[poporange] (4.28,-0.75) -- (4.83,-0.45);
\draw[poporange] (4.55,-1.5) -- (5.1,-1.2);
% labels
\node[popblue] at (1.9,2.6) {\scriptsize $\lambda_1$};
\node[poporange] at (5.6,-1.4) {\scriptsize $\lambda_2 < \lambda_1$};
% angles
\draw[popdark] (3.45,0) arc (180:124:0.55);
\node[popdark] at (3.62,0.35) {\scriptsize $i$};
\draw[popdark] (4,-0.55) arc (270:296:0.55);
\node[popdark] at (4.28,-0.72) {\scriptsize $r$};
\end{tikzpicture}
\end{center}
\legende{Refraction at a boundary: the wave slows in medium 2, so $\lambda$ decreases and the ray bends towards the normal. Frequency $f$ is unchanged.}
\end{popfigure}

\begin{attention}[The frequency never changes at a boundary]
A very common exam mistake is to say that the wavelength stays the same and the frequency changes when a wave enters a new medium. It is the \textbf{opposite}: $f$ is fixed by the source and \emph{cannot} change at the boundary; it is $v$ and $\lambda$ that adjust, keeping $v = f\lambda$ in each medium.
\end{attention}

\begin{exoresolu}[Water wave entering shallow water]
Straight waves of frequency $\SI{5.0}{Hz}$ travel in the deep part of a ripple tank where their speed is $\SI{0.30}{m.s^{-1}}$. They pass into a shallow region where the speed falls to $\SI{0.20}{m.s^{-1}}$. Find the wavelength in each region.
\tcblower
In the deep region:
\[
\lambda_1 = \frac{v_1}{f} = \frac{0.30}{5.0} = \SI{0.060}{m} = \SI{6.0}{cm}.
\]
In the shallow region, $f$ is unchanged:
\[
\lambda_2 = \frac{v_2}{f} = \frac{0.20}{5.0} = \SI{0.040}{m} = \SI{4.0}{cm}.
\]
The wavelength shortens from $\SI{6.0}{cm}$ to $\SI{4.0}{cm}$; the frequency remains $\SI{5.0}{Hz}$ throughout.
\end{exoresolu}

\courssub{Diffraction: Spreading Around Obstacles and Through Gaps}

Stand behind the corner of a building in Yaoundé while a friend talks to you from the other side: you hear him perfectly, although you cannot see him. Sound bends around the corner; light does not (noticeably). This bending or spreading of waves is called diffraction.

\begin{definition}[Diffraction]
\cle{Diffraction} is the spreading of a wave when it meets the edge of an obstacle or passes through an aperture whose size $a$ is comparable to, or smaller than, its wavelength $\lambda$.
\end{definition}

Huygens' principle makes this natural: the points of the wavefront that pass through the gap act as secondary sources, and their wavelets spread into the ``shadow'' region behind the obstacle. The ripple tank shows it beautifully:

\begin{itemize}
\item \textbf{Wide gap} ($a \gg \lambda$): the waves pass through almost straight, with only slight spreading at the edges. Diffraction is negligible.
\item \textbf{Narrow gap} ($a \lesssim \lambda$): the gap behaves almost like a point source and the waves emerge as semicircular wavefronts filling the whole region beyond. Diffraction is maximal.
\end{itemize}

\begin{popfigure}
\begin{center}
\begin{tikzpicture}[scale=0.85]
% ===== wide gap =====
\begin{scope}
% barriers
\fill[popmuted!50] (0.0,1.6) rectangle (0.25,3.2);
\fill[popmuted!50] (0.0,-0.4) rectangle (0.25,1.0);
% incoming plane wavefronts
\foreach \x in {-1.4,-0.9,-0.4} {\draw[popblue] (\x,-0.3) -- (\x,3.1);}
\draw[->, popblue, thick] (-1.7,1.4) -- (-1.1,1.4);
% transmitted: mostly straight, slight curve at edges
\foreach \x in {0.7,1.2,1.7} {\draw[poporange] (\x,0.05) -- (\x,2.55);}
\draw[poporange] (0.7,2.55) arc (90:60:0.55);
\draw[poporange] (1.2,2.55) arc (90:70:0.55);
\draw[poporange] (0.7,0.05) arc (270:300:0.55);
\draw[poporange] (1.2,0.05) arc (270:290:0.55);
\node[popink] at (0.1,3.5) {\scriptsize wide gap $a \gg \lambda$};
\end{scope}
% ===== narrow gap =====
\begin{scope}[xshift=5.2cm]
\fill[popmuted!50] (0.0,1.9) rectangle (0.25,3.2);
\fill[popmuted!50] (0.0,-0.4) rectangle (0.25,0.9);
\foreach \x in {-1.4,-0.9,-0.4} {\draw[popblue] (\x,-0.3) -- (\x,3.1);}
\draw[->, popblue, thick] (-1.7,1.4) -- (-1.1,1.4);
% semicircular wavefronts from gap centre (0.25,1.4)
\foreach \r in {0.5,1.0,1.5} {\draw[poporange] (0.25,1.4) ++(0:\r) arc (0:180:\r);}
\node[popink] at (0.1,3.5) {\scriptsize narrow gap $a \lesssim \lambda$};
\end{scope}
\end{tikzpicture}
\end{center}
\legende{Ripple-tank view of diffraction: a wide gap lets the wave through almost unchanged; a narrow gap ($a$ comparable to $\lambda$) acts like a point source and the wave spreads in all directions.}
\end{popfigure}

\begin{methode}[Will diffraction be noticeable?]
\begin{enumerate}
\item Estimate the wavelength $\lambda$ of the wave (use $\lambda = v/f$ if necessary).
\item Estimate the size $a$ of the aperture or obstacle.
\item Compare: if $a \lesssim \lambda$ (same order of magnitude or smaller), diffraction is \textbf{appreciable}; if $a \gg \lambda$, it is negligible and the wave travels essentially in straight lines.
\end{enumerate}
\end{methode}

This explains the street-corner observation quantitatively. Audible sound has wavelengths from about $\SI{2}{cm}$ to $\SI{17}{m}$; a doorway ($a \approx \SI{1}{m}$) is comparable to $\lambda$, so sound diffracts strongly through it and around corners. Visible light has $\lambda \approx \SI{5e-7}{m}$, utterly negligible compared with a doorway, so light goes straight through and casts sharp shadows. Diffraction of light is only seen with very narrow slits.

\courssub{The Single-Slit Diffraction Pattern}

Shine monochromatic light of wavelength $\lambda$ on a slit of width $a$, and observe the light on a distant screen. Instead of a single sharp bright line of width $a$ (the naive ``straight-line'' prediction), you obtain a \cle{diffraction pattern}: a broad, intense \cle{central maximum}, flanked by alternating dark and bright fringes whose intensity falls off rapidly.

\begin{propriete}[Single-slit pattern: minima positions]
For a slit of width $a$ illuminated by light of wavelength $\lambda$, the minima (dark fringes) occur at angles $\theta_n$ satisfying
\[
\boxed{\;a\sin\theta_n = n\lambda \quad (n = \pm 1, \pm 2, \pm 3, \dots)\;}
\]
The first minima ($n = \pm 1$) are at $\sin\theta = \pm\lambda/a$. The angular width of the central maximum is $\Delta\theta = \dfrac{2\lambda}{a}$ (for small angles). The central maximum is \textbf{twice as wide} as each secondary maximum, and the intensity decreases strongly away from the centre.
\end{propriete}

\textbf{Why the first minimum is at $\sin\theta=\lambda/a$ (qualitative derivation).} Divide the slit into two halves. For a direction $\theta$ such that the path difference between a ray from the top edge and a ray from the centre of the slit is exactly $\lambda/2$, these two contributions cancel by interference. But then \emph{every} point in the top half has a partner in the bottom half, a distance $a/2$ away, whose ray differs in path by $\lambda/2$ --- so the whole slit cancels in pairs and the screen is dark. The path difference across the half-slit is $\dfrac{a}{2}\sin\theta$, so the condition is $\dfrac{a}{2}\sin\theta = \dfrac{\lambda}{2}$, i.e.\ $\sin\theta = \lambda/a$. The same pairing argument applied to 4, 6, ... strips gives further minima at $\sin\theta = n\lambda/a$ ($n = 2, 3, \dots$).

\begin{popfigure}
\begin{center}
\begin{tikzpicture}[scale=1.0]
\begin{axis}[
  width=11cm, height=5.5cm,
  axis lines=middle,
  xlabel={$\sin\theta$}, ylabel={intensity $I$},
  xlabel style={at={(ticklabel* cs:1)}, anchor=west},
  ylabel style={at={(ticklabel* cs:1)}, anchor=south},
  xtick={-0.2,-0.1,0.1,0.2},
  xticklabels={$-\tfrac{2\lambda}{a}$,$-\tfrac{\lambda}{a}$,$\tfrac{\lambda}{a}$,$\tfrac{2\lambda}{a}$},
  ytick=\empty,
  xmin=-0.32, xmax=0.32, ymin=0, ymax=1.15,
  samples=400, domain=-0.31:0.31,
  declare function={sinc(\x) = (abs(\x)<1e-4) ? 1 : (sin(deg(pi*10*\x))/(pi*10*\x));},
]
\addplot[very thick, popblue] {sinc(x)^2};
\draw[dashed, popmuted] (axis cs:0.1,0) -- (axis cs:0.1,0.05);
\draw[dashed, popmuted] (axis cs:-0.1,0) -- (axis cs:-0.1,0.05);
\node[popink, align=center] at (axis cs:0,1.08) {\scriptsize central maximum};
\node[popink] at (axis cs:0.16,0.09) {\scriptsize secondary};
\end{axis}
\end{tikzpicture}
\end{center}
\legende{Single-slit intensity pattern: the first minima are at $\sin\theta = \pm\lambda/a$; the central maximum is twice as wide as the secondary maxima and much more intense.}
\end{popfigure}

On a screen at distance $D$ from the slit, with small angles ($\sin\theta \approx \tan\theta \approx \theta$), the linear width of the central maximum (distance between the first minima) is
\[
w = \frac{2\lambda D}{a}.
\]

\begin{exoresolu}[Width of the central maximum]
Light of wavelength $\SI{6.0e-7}{m}$ falls on a slit of width $\SI{0.10}{mm}$. A screen is placed $\SI{1.5}{m}$ away. (a) Find the angular position of the first minimum. (b) Find the width of the central bright fringe on the screen.
\tcblower
(a) First minimum ($n=1$):
\[
\sin\theta = \frac{\lambda}{a} = \frac{6.0\times 10^{-7}}{1.0\times 10^{-4}} = 6.0\times 10^{-3}
\quad\Longrightarrow\quad \theta \approx 6.0\times 10^{-3}\ \mathrm{rad} \approx \SI{0.34}{\degree}.
\]
(b) Width of the central maximum (distance between first minima):
\[
w = \frac{2\lambda D}{a} = \frac{2 \times 6.0\times 10^{-7} \times 1.5}{1.0\times 10^{-4}} = 1.8\times 10^{-2}\ \mathrm{m} = \SI{18}{mm}.
\]
Note how the \emph{narrower} slit gives a \emph{wider} pattern: $w \propto 1/a$. This inverse relationship is the signature of diffraction.
\end{exoresolu}

\begin{attention}[Common mistakes with the single slit]
\begin{itemize}
\item Do not confuse the single-slit formula $a\sin\theta = n\lambda$ (giving \emph{minima}) with the double-slit formula you will meet next section, where maxima occur at $d\sin\theta = n\lambda$. Same symbols, different physics: one slit gives minima, two slits give maxima.
\item Remember the factor 2: the central maximum spans from $-\lambda/a$ to $+\lambda/a$, so its angular width is $2\lambda/a$, not $\lambda/a$.
\item Diffraction becomes \emph{more} pronounced as the slit gets \emph{narrower} --- the opposite of what many students first guess.
\item The formula $a\sin\theta = n\lambda$ for minima is exact. The small-angle approximation $\theta \approx \sin\theta \approx \tan\theta = y/D$ is only valid when $\theta$ is small (typically $< 10^\circ$). Always check if the approximation is justified before using $y = n\lambda D/a$.
\end{itemize}
\end{attention}

\begin{savaistu}[Why telescopes and microscopes have limits]
Diffraction is not just a curiosity: it limits the \cle{resolving power} of every optical instrument. Light from a distant star passing through the circular aperture of a telescope is diffracted, so the star's image is a tiny disc (Airy disc) rather than a point; two stars closer together than about $1.22\lambda/D$ (in angular terms) cannot be distinguished. This is why bigger telescopes (larger aperture $D$) see finer detail. You will meet resolving power again in the extension work on optical instruments.
\end{savaistu}

\begin{aretenir}[Key points of this section]
\begin{itemize}
\item A wavefront joins points vibrating in phase; rays are perpendicular to wavefronts. Huygens: every point of a wavefront is a secondary source.
\item Reflection: $i = r$; $v$, $f$ and $\lambda$ unchanged. Phase reversal at rigid boundary, no phase change at free boundary. Applications: echoes, sonar ($2d = vt$).
\item Refraction: $f$ constant, $v$ and $\lambda$ change ($\lambda_1/\lambda_2 = v_1/v_2$); Snell's law $\sin i/\sin r = v_1/v_2$; wave bends towards normal when slowing down.
\item Diffraction is appreciable when aperture/obstacle size $a \lesssim \lambda$.
\item Single slit: minima at $a\sin\theta = n\lambda$ ($n=\pm1,\pm2,\dots$); central maximum angular width $2\lambda/a$, twice secondary maxima; linear width on screen $w = 2\lambda D/a$ (small angles).
\end{itemize}
\end{aretenir}

\begin{exercice}[Ripple tank and a gap]
Plane water waves of wavelength $\SI{1.5}{cm}$ approach a barrier containing a gap. Sketch and describe the waves beyond the barrier when the gap width is (a) $\SI{15}{cm}$, (b) $\SI{1.5}{cm}$, (c) $\SI{0.5}{cm}$. In each case state whether diffraction is appreciable.
\end{exercice}

\begin{exercice}[Refraction of sound in air layers]
A sound wave of frequency $\SI{680}{Hz}$ travels in cool air where its speed is $\SI{340}{m.s^{-1}}$, then enters a layer of warmer air where its speed is $\SI{350}{m.s^{-1}}$. Calculate the wavelength in each layer and state what happens to the frequency.
\end{exercice}

\begin{corrige}[Exercise 1]
Compare $a$ with $\lambda = \SI{1.5}{cm}$ in each case.
(a) $a = \SI{15}{cm} = 10\lambda \gg \lambda$: the waves pass through almost straight, with only slight spreading at the edges; diffraction negligible.
(b) $a = \SI{1.5}{cm} = \lambda$: strong spreading; the emerging wavefronts are markedly curved, nearly semicircular; diffraction appreciable.
(c) $a = \SI{0.5}{cm} = \lambda/3 < \lambda$: the gap acts essentially as a point source; semicircular wavefronts spread into the whole region beyond; diffraction maximal.
\end{corrige}

\begin{corrige}[Exercise 2]
The frequency is fixed by the source: $f = \SI{680}{Hz}$ in both layers.
Cool layer: $\lambda_1 = v_1/f = 340/680 = \SI{0.50}{m}$.
Warm layer: $\lambda_2 = v_2/f = 350/680 \approx \SI{0.51}{m}$.
The wavelength increases slightly in the warmer (faster) layer; the frequency is unchanged.
\end{corrige}

\coursec{Wave Phenomena II: Interference — Double and Multiple Slits, Measurement of Wavelength}

In the previous section we saw that a single slit spreads a wave into a diffraction pattern: the wave bends around edges and produces a central bright band flanked by weaker secondary bands. Diffraction showed us that a wave does not simply travel in straight lines. We now ask a deeper question: what happens when \emph{two} waves of the same kind overlap in the same region of space? The answer — \cle{interference} — is one of the most beautiful and most useful phenomena in physics. It is the phenomenon that allowed Young to prove that light is a wave, and it gives us, even in a school laboratory, a practical method to measure the incredibly small wavelength of light.

\courssub{The Superposition Principle}

Imagine two students at opposite ends of a long rope each sending a pulse towards the other. When the pulses meet, the rope does something remarkable: at the crossing point its displacement is momentarily larger (if the pulses are on the same side) or smaller (if they are on opposite sides) than either pulse alone. Then the pulses pass through each other and continue \emph{unchanged}, as if the meeting had never happened.

\begin{definition}[Superposition principle]
When two or more waves overlap at a point, the resultant displacement of the medium at that point and at that instant is the vector sum of the displacements that each wave would produce separately:
\[
\vec{y} = \vec{y}_1 + \vec{y}_2 + \cdots
\]
After overlapping, each wave continues to propagate as if the other were not present.
\end{definition}

For one-dimensional waves on a string, the displacements are along the same line, so the vector sum reduces to an algebraic sum: $y = y_1 + y_2$. Two crests meeting give a double-height crest; a crest meeting a trough of equal depth gives momentarily zero displacement — the string is flat for an instant, yet both waves are still there.

\begin{attention}[Superposition is not a new force]
Nothing ``pushes'' the waves together. Superposition simply states that the response of a linear medium to two simultaneous disturbances is the sum of the individual responses. It holds for small-amplitude mechanical waves, sound and light.
\end{attention}

\courssub{Coherent Sources and the Conditions for Interference}

Superposition alone is not enough to produce a \emph{stable, observable} pattern. If you switch on two independent torches and cross their beams, you see no dark and bright bands on the wall — just uniform illumination. Why?

\begin{definition}[Coherent sources]
Two sources are \cle{coherent} if they emit waves of the \textbf{same frequency} (and hence same wavelength) and maintain a \textbf{constant phase relationship} (a fixed phase difference) with each other.
\end{definition}

The reason is simple. The interference at a point depends on the phase difference between the two arriving waves. If that phase difference wanders randomly with time — as it does for two independent lamps, whose atoms emit in random, uncorrelated bursts — then any point on the screen is alternately bright and dark millions of times per second, and the eye records only the average: uniform light. A stable pattern requires:

\begin{itemize}
\item \textbf{same frequency} for both waves;
\item \textbf{constant phase difference} (coherence);
\item \textbf{comparable amplitudes}, so that destructive interference can produce nearly complete cancellation;
\item for light, the two waves must also have \textbf{the same plane of polarisation} (or at least a common component).
\end{itemize}

\begin{attention}[Two independent lamps never interfere]
A classic GCE trap: ``Two identical lamps are placed side by side; explain why no interference fringes are observed.'' Answer: the lamps are \cle{incoherent} — their emissions have random, rapidly changing phase differences, so no stable pattern exists. In practice, coherence is achieved by deriving both waves from a \emph{single} source (e.g.\ one slit illuminating two slits in Young's experiment).
\end{attention}

\courssub{Constructive and Destructive Interference}

Consider two coherent sources $S_1$ and $S_2$ vibrating in phase, and a point $P$ in the medium. The wave from $S_1$ travels a distance $S_1P$; the wave from $S_2$ travels $S_2P$. The \cle{path difference} is
\[
\delta = |S_2P - S_1P|.
\]

\begin{itemize}
\item If $\delta$ is a whole number of wavelengths, the two waves arrive \emph{in phase}: crest meets crest. The oscillations reinforce — \textbf{constructive interference}.
\item If $\delta$ is an odd number of half-wavelengths, the waves arrive \emph{in opposite phase}: crest meets trough. They cancel — \textbf{destructive interference}.
\end{itemize}

\begin{propriete}[Conditions for interference (proof examinable)]
For two coherent sources vibrating in phase, at a point $P$ with path difference $\delta = |S_2P - S_1P|$:
\[
\boxed{\text{Constructive (bright/antinode): } \delta = k\lambda, \qquad k = 0, 1, 2, \dots}
\]
\[
\boxed{\text{Destructive (dark/node): } \delta = (2k+1)\dfrac{\lambda}{2}, \qquad k = 0, 1, 2, \dots}
\]
\end{propriete}

\textbf{Justification.} A path difference $\delta$ corresponds to a phase difference $\Delta\varphi = \dfrac{2\pi\delta}{\lambda}$. The resultant amplitude of two waves of equal amplitude $A$ is $2A\left|\cos\dfrac{\Delta\varphi}{2}\right|$, which is maximal ($2A$) when $\Delta\varphi = 2k\pi$, i.e.\ $\delta = k\lambda$, and zero when $\Delta\varphi = (2k+1)\pi$, i.e.\ $\delta = (2k+1)\lambda/2$.

\courssub{Two-Source Interference in a Ripple Tank}

The clearest way to \emph{see} interference is the ripple tank: two dippers attached to the same vibrating bar touch the water surface, so they are automatically coherent (same frequency, fixed phase relation). Circular waves spread from each dipper and overlap.

\begin{popfigure}
\begin{center}
\begin{tikzpicture}[scale=1.0]
% water surface
\fill[popblueL] (-4.5,-2.6) rectangle (4.5,2.6);
% sources
\fill[popdark] (0,1.4) circle (0.09) node[above, text=popdark]{\small $S_1$};
\fill[popdark] (0,-1.4) circle (0.09) node[below, text=popdark]{\small $S_2$};
% circular wavefronts from S1
\foreach \r in {0.5,1.0,1.5,2.0,2.5,3.0,3.5,4.0}
  \draw[popblue, thin] (0,1.4) circle (\r);
% circular wavefronts from S2
\foreach \r in {0.5,1.0,1.5,2.0,2.5,3.0,3.5,4.0}
  \draw[poporange, thin] (0,-1.4) circle (\r);
% antinodal lines (hyperbola-like, drawn as curves)
\draw[popgreen, very thick] (0,-2.6) -- (0,2.6);
\draw[popgreen, very thick] (0.9,-2.6) .. controls (0.55,0) .. (0.9,2.6);
\draw[popgreen, very thick] (-0.9,-2.6) .. controls (-0.55,0) .. (-0.9,2.6);
\draw[popgreen, very thick] (2.0,-2.6) .. controls (1.25,0) .. (2.0,2.6);
\draw[popgreen, very thick] (-2.0,-2.6) .. controls (-1.25,0) .. (-2.0,2.6);
% nodal lines
\draw[poppink, very thick, dashed] (0.45,-2.6) .. controls (0.28,0) .. (0.45,2.6);
\draw[poppink, very thick, dashed] (-0.45,-2.6) .. controls (-0.28,0) .. (-0.45,2.6);
\draw[poppink, very thick, dashed] (1.4,-2.6) .. controls (0.85,0) .. (1.4,2.6);
\draw[poppink, very thick, dashed] (-1.4,-2.6) .. controls (-0.85,0) .. (-1.4,2.6);
\draw[poppink, very thick, dashed] (2.9,-2.6) .. controls (1.8,0) .. (2.9,2.6);
\draw[poppink, very thick, dashed] (-2.9,-2.6) .. controls (-1.8,0) .. (-2.9,2.6);
% labels
\node[text=popgreen, align=center, font=\small] at (3.6,1.9) {antinodal\\ lines};
\node[text=poppink, align=center, font=\small] at (3.6,-1.9) {nodal\\ lines};
\end{tikzpicture}
\end{center}
\legende{Interference of two coherent sources in a ripple tank. Along the solid green (antinodal) lines, $\delta = k\lambda$ and the water oscillates strongly; along the dashed pink (nodal) lines, $\delta = (2k+1)\lambda/2$ and the surface stays almost flat.}
\end{popfigure}

The overlap region shows a striking pattern:

\begin{itemize}
\item \textbf{Antinodal lines} (solid curves): loci of points where $\delta = k\lambda$; the water oscillates with maximum amplitude. The central line ($\delta = 0$, the perpendicular bisector of $S_1S_2$) is always antinodal for in-phase sources.
\item \textbf{Nodal lines} (dashed curves): loci of points where $\delta = (2k+1)\lambda/2$; the water surface remains almost still at all times.
\end{itemize}

Geometrically, each nodal or antinodal line is a hyperbola with foci at $S_1$ and $S_2$, since a hyperbola is precisely the locus of points whose difference of distances to two fixed points is constant.

\begin{experience}[Observing nodal lines in the school laboratory]
Fill a ripple tank with a shallow layer of water, fix two pointed dippers to the same vibrating bar so that they touch the surface, and illuminate from above with a stroboscopic lamp. On the white screen below the tank you see alternating regions of violent agitation and calm water. Sliding the stroboscope frequency slightly reveals the waves ``frozen'' in motion. Counting nodal lines and measuring their positions lets you estimate the wavelength of the water waves — the same logic we shall apply to light.
\end{experience}

\courssub{Young's Double-Slit Experiment: Geometry and Fringe Pattern}

In 1801 Thomas Young performed the historic experiment that demonstrated the wave nature of light. Light from a single narrow slit $S$ (this guarantees coherence) falls on two very close, parallel slits $S_1$ and $S_2$. Each slit diffracts the light, and the two diffracted beams overlap on a screen placed a large distance $D$ away. Where they overlap, interference produces a pattern of \textbf{equally spaced, parallel bright and dark fringes}.

\begin{popfigure}
\begin{center}
\begin{tikzpicture}[scale=1.0]
% barrier with slits
\fill[popmuted] (-0.08,-2.2) rectangle (0.08,-0.45);
\fill[popmuted] (-0.08,-0.15) rectangle (0.08,0.15);
\fill[popmuted] (-0.08,0.45) rectangle (0.08,2.2);
\node[left, text=popdark] at (-0.1,0.3) {$S_1$};
\node[left, text=popdark] at (-0.1,-0.3) {$S_2$};
% slit separation a
\draw[<->, poporange, thick] (0.35,-0.3) -- (0.35,0.3) node[midway, right, text=poporange]{$a$};
% screen
\fill[popmuted] (6,-2.2) rectangle (6.16,2.2);
\node[below, text=popdark] at (6.08,-2.25) {screen};
% distance D
\draw[<->, popblue, thick] (0,-2.0) -- (6,-2.0) node[midway, below, text=popblue]{$D$};
% central axis
\draw[dashed, gray] (0,0) -- (6,0);
% rays to point P
\coordinate (P) at (6,1.3);
\draw[popgreen, thick] (0,0.3) -- (P);
\draw[popgreen, thick] (0,-0.3) -- (P);
\fill[popgreen] (P) circle (0.07) node[right, text=popgreen]{$P$};
% x on screen
\draw[<->, poppurple, thick] (6.4,0) -- (6.4,1.3) node[midway, right, text=poppurple]{$x$};
% path difference construction
\draw[poporange, thick] (0,-0.3) -- (0.55,0.19);
\draw[poporange, dashed] (0,0.3) -- (0.55,0.19);
\draw[poporange] (0.28,0.24) arc[start angle=120, end angle=180, radius=0.35];
\node[text=poporange, font=\small] at (1.15,-0.05) {$\delta \approx a\sin\theta$};
% angle theta
\draw[popblue] (1.6,0) arc[start angle=0, end angle=12.2, radius=1.6];
\node[text=popblue, font=\small] at (1.95,0.22) {$\theta$};
\end{tikzpicture}
\end{center}
\legende{Geometry of Young's double-slit experiment. The slits are separated by $a$; the screen is at distance $D \gg a$. For a point $P$ at position $x$, the path difference is $\delta = a\sin\theta \approx \dfrac{ax}{D}$.}
\end{popfigure}

\textbf{Derivation of the fringe spacing (examinable).} Let $a = S_1S_2$ be the slit separation and $D$ the slit-to-screen distance, with $D \gg a$. For a point $P$ at distance $x$ from the central axis, the ray from $S_2$ travels an extra distance
\[
\delta = a\sin\theta.
\]
Since $D \gg a$, the angle $\theta$ is small, so $\sin\theta \approx \tan\theta = \dfrac{x}{D}$, giving
\[
\delta \approx \dfrac{ax}{D}.
\]
Bright fringes occur where $\delta = k\lambda$, i.e.\ at positions $x_k = \dfrac{k\lambda D}{a}$. The distance between two consecutive bright (or dark) fringes — the \cle{fringe spacing} or \emph{interfringe} — is therefore constant:

\begin{propriete}[Fringe spacing for Young's double slit (proof examinable)]
\[
\boxed{i = \dfrac{\lambda D}{a}}
\]
where $i$ is the fringe spacing (m), $\lambda$ the wavelength (m), $D$ the slit-to-screen distance (m) and $a$ the slit separation (m). The fringes are equally spaced and parallel to the slits; the central fringe ($k=0$) is bright.
\end{propriete}

Note the physical meaning: $i$ grows with $\lambda$ (red light gives wider fringes than blue), grows with $D$, and shrinks when the slits are moved apart.

\begin{attention}[Do not confuse $a$ with the slit width]
The most frequent error at the GCE: $a$ is the \textbf{separation between the centres of the two slits}, not the width of one slit. A second classic error is forgetting to convert units: $a$ and $i$ are usually given in millimetres and must be converted to metres before applying $i = \lambda D/a$. Finally, $D$ is measured from the \emph{slits} to the screen, not from the source slit $S$.
\end{attention}

\begin{exoresolu}[Computing a wavelength from Young's fringes]
In a Young's double-slit experiment, the slits are separated by $a = 0.40\ \mathrm{mm}$ and the screen is placed at $D = 1.50\ \mathrm{m}$. The distance between the centres of the first and eleventh bright fringes (10 fringe spacings) is measured as $22.5\ \mathrm{mm}$. Calculate the wavelength of the light used.
\tcblower
\textbf{Solution.} The measured width corresponds to 10 interfringes, so
\[
i = \dfrac{22.5\ \mathrm{mm}}{10} = 2.25\ \mathrm{mm} = 2.25\times 10^{-3}\ \mathrm{m}.
\]
Converting: $a = 0.40\ \mathrm{mm} = 4.0\times 10^{-4}\ \mathrm{m}$. From $i = \dfrac{\lambda D}{a}$:
\[
\lambda = \dfrac{ia}{D} = \dfrac{(2.25\times 10^{-3})(4.0\times 10^{-4})}{1.50} = 6.0\times 10^{-7}\ \mathrm{m} = 600\ \mathrm{nm}.
\]
This is orange light — a sensible value, since visible wavelengths lie roughly between $400\ \mathrm{nm}$ and $750\ \mathrm{nm}$.
\end{exoresolu}

\courssub{Experimental Measurement of Wavelength with the Double Slit}

\begin{methode}[Measuring $\lambda$ with Young's double slit]
\begin{enumerate}
\item Set up the source (laser or filtered lamp), the single slit, the double slit and the screen on an optical bench; align them so the fringes are sharp and horizontal.
\item Measure the slit-to-screen distance $D$ with a metre rule.
\item Using a ruler or a travelling microscope, measure the total width $L$ occupied by $N$ consecutive fringe spacings (take $N$ as large as possible, e.g.\ $N = 10$). This means measuring the distance between the centres of the first and the $(N+1)$-th bright fringe.
\item Compute the mean fringe spacing: $i = \dfrac{L}{N}$.
\item Compute the wavelength: $\lambda = \dfrac{ia}{D}$.
\item Repeat for several values of $D$ and average, or plot $i$ against $D$: the slope of the straight line is $\lambda/a$.
\end{enumerate}
\end{methode}

\textbf{Precautions and uncertainties.}
\begin{itemize}
\item Measuring $N$ fringes at once divides the reading error on $i$ by $N$ — this is why we never measure a single fringe.
\item The uncertainty on $\lambda$ combines the relative uncertainties: $\dfrac{\Delta\lambda}{\lambda} = \dfrac{\Delta i}{i} + \dfrac{\Delta a}{a} + \dfrac{\Delta D}{D}$.
\item Ensure the room is dark and the fringes are viewed perpendicular to the screen to avoid parallax.
\end{itemize}

\courssub{Multiple Slits: The Diffraction Grating}

What happens if we use not two slits but $N$ identical, equally spaced slits? Each slit contributes a wave; at a direction $\theta$, consecutive slits differ in path by $d\sin\theta$, where $d$ is the \cle{grating spacing} (distance between the centres of adjacent slits). All $N$ waves add constructively in directions where every pair is in phase:

\begin{propriete}[Grating condition for principal maxima (proof examinable)]
For a grating of spacing $d$ illuminated at normal incidence, the principal maxima occur at angles $\theta_k$ such that
\[
\boxed{d\sin\theta_k = k\lambda, \qquad k = 0, 1, 2, \dots}
\]
$k$ is called the \textbf{order} of the maximum. The central maximum ($k=0$) lies on the axis for all wavelengths.
\end{propriete}

As $N$ increases, two things happen: the principal maxima become \textbf{brighter} (amplitude $N$ times that of one slit, intensity $N^2$ times) and \textbf{much sharper}, because between two principal maxima the $N$ waves cancel each other almost completely. A grating with thousands of lines therefore produces extremely fine, intense, well-separated bright lines — ideal for precise wavelength measurement.

\begin{popfigure}
\begin{center}
\begin{tikzpicture}
\begin{axis}[
  width=12cm, height=6cm,
  xlabel={position on screen $x$ (arbitrary units)},
  ylabel={intensity $I$},
  xmin=-6, xmax=6, ymin=0, ymax=1.25,
  xtick=\empty, ytick=\empty,
  axis lines=left,
  legend style={at={(0.02,0.95)}, anchor=north west, draw=none, font=\small}
]
% double slit: cos^2 pattern, equally spaced
\addplot[popblue, very thick, domain=-6:6, samples=400]
  {cos(deg(3.14159*x/1.5))^2};
\addlegendentry{double slit: broad, equally spaced fringes}
% grating: sharp maxima (approximate with high power)
\addplot[poporange, very thick, domain=-6:6, samples=600]
  {cos(deg(3.14159*x/1.5))^16};
\addlegendentry{grating ($N$ large): sharp, bright principal maxima}
\end{axis}
\end{tikzpicture}
\end{center}
\legende{Intensity patterns. The double slit gives broad $\cos^2$ fringes of equal spacing $i = \lambda D/a$; a grating with many slits concentrates the same energy into very narrow, intense principal maxima at the same positions $d\sin\theta_k = k\lambda$.}
\end{popfigure}

\textbf{Measuring $\lambda$ with a grating.} Shine the monochromatic beam normally onto the grating, measure on the screen the distance $y_k$ of the $k$-th order spot from the central spot and the grating-to-screen distance $L$. Then $\tan\theta_k = y_k/L$, and
\[
\lambda = \dfrac{d\sin\theta_k}{k}.
\]
Because the maxima are so sharp, their positions can be located very precisely — this is why gratings, not double slits, are used in real spectrometers. The grating spacing is usually quoted as a number of lines per millimetre $n$; then $d = \dfrac{1}{n}$ (remember to convert: $n = 600$ lines/mm gives $d = \dfrac{1}{600}\ \mathrm{mm} = 1.67\times 10^{-6}\ \mathrm{m}$).

\begin{savaistu}[Gratings around us]
A CD or DVD acts as a reflection grating: its spiral track spacing is of the order of a micrometre, comparable to the wavelength of light, which is why it splits white light into rainbow colours. The same physics explains the iridescent colours of some butterfly wings and of the feathers of the sunbirds common in Cameroonian gardens.
\end{savaistu}

\courssub{Synthesis: Single Slit, Double Slit, Grating}

\begin{aretenir}[Comparing the three patterns — a classic GCE synthesis]
\begin{center}
\begin{tabularx}{\linewidth}{|l|X|X|X|}
\hline
\rowcolor{popblueL}
 & \textbf{Single slit} & \textbf{Double slit} & \textbf{Grating ($N$ slits)} \\
\hline
Phenomenon & Diffraction & Interference (2 waves) & Interference ($N$ waves) \\
\hline
Pattern & Wide central maximum, weaker secondary maxima, unequal spacing & Equally spaced fringes of nearly equal brightness & Very sharp, very bright principal maxima, widely separated \\
\hline
Condition & Minima: $b\sin\theta = k\lambda\ (k=\pm1,\pm2,\dots)$ ($b$ = slit width) & Bright: $\delta = k\lambda$; dark: $\delta = (2k+1)\lambda/2$; $i = \lambda D/a$ & Maxima: $d\sin\theta = k\lambda$ \\
\hline
Use & Shows light bends around apertures & Measures $\lambda$ (school lab) & Precise spectroscopy, measuring $\lambda$ \\
\hline
\end{tabularx}
\end{center}
\end{aretenir}

\begin{exercice}[Fringes and grating]
A laser beam falls normally on a double slit with $a = 0.25\ \mathrm{mm}$; the screen is at $D = 2.00\ \mathrm{m}$ and the fringe spacing is $i = 5.2\ \mathrm{mm}$.
\begin{enumerate}
\item Calculate the wavelength $\lambda$ of the laser.
\item The same laser illuminates a grating of $500$ lines per millimetre at normal incidence. Calculate the angle of the first-order maximum.
\item State two differences between the double-slit pattern and the grating pattern.
\end{enumerate}
\end{exercice}

\begin{corrige}[Exercise 1]
\textbf{1.} $\lambda = \dfrac{ia}{D} = \dfrac{(5.2\times 10^{-3})(2.5\times 10^{-4})}{2.00} = 6.5\times 10^{-7}\ \mathrm{m} = 650\ \mathrm{nm}$ (red laser).\par
\textbf{2.} $d = \dfrac{1}{500}\ \mathrm{mm} = 2.0\times 10^{-6}\ \mathrm{m}$. From $d\sin\theta_1 = \lambda$: $\sin\theta_1 = \dfrac{6.5\times 10^{-7}}{2.0\times 10^{-6}} = 0.325$, so $\theta_1 \approx 19^{\circ}$.\par
\textbf{3.} The grating maxima are much sharper and much brighter than the double-slit fringes; for a typical grating with many lines per mm, the maxima are more widely spaced angularly than the double-slit fringes.
\end{corrige}

\begin{aretenir}[Key points of this section]
\begin{itemize}
\item Superposition: the resultant displacement is the vector sum of the individual displacements.
\item Observable interference requires \textbf{coherent} sources: same frequency, constant phase difference.
\item Constructive: $\delta = k\lambda$; destructive: $\delta = (2k+1)\lambda/2$.
\item Young's double slit: equally spaced fringes with $i = \dfrac{\lambda D}{a}$; measure $N$ fringes to reduce uncertainty, then $\lambda = \dfrac{ia}{D}$.
\item Grating: principal maxima at $d\sin\theta = k\lambda$; more slits $\Rightarrow$ sharper, brighter maxima $\Rightarrow$ more precise wavelength measurement.
\end{itemize}
\end{aretenir}

\coursec{Doppler Effect and Stationary Waves}

Having explored how waves interfere when they meet, we now turn to two fascinating phenomena that arise from the relative motion of sources and observers, and from the superposition of waves travelling in opposite directions. The \emph{Doppler effect} explains why the pitch of a siren changes as an ambulance speeds past you on the streets of Yaoundé, while \emph{stationary waves} are the secret behind the music produced by a guitar string or a flute. Both topics are favourites in GCE examinations and require a clear grasp of sign conventions and physical interpretation.

\courssub{The Doppler Effect: Apparent Change of Frequency}

\begin{definition}[The Doppler Effect]
The \cle{Doppler effect} is the apparent change in the frequency (and wavelength) of a wave as perceived by an observer when there is relative motion between the source of the wave and the observer. The actual frequency emitted by the source does not change; only the frequency received by the observer changes.
\end{definition}

\textbf{Everyday examples.}\par
You have certainly heard the sudden drop in pitch of a motorbike's horn as it overtakes you on the highway, or the changing note of a train whistle as it approaches and then recedes from a platform. In Cameroon, the siren of a \textsc{samu} ambulance speeding through traffic provides a perfect real-life illustration: the pitch sounds higher as it approaches and lower once it has passed.

\textbf{Qualitative explanation with wavefronts.}\par
Consider a source emitting spherical wavefronts at a constant frequency $f$ in a medium where the wave speed is $v$. If the source is stationary, the wavefronts are concentric circles spaced by the wavelength $\lambda = v/f$. When the source moves towards a stationary observer, each successive wavefront is emitted from a position closer to the observer. The wavefronts bunch up in front of the source (shorter effective wavelength) and spread out behind it (longer effective wavelength). The observer in front therefore receives more wavefronts per second $\rightarrow$ higher frequency; the observer behind receives fewer $\rightarrow$ lower frequency.

\begin{popfigure}
\begin{tikzpicture}[scale=0.7]
% Moving source Doppler wavefronts (source moves to the right)
\draw[popblue, thick] (-2.16,0) circle (3.6);
\draw[popblue, thick] (-1.44,0) circle (2.4);
\draw[popblue, thick] (-0.72,0) circle (1.2);
\fill[poporange] (-2.16,0) circle (0.08);
\fill[poporange] (-1.44,0) circle (0.08);
\fill[poporange] (-0.72,0) circle (0.08);
% Current source position (t=0)
\fill[popdark] (0,0) circle (0.1);
\node[popdark, above] at (0,0.3) {Source at $t=0$};
% Observer positions
\node[popgreen, below] at (4.5,-0.5) {Observer A (approach)};
\node[popgreen, below] at (-5,-0.5) {Observer B (recession)};
% Arrow for source velocity
\draw[->, thick, popdark] (-0.5,1.5) -- (0.5,1.5) node[midway, above] {$\vec{v}_s$};
% Wavelength labels
\draw[<->, popmuted] (1.2,0) -- (2.4,0) node[midway, above, popmuted] {$\lambda' < \lambda$};
\draw[<->, popmuted] (-3.84,0) -- (-2.64,0) node[midway, above, popmuted] {$\lambda' > \lambda$};
\end{tikzpicture}
\legende{Wavefronts emitted by a source moving to the right. Wavefronts are compressed ahead (shorter $\lambda'$) and stretched behind (longer $\lambda'$).}
\end{popfigure}

\textbf{Moving source, stationary observer.}\par
Let the source move with speed $v_s$ towards a stationary observer. In one period $T = 1/f$, the source travels a distance $v_s T$. The wavelength in front of the source is reduced to $\lambda' = \lambda - v_s T = \frac{v}{f} - \frac{v_s}{f} = \frac{v - v_s}{f}$. The observer, stationary in the medium, receives waves travelling at speed $v$, so the observed frequency is
\[
f' = \frac{v}{\lambda'} = \frac{v}{(v - v_s)/f} = f\,\frac{v}{v - v_s}.
\]
If the source moves \emph{away} from the observer, the wavelength is stretched: $\lambda' = \lambda + v_s T$, giving
\[
f' = f\,\frac{v}{v + v_s}.
\]
These two cases are combined into one formula using the upper sign for approach and the lower for recession:
\[
\boxed{f' = f\,\frac{v}{v \mp v_s}}
\]

\textbf{Moving observer, stationary source.}\par
Now the source is stationary, emitting wavelength $\lambda = v/f$. The observer moves with speed $v_o$ towards the source. The relative speed between the observer and the wavefronts is $v + v_o$ (approach) or $v - v_o$ (recession). The observed frequency is the rate at which the observer crosses wavefronts:
\[
f' = \frac{v \pm v_o}{\lambda} = \frac{v \pm v_o}{v/f} = f\,\frac{v \pm v_o}{v}.
\]
Upper sign for observer moving towards the source, lower sign for moving away.

\begin{propriete}[Doppler Effect Formulae for Sound in Air]
Let $v$ be the speed of sound in air ($\approx \SI{340}{m.s^{-1}}$ at room temperature), $f$ the source frequency, $v_s$ the speed of the source, and $v_o$ the speed of the observer.
\begin{itemize}
\item \textbf{Moving source, stationary observer:} $f' = f\,\dfrac{v}{v \mp v_s}$ \quad ($-$ for approach, $+$ for recession).
\item \textbf{Moving observer, stationary source:} $f' = f\,\dfrac{v \pm v_o}{v}$ \quad ($+$ for approach, $-$ for recession).
\end{itemize}
\textbf{Note:} These formulae assume the motion is along the line joining source and observer. If the motion is at an angle, the component of velocity along that line must be used.
\end{propriete}

\begin{methode}[Choosing the Correct Sign — The ``Approach = Higher Frequency'' Rule]
\begin{enumerate}
\item Identify whether the source or the observer is moving (or both).
\item Ask: \emph{Are the source and observer getting closer or moving apart?}
\item If they \textbf{approach} each other, the observed frequency \textbf{increases} ($f' > f$). Choose the sign that makes the fraction larger than 1.
\item If they \textbf{recede} from each other, the observed frequency \textbf{decreases} ($f' < f$). Choose the sign that makes the fraction smaller than 1.
\item For a moving source: denominator $v \mp v_s$; approach $\rightarrow$ minus (smaller denominator $\rightarrow$ larger $f'$).
\item For a moving observer: numerator $v \pm v_o$; approach $\rightarrow$ plus (larger numerator $\rightarrow$ larger $f'$).
\end{enumerate}
\end{methode}

\begin{attention}[Common Mistake: Confusing $v$ and $v_s$]
In the Doppler formulae, \cle{$v$ is the speed of the wave in the medium} (speed of sound in air, $\approx \SI{340}{m.s^{-1}}$), \emph{not} the speed of the source. Students often mistakenly substitute the source speed for $v$. Always write down the known values with their symbols before plugging into the formula.
\end{attention}

\begin{exoresolu}[Ambulance Siren — Approach and Recession]
An ambulance sounds its siren at a frequency of $\SI{800}{Hz}$ while travelling at $\SI{30}{m.s^{-1}}$ on a straight road. A pedestrian stands by the roadside. Take the speed of sound in air as $\SI{340}{m.s^{-1}}$.
\begin{enumerate}
\item Calculate the frequency heard by the pedestrian as the ambulance \emph{approaches}.
\item Calculate the frequency heard as the ambulance \emph{recedes} after passing.
\item What is the beat frequency heard by the pedestrian if the ambulance passes very close? (Assume the transition is instantaneous for this calculation.)
\end{enumerate}
\tcblower
\begin{itemize}
\item \textbf{Approach:} Source moving towards stationary observer. $f' = f\,\dfrac{v}{v - v_s} = 800 \times \dfrac{340}{340 - 30} = 800 \times \dfrac{340}{310} \approx \SI{877.4}{Hz}$.
\item \textbf{Recession:} Source moving away from stationary observer. $f' = f\,\dfrac{v}{v + v_s} = 800 \times \dfrac{340}{340 + 30} = 800 \times \dfrac{340}{370} \approx \SI{735.1}{Hz}$.
\item \textbf{Beat frequency:} $f_{\text{beat}} = |f'_{\text{approach}} - f'_{\text{recede}}| \approx 877.4 - 735.1 = \SI{142.3}{Hz}$.
\end{itemize}
\textbf{Comment:} The pedestrian hears a distinct drop in pitch of about 142 Hz as the ambulance passes. This large change is typical for fast-moving sources.
\end{exoresolu}

\begin{exoresolu}[Moving Observer — Motorcyclist and Factory Siren]
A factory siren emits a steady note of $\SI{500}{Hz}$. A motorcyclist rides towards the factory at $\SI{20}{m.s^{-1}}$, then turns around and rides away at the same speed. Speed of sound $v = \SI{340}{m.s^{-1}}$. Find the frequencies heard in both cases.
\tcblower
\begin{itemize}
\item \textbf{Towards factory (approach):} Observer moving towards stationary source. $f' = f\,\dfrac{v + v_o}{v} = 500 \times \dfrac{340 + 20}{340} = 500 \times \dfrac{360}{340} \approx \SI{529.4}{Hz}$.
\item \textbf{Away from factory (recession):} Observer moving away from stationary source. $f' = f\,\dfrac{v - v_o}{v} = 500 \times \dfrac{340 - 20}{340} = 500 \times \dfrac{320}{340} \approx \SI{470.6}{Hz}$.
\end{itemize}
\textbf{Note:} The shift is smaller than in the moving-source case for the same speed because the wavelength in the medium is unchanged; only the encounter rate changes.
\end{exoresolu}

\courssub{Stationary Waves: Formation and Characteristics}

\begin{definition}[Stationary (Standing) Wave]
A \cle{stationary wave} is formed by the superposition of two progressive waves of the same frequency, wavelength, and amplitude travelling in opposite directions. Unlike a progressive wave, a stationary wave does not transport energy along the medium; the energy is confined between fixed points.
\end{definition}

\textbf{Formation.}\par
The most common way to produce a stationary wave is to send a progressive wave towards a fixed boundary (e.g., a string tied to a wall). The wave reflects with a phase reversal (a crest returns as a trough). The incident and reflected waves superpose. At certain points the two waves always cancel $\rightarrow$ \cle{nodes} (zero amplitude). At points midway between nodes, they always reinforce $\rightarrow$ \cle{antinodes} (maximum amplitude).

\begin{popfigure}
\begin{tikzpicture}[scale=0.9]
% Stationary wave snapshots at different times
% t = 0 (maximum displacement)
\draw[popline, thin] (0,0) -- (6,0);
\draw[popblue, thick, domain=0:6, samples=100, smooth] plot (\x, {1.2*sin(60*\x)});
\fill[popdark] (0,0) circle (0.07);
\fill[popdark] (3,0) circle (0.07);
\fill[popdark] (6,0) circle (0.07);
\node[popdark, below] at (0,-0.3) {N};
\node[popdark, below] at (3,-0.3) {N};
\node[popdark, below] at (6,-0.3) {N};
\node[poporange, above] at (1.5,1.4) {A};
\node[poporange, above] at (4.5,1.4) {A};
\draw[<->, popmuted] (0,-0.8) -- (3,-0.8) node[midway, below, popmuted] {$\lambda/2$};
\draw[<->, popmuted] (3,-0.8) -- (6,-0.8) node[midway, below, popmuted] {$\lambda/2$};
\node[popdark, left] at (-0.5,0) {$t=0$};
% t = T/4 (flat)
\draw[popline, thin] (0,-3) -- (6,-3);
\draw[popblue, thick] (0,-3) -- (6,-3);
\node[popdark, left] at (-0.5,-3) {$t=T/4$};
% t = T/2 (inverted)
\draw[popline, thin] (0,-6) -- (6,-6);
\draw[popblue, thick, domain=0:6, samples=100, smooth] plot (\x, {-1.2*sin(60*\x)-6});
\node[popdark, left] at (-0.5,-6) {$t=T/2$};
% t = 3T/4 (flat)
\draw[popline, thin] (0,-9) -- (6,-9);
\draw[popblue, thick] (0,-9) -- (6,-9);
\node[popdark, left] at (-0.5,-9) {$t=3T/4$};
% Legend
\node[popdark, align=center, text width=3.2cm] at (8.2,-4.5) {N = Node (zero amplitude)\\A = Antinode (max amplitude)};
\end{tikzpicture}
\legende{Snapshots of a stationary wave on a string at four instants. Nodes (N) are fixed; antinodes (A) oscillate with maximum amplitude. Distance between adjacent nodes = $\lambda/2$.}
\end{popfigure}

\begin{propriete}[Characteristics of a Stationary Wave]
\begin{itemize}
\item \textbf{Nodes (N):} Points of zero amplitude at all times. Adjacent nodes are separated by $\lambda/2$.
\item \textbf{Antinodes (A):} Points of maximum amplitude. Adjacent antinodes are separated by $\lambda/2$. A node and the next antinode are $\lambda/4$ apart.
\item \textbf{Phase:} All points between two adjacent nodes vibrate \textbf{in phase}. Points on opposite sides of a node vibrate in \textbf{antiphase} (phase difference $\pi$).
\item \textbf{Amplitude:} Varies with position: $A(x) = 2A_0 \sin(kx)$ (for a string fixed at $x=0$). Maximum at antinodes, zero at nodes.
\item \textbf{Energy:} No net energy transport along the medium. Energy oscillates between kinetic and potential forms within each loop.
\end{itemize}
\end{propriete}

\begin{attention}[Stationary vs. Progressive Waves — Key Differences]
\begin{center}
\begin{tabularx}{\linewidth}{|l|X|X|}\hline
\rowcolor{popblueL}
\textbf{Feature} & \textbf{Progressive Wave} & \textbf{Stationary Wave} \\ \hline
Energy transport & Yes, in direction of propagation & No net transport \\ \hline
Amplitude & Same for all points (in ideal medium) & Varies with position (nodes to antinodes) \\ \hline
Phase & Changes continuously with distance & Constant between nodes; flips by $\pi$ at each node \\ \hline
Wave profile & Moves along the medium & Stays fixed; only amplitude oscillates \\ \hline
\end{tabularx}
\end{center}
\textbf{Exam trap:} Never describe a stationary wave as ``travelling'' or say it transfers energy from one end to the other.
\end{attention}

\courssub{Stationary Waves on a Stretched String: Harmonics}

When a string of length $L$ is fixed at both ends, stationary waves can only form if the boundary conditions are satisfied: nodes at both ends. This restricts the possible wavelengths to
\[
L = n\,\frac{\lambda_n}{2} \quad (n = 1, 2, 3, \dots)
\]
where $n$ is the \cle{harmonic number} (or mode number). The allowed wavelengths and frequencies are
\[
\lambda_n = \frac{2L}{n}, \qquad f_n = \frac{v}{\lambda_n} = n\,\frac{v}{2L} = n f_1.
\]
The fundamental frequency ($n=1$) is $f_1 = \dfrac{v}{2L}$. Higher harmonics are integer multiples of the fundamental.

\begin{propriete}[Frequencies of a String Fixed at Both Ends]
For a string of length $L$, mass per unit length $\mu$, and tension $T$, the wave speed is $v = \sqrt{T/\mu}$. The natural frequencies (harmonics) are
\[
\boxed{f_n = \frac{n}{2L}\sqrt{\frac{T}{\mu}} \qquad (n = 1, 2, 3, \dots)}
\]
The fundamental ($n=1$) is also called the \emph{first harmonic}; $n=2$ is the second harmonic (first overtone), etc.
\end{propriete}

\begin{popfigure}
\begin{tikzpicture}[scale=0.7]
% First harmonic n = 1
\draw[popline, thin] (0,0) -- (8,0);
\draw[popblue, thick, domain=0:8, samples=100, smooth] plot (\x, {0.8*sin(22.5*\x)});
\draw[popblue, dashed, domain=0:8, samples=100, smooth] plot (\x, {-0.8*sin(22.5*\x)});
\fill[popdark] (0,0) circle (0.07);
\fill[popdark] (8,0) circle (0.07);
\node[poporange] at (4,0) {$\bullet$};
\node[popdark, left] at (-0.5,0) {$n=1$ (fundamental)};
\draw[<->, popmuted] (0,-1.2) -- (8,-1.2) node[midway, below, popmuted] {$\lambda_1 = 2L$};
% Second harmonic n = 2
\draw[popline, thin] (0,-3) -- (8,-3);
\draw[popblue, thick, domain=0:8, samples=100, smooth] plot (\x, {0.8*sin(45*\x)-3});
\draw[popblue, dashed, domain=0:8, samples=100, smooth] plot (\x, {-0.8*sin(45*\x)-3});
\fill[popdark] (0,-3) circle (0.07);
\fill[popdark] (4,-3) circle (0.07);
\fill[popdark] (8,-3) circle (0.07);
\node[poporange] at (2,-3) {$\bullet$};
\node[poporange] at (6,-3) {$\bullet$};
\node[popdark, left] at (-0.5,-3) {$n=2$};
\draw[<->, popmuted] (0,-4.2) -- (8,-4.2) node[midway, below, popmuted] {$\lambda_2 = L$};
% Third harmonic n = 3
\draw[popline, thin] (0,-6) -- (8,-6);
\draw[popblue, thick, domain=0:8, samples=100, smooth] plot (\x, {0.8*sin(67.5*\x)-6});
\draw[popblue, dashed, domain=0:8, samples=100, smooth] plot (\x, {-0.8*sin(67.5*\x)-6});
\fill[popdark] (0,-6) circle (0.07);
\fill[popdark] (2.667,-6) circle (0.07);
\fill[popdark] (5.333,-6) circle (0.07);
\fill[popdark] (8,-6) circle (0.07);
\node[poporange] at (1.333,-6) {$\bullet$};
\node[poporange] at (4,-6) {$\bullet$};
\node[poporange] at (6.667,-6) {$\bullet$};
\node[popdark, left] at (-0.5,-6) {$n=3$};
\draw[<->, popmuted] (0,-7.2) -- (5.333,-7.2) node[midway, below, popmuted] {$\lambda_3 = 2L/3$};
\end{tikzpicture}
\legende{First three harmonics of a string fixed at both ends. Nodes (black dots) at the ends and interior; antinodes (orange dots) at maximum displacement. $\lambda_n = 2L/n$. Distance between adjacent nodes = $\lambda_n/2$.}
\end{popfigure}

\begin{experience}[Sonometer / Melde's Experiment — Measuring Wavelength or Frequency]
\textbf{Apparatus:} A long wire stretched over two bridges on a wooden box (sonometer), tension adjusted by hanging masses, a tuning fork of known frequency (or a signal generator with a vibrator).
\textbf{Procedure:}
\begin{enumerate}
\item Set the tension $T$ by hanging a known mass $m$ ($T = mg$).
\item Place a small paper rider on the wire at an antinode position.
\item Excite the wire (bow it, or use a vibrating tuning fork placed on the bridge, or a signal generator driving a vibrator attached to the wire).
\item Adjust the distance $L$ between the bridges until the rider is thrown off violently $\rightarrow$ resonance at the fundamental (or a harmonic).
\item Measure $L$ for the fundamental. Then $\lambda = 2L$. Knowing $f$ of the tuning fork, $v = f\lambda$. Alternatively, vary $T$ and plot $f^2$ vs $T$ to find $\mu$.
\end{enumerate}
\textbf{Exam relevance:} You may be asked to describe the experiment, explain why the rider falls at resonance, or derive the relationship between $f$, $L$, $T$, and $\mu$.
\end{experience}

\begin{savaistu}[Musical Instruments and Resonance in Pipes]
The principles of stationary waves on strings apply directly to \textbf{string instruments} (guitar, violin, harp, \emph{mvet} — the traditional Cameroonian harp-lute). The player changes the effective length $L$ by pressing the string against the fingerboard, thus changing the pitch ($f \propto 1/L$). Tuning is done by adjusting tension $T$ ($f \propto \sqrt{T}$).
For \textbf{wind instruments} (flute, trumpet, \emph{algaita}), stationary waves form in air columns. The boundary conditions differ: an open end is an antinode (pressure node), a closed end is a node (pressure antinode). This leads to different harmonic series (only odd harmonics for a closed pipe). While pipe resonance is an extension beyond the core string syllabus, it frequently appears in GCE ``applications'' questions.
\end{savaistu}

\begin{exoresolu}[Guitar String — Harmonics and Finger Placement]
A guitar string of length $\SI{0.65}{m}$ and mass per unit length $\SI{5.0e-4}{kg.m^{-1}}$ is under a tension of $\SI{80}{N}$.
\begin{enumerate}
\item Calculate the fundamental frequency.
\item The guitarist presses the string at the 12th fret, exactly halfway along the string. What is the new fundamental frequency?
\item If the string is lightly touched at the midpoint (without pressing to the fretboard) and plucked, what frequency is heard? (This produces a harmonic.)
\end{enumerate}
\tcblower
\begin{enumerate}
\item Wave speed $v = \sqrt{T/\mu} = \sqrt{80 / (5.0\times10^{-4})} = \sqrt{1.6\times10^5} = \SI{400}{m.s^{-1}}$.
  Fundamental $f_1 = \dfrac{v}{2L} = \dfrac{400}{2 \times 0.65} = \dfrac{400}{1.3} \approx \SI{307.7}{Hz}$.
\item Pressing at the 12th fret makes the vibrating length $L' = L/2 = \SI{0.325}{m}$. New fundamental $f_1' = \dfrac{v}{2L'} = \dfrac{400}{2 \times 0.325} = \dfrac{400}{0.65} \approx \SI{615.4}{Hz}$ (one octave higher).
\item Lightly touching the midpoint forces a node there. The string vibrates in its second harmonic ($n=2$) with the full length $L$. Frequency $f_2 = 2 f_1 = \SI{615.4}{Hz}$. The pitch is the same as in (ii), but the timbre is different (no fundamental component).
\end{enumerate}
\end{exoresolu}

\begin{exercice}[Doppler Effect and Stationary Waves Practice]
\begin{enumerate}
\item A police car siren emits a frequency of $\SI{600}{Hz}$. The car moves at $\SI{25}{m.s^{-1}}$ towards a stationary observer. Speed of sound $v = \SI{340}{m.s^{-1}}$. Calculate the observed frequency.
\item An observer moves at $\SI{10}{m.s^{-1}}$ away from a stationary source emitting $\SI{1000}{Hz}$. Find the observed frequency.
\item A string of length $\SI{1.2}{m}$ fixed at both ends vibrates in its third harmonic. If the wave speed on the string is $\SI{240}{m.s^{-1}}$, find the frequency of this vibration.
\item In a stationary wave on a string, the distance between two adjacent nodes is $\SI{0.15}{m}$. What is the wavelength? If the frequency is $\SI{50}{Hz}$, what is the wave speed?
\item \textbf{Challenge:} A source emits frequency $f$ and moves at speed $v_s$ towards a stationary wall. The sound reflects off the wall and returns to the source. What is the frequency of the reflected wave as ``heard'' by the source? (Hint: treat the wall as a stationary observer that becomes a stationary source re-emitting the frequency it receives.)
\end{enumerate}
\end{exercice}

\begin{corrige}[Exercise n]
\begin{enumerate}
\item Moving source, stationary observer, approach: $f' = 600 \times \frac{340}{340-25} = 600 \times \frac{340}{315} \approx \SI{647.6}{Hz}$.
\item Moving observer, stationary source, recession: $f' = 1000 \times \frac{340-10}{340} = 1000 \times \frac{330}{340} \approx \SI{970.6}{Hz}$.
\item Third harmonic $\rightarrow n=3$. $f_3 = 3 \times \frac{v}{2L} = 3 \times \frac{240}{2 \times 1.2} = 3 \times 100 = \SI{300}{Hz}$.
\item Node-to-node distance $= \lambda/2 = 0.15 \Rightarrow \lambda = \SI{0.30}{m}$. $v = f\lambda = 50 \times 0.30 = \SI{15}{m.s^{-1}}$.
\item \textbf{Step 1:} Wall receives frequency $f_1 = f \frac{v}{v - v_s}$ (source approaches wall). \textbf{Step 2:} Wall reflects this frequency $f_1$ as a stationary source. Source now moves towards this stationary source (the wall) with speed $v_s$. Frequency heard by moving observer (the original source) approaching a stationary source: $f_2 = f_1 \frac{v + v_s}{v} = f \frac{v}{v - v_s} \frac{v + v_s}{v} = f \frac{v + v_s}{v - v_s}$.
\end{enumerate}
\end{corrige}

\begin{aretenir}[Key Points for GCE Revision]
\begin{itemize}
\item \textbf{Doppler effect:} Apparent frequency change due to relative motion. Use $f' = f \frac{v}{v \mp v_s}$ (moving source) or $f' = f \frac{v \pm v_o}{v}$ (moving observer). $v$ = wave speed in medium.
\item \textbf{Sign rule:} Approach $\rightarrow$ frequency increases (choose sign to make fraction $>1$). Recession $\rightarrow$ frequency decreases (fraction $<1$).
\item \textbf{Stationary wave:} Superposition of two identical opposite waves. Nodes (zero amplitude) and antinodes (max amplitude). Node-to-node = $\lambda/2$, node-to-antinode = $\lambda/4$.
\item \textbf{No energy transport} in a stationary wave. Amplitude depends on position. Phase constant between nodes.
\item \textbf{String fixed at both ends:} $\lambda_n = 2L/n$, $f_n = n v/(2L) = n f_1$. $v = \sqrt{T/\mu}$.
\item \textbf{Sonometer:} Resonance when $L = n\lambda/2$. Used to measure $f$, $v$, or $\mu$.
\end{itemize}
\end{aretenir}

\newpage

\coursec{Integration activity}

\coursec{Integration Activity: Mechanical Waves}

\courssub{Aims of the Activity}
This integration activity is designed to consolidate your mastery of the core concepts of Mechanical Waves covered in this chapter: the \cle{wave equation} ($v = f\lambda$), \cle{interference} (double-slit geometry), \cle{stationary waves} (resonance in air columns), and the \cle{Doppler effect}. You will mobilise these concepts in a realistic laboratory setting typical of the GCE Advanced Level Practical Paper (Paper 3). The activity develops three key competencies: \textbf{experimental design and data collection} (handling apparatus, reading instruments, estimating uncertainties), \textbf{mathematical modelling} (linearising equations, graphical analysis, gradient calculation), and \textbf{critical evaluation} (discussing systematic errors, comparing with accepted values, proposing improvements). By the end of this session, you should be able to plan and execute measurements to determine the speed of sound in air and the wavelength of laser light, and to predict the outcome of variations in experimental parameters using the Doppler formula and interference conditions.

\courssub{Situation}
\textbf{Context:} You are a student in Upper Sixth at \emph{Government Bilingual High School (GBHS) Molyko, Buea}. The Physics laboratory is equipped for the GCE Practical examination. Your teacher has set up two independent stations (Station A and Station B) to be completed in a single 2-hour session. You work in groups of three, rotating between stations after one hour. A written report following the GCE Practical marking scheme (Observations, Calculations, Graph, Conclusion, Evaluation) is required for assessment.

\textbf{Station A: Determination of the Speed of Sound in Air using a Resonance Tube}
A tuning fork of known frequency $f = \SI{512}{\hertz}$ (marked on the stem) is struck gently and held just above the open end of a vertical glass tube dipped into a large cylindrical jar of water. The tube can be raised or lowered vertically, varying the length $L$ of the air column above the water surface. The water surface acts as a \cle{closed end} (displacement node, pressure antinode) and the open end acts as an approximate \cle{displacement antinode} (pressure node). When the air column resonates, a loud sound is heard. The first resonance (fundamental mode) occurs when $L_1 \approx \lambda/4$, and the second resonance (first overtone) occurs when $L_2 \approx 3\lambda/4$. The \cle{end correction} $e$ accounts for the fact that the antinode lies slightly outside the tube: $L_1 + e = \lambda/4$ and $L_2 + e = 3\lambda/4$. Subtracting eliminates $e$: $\lambda = 2(L_2 - L_1)$. The speed of sound is then $v = f\lambda = 2f(L_2 - L_1)$.

\textbf{Apparatus:} Tuning fork ($f = \SI{512}{\hertz}$), rubber hammer, metre rule ($\pm \SI{0.1}{\centi\metre}$), set squares, tall glass tube (diameter $\approx \SI{3}{\centi\metre}$), large water jar, water, thermometer ($\pm \SI{0.5}{\ensuremath{{}^{\circ}\mathrm{C}}}$).

\textbf{Station B: Determination of the Wavelength of Laser Light using a Double Slit}
A red laser pointer ($\lambda \approx \SI{650}{\nano\metre}$) illuminates a double slit of known separation $a = \SI{0.30}{\milli\metre}$ (provided by the manufacturer). The interference pattern is projected onto a vertical screen placed at a distance $D$ from the slits. Bright fringes (maxima) are observed. The fringe spacing $i$ (distance between adjacent bright fringes) is given by $i = \lambda D / a$ for small angles ($\sin\theta \approx \tan\theta = y/D$). To reduce percentage uncertainty in $i$, you measure the total width $W_N$ of $N$ fringe spacings (i.e., distance between the centres of the 1st and $(N+1)$-th bright fringes on one side of the central maximum). Then $i = W_N / N$ and $\lambda = i a / D = W_N a / (N D)$.

\textbf{Apparatus:} Red laser pointer (class II, $\lambda \approx \SI{650}{\nano\metre}$), double slit slide ($a = \SI{0.30 \pm 0.01}{\milli\metre}$), metre rule ($\pm \SI{0.1}{\centi\metre}$), screen (white paper on cardboard), optical bench or retort stands and clamps, vernier callipers ($\pm \SI{0.01}{\milli\metre}$) to verify $a$ if needed.

\textbf{Safety:} \textbf{NEVER} look directly into the laser beam or its reflection. Ensure the beam terminates on the screen. Keep water away from electrical sockets at Station A.

\courssub{Tasks}
\begin{enumerate}
    \item \textbf{Station A -- Data Collection:} 
    \begin{enumerate}
        \item Strike the tuning fork on the rubber hammer (not the bench). Hold it horizontally $\approx \SI{1}{\centi\metre}$ above the tube. Raise the tube slowly from the water until the first loud resonance is heard. Use the set square against the metre rule to read the length $L_1$ of the air column. Repeat \textbf{three times} and record $L_1$ values.
        \item Continue raising the tube to find the second resonance (next loud sound). Read $L_2$ three times. Record all values in a table with units and uncertainties.
        \item Measure and record the room temperature $\theta$ using the thermometer.
    \end{enumerate}
    \item \textbf{Station A -- Analysis:} 
    \begin{enumerate}
        \item Calculate the mean values $\overline{L_1}$ and $\overline{L_2}$ and their absolute uncertainties $\Delta L_1$, $\Delta L_2$ (use half the range or instrument precision, whichever is larger).
        \item Determine the wavelength $\lambda = 2(\overline{L_2} - \overline{L_1})$ and its absolute uncertainty $\Delta \lambda$ using simple error propagation for subtraction and multiplication by a constant.
        \item Calculate the experimental speed of sound $v_{\text{exp}} = f\lambda$ and its uncertainty $\Delta v_{\text{exp}}$.
        \item The theoretical speed of sound in dry air at temperature $\theta$ ($^\circ\mathrm{C}$) is $v_{\text{th}} = \SI{331.3}{\metre\per\second} \sqrt{1 + \frac{\theta}{273.15}}$. Compute $v_{\text{th}}$ for your measured $\theta$.
        \item Compare $v_{\text{exp}}$ and $v_{\text{th}}$. Do they agree within experimental uncertainty? Comment on the end correction $e = \lambda/4 - \overline{L_1}$.
    \end{enumerate}
    \item \textbf{Station B -- Data Collection:} 
    \begin{enumerate}
        \item Set up the laser, double slit, and screen. Adjust $D$ to be at least $\SI{2.0}{\metre}$ (measure $D$ from the slit plane to the screen). Record $D$ and its uncertainty.
        \item Observe the interference pattern. Choose a large $N$ (e.g., $N=10$ or $20$) to measure the total width $W_N$ across $N$ fringes on \textbf{one side} of the central maximum (e.g., from $m=1$ to $m=N+1$). Measure $W_N$ three times. Record $N$, $W_N$ values, and $a$.
    \end{enumerate}
    \item \textbf{Station B -- Analysis:} 
    \begin{enumerate}
        \item Calculate mean $\overline{W_N}$ and its uncertainty $\Delta W_N$.
        \item Compute the fringe spacing $i = \overline{W_N}/N$ and its uncertainty.
        \item Determine the laser wavelength $\lambda_{\text{laser}} = i a / D$ and its absolute uncertainty $\Delta \lambda_{\text{laser}}$ (propagate relative uncertainties: $\frac{\Delta \lambda}{\lambda} = \frac{\Delta i}{i} + \frac{\Delta a}{a} + \frac{\Delta D}{D}$).
        \item Express $\lambda_{\text{laser}}$ in nanometres. Compare with the typical value $\SI{650}{\nano\metre}$.
    \end{enumerate}
    \item \textbf{Synthesis Questions (Theory Application):} 
    \begin{enumerate}
        \item \textbf{Double slit modification:} If the double slit separation $a$ is doubled (using a different slide) while $D$ and $\lambda$ remain constant, describe qualitatively and quantitatively what happens to the fringe spacing $i$ and the total width $W_N$ for the same $N$. Sketch the expected intensity distribution (intensity vs position) for the two cases on the same axes, labelling key features.
        \item \textbf{Doppler Effect:} Return to Station A. Suppose the tuning fork (source) is stationary, but an observer (student) walks \textbf{towards} the fork at a constant speed $v_o = \SI{2.0}{\metre\per\second}$. The speed of sound in the lab is your $v_{\text{exp}}$. Calculate the frequency $f'$ heard by the moving observer. State the formula used.
        \item \textbf{Stationary Wave Check:} For the resonance tube at Station A, the first resonance length $L_1$ corresponds to a stationary wave pattern. Draw the displacement wave profile inside the tube at resonance, labelling the node (N) at the water surface, the antinode (A) near the open end, and indicating the distance $L_1 + e = \lambda/4$. What is the phase difference between the motion of air particles at the node and at the antinode?
    \end{enumerate}
    \item \textbf{Report Writing:} Write a concise practical report (max 2 pages) structured as: \emph{Title, Aim, Apparatus, Procedure (key steps), Observations (tables), Calculations (with uncertainty propagation), Graph (optional: $L_n$ vs $n$ for Station A to find $\lambda/2$ from gradient), Conclusion (final values with uncertainties, comparison with theory), Evaluation (two major sources of error, two improvements)}.
\end{enumerate}

\courssub{Model Answer}
\begin{exoresolu}[Model Answer -- Integration Activity: Measuring $v_{\text{sound}}$ and $\lambda_{\text{light}}$]
\textbf{Station A -- Sample Data (Simulated for Marking Scheme)}
Room temperature $\theta = \SI{27.0 \pm 0.5}{\ensuremath{{}^{\circ}\mathrm{C}}}$.
Frequency $f = \SI{512}{\hertz}$ (exact for calculation).

\begin{center}
\begin{tabular}{|c|c|c|c|}
\hline
\rowcolor{popblueL}
Resonance & Trial 1 (cm) & Trial 2 (cm) & Trial 3 (cm) \\
\hline
1st ($L_1$) & 15.8 & 16.0 & 15.9 \\
2nd ($L_2$) & 50.2 & 50.5 & 50.3 \\
\hline
\end{tabular}
\end{center}

\textbf{Analysis A}
\begin{align*}
\overline{L_1} &= \frac{15.8+16.0+15.9}{3} = \SI{15.9}{cm} = \SI{0.159}{m} \\
\Delta L_1 &= \max(\SI{0.1}{cm}, \frac{16.0-15.8}{2}\ \text{cm}) = \SI{0.1}{cm} = \SI{0.001}{m} \\
\overline{L_2} &= \frac{50.2+50.5+50.3}{3} = \SI{50.33}{cm} \approx \SI{50.3}{cm} = \SI{0.503}{m} \\
\Delta L_2 &= \max(\SI{0.1}{cm}, \frac{50.5-50.2}{2}\ \text{cm}) = \SI{0.15}{cm} \approx \SI{0.2}{cm} = \SI{0.002}{m} \\
\Delta(L_2 - L_1) &= \Delta L_1 + \Delta L_2 = \SI{0.003}{m} \quad \text{(simple sum for maximum error)} \\
\lambda &= 2(\overline{L_2} - \overline{L_1}) = 2(0.503 - 0.159) = 2(0.344) = \SI{0.688}{m} \\
\Delta \lambda &= 2 \times \Delta(L_2 - L_1) = \SI{0.006}{m} \\
\lambda &= \SI{0.688 \pm 0.006}{m} \\
v_{\text{exp}} &= f\lambda = 512 \times 0.688 = \SI{352.3}{m/s} \\
\frac{\Delta v_{\text{exp}}}{v_{\text{exp}}} &= \frac{\Delta \lambda}{\lambda} = \frac{0.006}{0.688} \approx 0.0087 \\
\Delta v_{\text{exp}} &= 352.3 \times 0.0087 \approx \SI{3.1}{m/s} \\
v_{\text{exp}} &= \SI{352 \pm 3}{m/s}
\end{align*}

\textbf{Theoretical Speed at $\theta = \SI{27}{\ensuremath{{}^{\circ}\mathrm{C}}}$}
\[ v_{\text{th}} = 331.3 \sqrt{1 + \frac{27}{273.15}} = 331.3 \sqrt{1.0988} = 331.3 \times 1.0482 = \SI{347.3}{m/s} \]

\textbf{Comparison:} $v_{\text{exp}} = \SI{352 \pm 3}{m/s}$ (range $\SI{349}{m/s}$ to $\SI{355}{m/s}$), $v_{\text{th}} = \SI{347.3}{m/s}$. The theoretical value lies \textbf{outside} the experimental uncertainty range. The values do \textbf{not} agree within the estimated experimental uncertainty. The discrepancy ($\approx \SI{5}{m/s}$) is likely due to systematic errors: the end correction approximation ($e \approx 0.3d$ assumes ideal flange), temperature gradients in the tube (air column warmer than room thermometer reading), and the subjective detection of maximum loudness by ear.

\textbf{End Correction:}
\[ e = \frac{\lambda}{4} - \overline{L_1} = \frac{0.688}{4} - 0.159 = 0.172 - 0.159 = \SI{0.013}{m} = \SI{1.3}{cm} \]
Theoretical estimate for an unflanged pipe: $e \approx 0.6r = 0.3d = 0.3 \times \SI{3}{cm} = \SI{0.9}{cm}$. The experimental value ($\SI{1.3}{cm}$) is of the same order, slightly larger, consistent with the effective length being longer than the geometric length.

\tcblower

\textbf{Station B -- Sample Data}
Slit separation $a = \SI{0.30 \pm 0.01}{mm} = \SI{3.0 \pm 0.1 \times 10^{-4}}{m}$.
Slit-to-screen distance $D = \SI{2.50 \pm 0.001}{m}$ (measured with metre rule, precision $\SI{0.1}{cm}$).
Number of fringe spacings $N = 20$ (measured from $m=1$ to $m=21$ on one side).
$W_{20}$ trials: $\SI{10.7}{cm}$, $\SI{10.9}{cm}$, $\SI{10.8}{cm}$.

\textbf{Analysis B}
\begin{align*}
\overline{W_{20}} &= \frac{10.7+10.9+10.8}{3} = \SI{10.8}{cm} = \SI{0.108}{m} \\
\Delta W_{20} &= \max(\SI{0.1}{cm}, \frac{10.9-10.7}{2}\ \text{cm}) = \SI{0.1}{cm} = \SI{0.001}{m} \\
i &= \frac{\overline{W_{20}}}{N} = \frac{0.108}{20} = \SI{5.40 \times 10^{-3}}{m} = \SI{5.40}{mm} \\
\frac{\Delta i}{i} &= \frac{\Delta W_{20}}{\overline{W_{20}}} = \frac{0.001}{0.108} \approx 0.00926 \\
\Delta i &= 5.40 \times 10^{-3} \times 0.00926 \approx \SI{5.0 \times 10^{-5}}{m} \\
\lambda_{\text{laser}} &= \frac{i a}{D} = \frac{(5.40 \times 10^{-3}) \times (3.0 \times 10^{-4})}{2.50} = \SI{6.48 \times 10^{-7}}{m} = \SI{648}{nm}
\end{align*}

\textbf{Uncertainty Propagation (Relative Uncertainties):}
\begin{align*}
\frac{\Delta a}{a} &= \frac{0.01}{0.30} = 0.0333 \quad (3.33\%) \\
\frac{\Delta D}{D} &= \frac{0.001}{2.50} = 0.0004 \quad (0.04\%) \\
\frac{\Delta i}{i} &= 0.0093 \quad (0.93\%) \\
\frac{\Delta \lambda}{\lambda} &= \frac{\Delta i}{i} + \frac{\Delta a}{a} + \frac{\Delta D}{D} \approx 0.0093 + 0.0333 + 0.0004 = 0.0430 \quad (4.30\%) \\
\Delta \lambda_{\text{laser}} &= 648 \times 0.0430 \approx \SI{28}{nm} \\
\lambda_{\text{laser}} &= \SI{648 \pm 28}{nm}
\end{align*}

\textbf{Comparison:} The experimental value $\SI{648 \pm 28}{nm}$ (range $\SI{620}{nm}$ to $\SI{676}{nm}$) agrees well with the typical manufacturer value of $\SI{650}{nm}$.

\textbf{Synthesis Answers}
\begin{enumerate}
    \item \textbf{Double Slit Modification ($a \to 2a$):}
    \begin{itemize}
        \item \textbf{Quantitative:} Fringe spacing $i = \lambda D / a$. If $a$ doubles, $i$ \textbf{halves} ($i' = i/2$). Total width for the same $N$: $W_N' = N i' = N i / 2 = W_N / 2$. The pattern shrinks towards the centre.
        \item \textbf{Qualitative:} The fringes become twice as close together. The central diffraction envelope (governed by single slit width $b$, unchanged) remains the same width, so more interference fringes fit inside the central maximum.
        \item \textbf{Sketch:} Two curves of Intensity $I$ vs Position $y$ on same axes.
    \end{itemize}
    \begin{popfigure}
    \begin{tikzpicture}[scale=0.8, domain=-15:15, samples=400]
        % Axes
        \draw[->] (-16,0) -- (16,0) node[right] {$y$ (mm)};
        \draw[->] (0,-0.5) -- (0,5) node[above] {$I$ (a.u.)};
        % Envelope (single slit diffraction, width b constant)
        \def\envelope{2.5*(sin(180*\x/5)/(180*\x/5))^2} % approximate
        % Case 1: a = a0, fringe spacing i0 = 5.4 mm
        \def\iA{5.4}
        \def\patternA{2.5*(sin(180*\x/5)/(180*\x/5))^2 * (cos(180*\x/\iA))^2}
        % Case 2: a = 2a0, fringe spacing i0/2 = 2.7 mm
        \def\iB{2.7}
        \def\patternB{2.5*(sin(180*\x/5)/(180*\x/5))^2 * (cos(180*\x/\iB))^2}
        % Plot envelope (dashed)
        \draw[dashed, popmuted, thick] plot (\x, {2.5*(sin(180*\x/5)/(180*\x/5))^2});
        % Plot pattern A (original)
        \draw[popblue, very thick] plot (\x, {\patternA});
        % Plot pattern B (doubled a)
        \draw[poporange, very thick, dashed] plot (\x, {\patternB});
        % Labels
        \node[popblue] at (12,4) {$a = a_0$};
        \node[poporange] at (12,3.3) {$a = 2a_0$};
        \node[popmuted] at (-12,4) {Diffraction envelope};
    \end{tikzpicture}
    \legende{Intensity distribution for double slit: original separation $a_0$ (blue solid) and doubled separation $2a_0$ (orange dashed). The single-slit diffraction envelope (grey dashed) is unchanged. Fringe spacing halves when $a$ doubles.}
    \end{popfigure}
    \item \textbf{Doppler Effect (Moving Observer):}
    Formula for stationary source, observer moving towards source:
    \[ f' = f \left( \frac{v + v_o}{v} \right) \]
    Using $v = v_{\text{exp}} = \SI{352}{m/s}$, $v_o = \SI{2.0}{m/s}$, $f = \SI{512}{Hz}$:
    \[ f' = 512 \left( \frac{352 + 2.0}{352} \right) = 512 \times \frac{354}{352} = 512 \times 1.00568 = \SI{514.9}{Hz} \approx \SI{515}{Hz} \]
    The observer hears a frequency \textbf{higher} than the source frequency by about $\SI{3}{Hz}$.
    \item \textbf{Stationary Wave Profile (First Resonance):}
    \begin{popfigure}
    \begin{tikzpicture}[scale=1.2]
        % Tube outline
        \draw[thick] (0,0) -- (0,3) -- (1,3) -- (1,0);
        \draw[thick] (0,0) -- (1,0); % water surface
        % Water
        \fill[popblueL, opacity=0.3] (0,0) rectangle (1,0.3);
        \node at (0.5,0.15) {Water};
        % Displacement wave (quarter sine)
        \draw[popblue, very thick, domain=0:1.57] plot ({0.5+0.4*sin(\x r)}, {0.3+2.7*\x/1.57});
        % Node and Antinode labels
        \node[left] at (0.5,0.3) {N (Node)};
        \node[right] at (0.9,3.0) {A (Antinode)};
        % Distance label
        \draw[<->, popdark] (1.2,0.3) -- (1.2,3.0) node[midway, right] {$L_1+e = \lambda/4$};
        % Arrow for particle motion at antinode
        \draw[<->, poporange, thick] (0.9,2.8) -- (0.9,3.2);
        \node[poporange, right] at (0.9,3.0) {Max amplitude};
        % Phase difference note
        \node[align=center, popdark] at (2.5,1.5) {Phase difference between\\ node and antinode = $\pi/2$ (90°)};
    \end{tikzpicture}
    \legende{Displacement stationary wave in resonance tube at first resonance. Water surface = displacement node (N). Open end ≈ displacement antinode (A). Distance $L_1+e = \lambda/4$.}
    \end{popfigure}
    The phase difference between the motion of air particles at the node (zero displacement) and at the antinode (maximum displacement) is \textbf{$\pi/2$ radians (90°)}. At the node, particles are always at rest; at the antinode, they oscillate with maximum amplitude. In a standing wave, all particles between two adjacent nodes oscillate in phase; the node is a point of zero amplitude, and the phase changes by $\pi$ across a node. Since the antinode is $\lambda/4$ from the node, the phase difference is $\pi/2$.
\end{enumerate}
\end{exoresolu}

\newpage

\coursec{Methods and worked examples}

\courssub{Method 1 — Using the Progressive Wave Equation $y(x,t)$}

\begin{methode}[Reading and exploiting a progressive wave equation]
A one-dimensional progressive harmonic wave is described by
\[
y(x,t)=a\sin\!\left[\omega t - kx + \varphi\right]
\qquad\text{(wave travelling towards increasing }x\text{)}.
\]
To exploit it, proceed step by step:
\begin{enumerate}
\item \textbf{Identify the constants.} Compare with the standard form: $a$ is the \cle{amplitude} (m), $\omega$ the \cle{angular frequency} (rad\,s$^{-1}$), $k$ the \cle{wave number} (rad\,m$^{-1}$), $\varphi$ the \cle{phase constant}.
\item \textbf{Deduce the period and frequency:} $T=\dfrac{2\pi}{\omega}$, $f=\dfrac{1}{T}=\dfrac{\omega}{2\pi}$.
\item \textbf{Deduce the wavelength:} $\lambda=\dfrac{2\pi}{k}$.
\item \textbf{Deduce the speed:} $v=f\lambda=\dfrac{\omega}{k}$.
\item \textbf{Direction of propagation:} sign $-$ before $kx$ $\Rightarrow$ propagation towards $+x$; sign $+$ $\Rightarrow$ towards $-x$.
\item \textbf{Displacement of a given point:} substitute the values of $x$ and $t$ asked for.
\item \textbf{Phase difference} between two points $x_1$ and $x_2$: $\Delta\varphi = k\,|x_2-x_1| = \dfrac{2\pi}{\lambda}|x_2-x_1|$.
\end{enumerate}
\textbf{Reflexes:} always convert wavelengths to metres; check that $\omega/k$ has the unit of a speed; two points separated by $\lambda$ vibrate \emph{in phase}, by $\lambda/2$ \emph{in opposition of phase}.
\end{methode}

\begin{exoresolu}[Exploiting a wave equation (GCE style)]
A transverse wave travels along a rope. Its displacement is given (in SI units) by
\[
y(x,t)=0.04\sin\left(100\pi t - 2\pi x\right).
\]
\begin{enumerate}
\item State the amplitude, the frequency, the wavelength and the speed of the wave.
\item In which direction does the wave travel?
\item Calculate the displacement of the point $x=0.75$ m at time $t=0.02$ s.
\item Find the phase difference between the points $x_1=0.25$ m and $x_2=1.00$ m. What can you say about their motion?
\end{enumerate}
\tcblower
\textbf{1. Constants.} Comparing with $y=a\sin(\omega t - kx)$:
\[
a = 0.04\ \mathrm{m},\qquad \omega = 100\pi\ \mathrm{rad\,s^{-1}},\qquad k = 2\pi\ \mathrm{rad\,m^{-1}}.
\]
Frequency:
\[
f=\frac{\omega}{2\pi}=\frac{100\pi}{2\pi}=50\ \mathrm{Hz}.
\]
Wavelength:
\[
\lambda=\frac{2\pi}{k}=\frac{2\pi}{2\pi}=1.0\ \mathrm{m}.
\]
Speed:
\[
v=f\lambda = 50\times 1.0 = 50\ \mathrm{m\,s^{-1}}.
\]
\textbf{2. Direction.} The sign before $kx$ is negative, so the wave propagates towards \textbf{increasing $x$}.

\textbf{3. Displacement at $x=0.75$ m, $t=0.02$ s.}
\[
y = 0.04\sin\left(100\pi\times 0.02 - 2\pi\times 0.75\right)
=0.04\sin\left(2\pi - 1.5\pi\right)=0.04\sin(0.5\pi).
\]
\[
y = 0.04\times 1 = +0.04\ \mathrm{m}.
\]
The point is at its maximum positive displacement (a crest) at that instant.

\textbf{4. Phase difference.}
\[
\Delta\varphi = k\,|x_2-x_1| = 2\pi(1.00-0.25)=1.5\pi\ \mathrm{rad}.
\]
Since $|x_2-x_1|=0.75\ \mathrm{m}=\dfrac{3\lambda}{4}$, the two points are neither in phase nor in opposition; the point at $x_2$ \emph{lags} behind the point at $x_1$ by $\dfrac{3T}{4}$ because the wave reaches it later.
\end{exoresolu}

\courssub{Method 2 — Graphical Representation: $y(t)$ at a Point and $y(x)$ at an Instant}

\begin{methode}[Distinguishing and using the two graphs of a wave]
\begin{enumerate}
\item \textbf{Identify the graph.} $y$ versus $t$ (fixed $x$) shows the \emph{vibration of one point}: read the period $T$ on the time axis. $y$ versus $x$ (fixed $t$) shows the \emph{shape of the medium} (a ``photograph''): read the wavelength $\lambda$ on the space axis.
\item \textbf{Read the amplitude} as the maximum ordinate in both cases.
\item \textbf{Link the two:} $v=\dfrac{\lambda}{T}$.
\item \textbf{Direction of motion of a point} on a $y(x)$ snapshot: shift the whole profile slightly in the direction of propagation; the point moves towards its new ordinate.
\item \textbf{Delay between two points:} the disturbance at $x_1$ reaches $x_2$ after $\tau=\dfrac{|x_2-x_1|}{v}$; the second point reproduces the motion of the first with this delay.
\end{enumerate}
\textbf{Reflex:} never read a wavelength on a $y(t)$ graph nor a period on a $y(x)$ graph — this is the most frequent error at the GCE.
\end{methode}

\begin{attention}[The two graphs are not interchangeable]
A $y(t)$ curve and a $y(x)$ curve may look identical (both are sinusoids), but their horizontal axes have different meanings and different units (seconds versus metres). Before reading anything, ask yourself: \emph{is this a film of one point, or a photograph of the whole medium?} Reading $T$ on a space graph or $\lambda$ on a time graph costs marks every year.
\end{attention}

\begin{exoresolu}[Reading wave graphs]
A wave propagates along the $x$-axis with speed $v=2.0\ \mathrm{m\,s^{-1}}$. The displacement of the \emph{source point} $O$ ($x=0$) is a sinusoid of amplitude $3$ cm and period $T=0.4$ s, with $y_O(0)=0$ and increasing.
\begin{enumerate}
\item Write the equation $y_O(t)$ of the source.
\item Deduce the equation $y_M(t)$ of a point $M$ at $x=1.2$ m.
\item Describe the profile $y(x)$ of the medium at $t=0.3$ s, stating the wavelength.
\end{enumerate}
\tcblower
\textbf{1. Equation of the source.} $\omega=\dfrac{2\pi}{T}=\dfrac{2\pi}{0.4}=5\pi\ \mathrm{rad\,s^{-1}}$. Since $y_O(0)=0$ and increasing:
\[
y_O(t)=0.03\sin(5\pi t)\quad (\mathrm{m}).
\]
\textbf{2. Equation of $M$.} The delay is
\[
\tau=\frac{x}{v}=\frac{1.2}{2.0}=0.6\ \mathrm{s}=1.5\,T.
\]
Hence
\[
y_M(t)=y_O(t-\tau)=0.03\sin\left[5\pi(t-0.6)\right]=0.03\sin(5\pi t - 3\pi).
\]
Using $\sin(\theta-3\pi)=\sin(\theta-\pi-2\pi)=\sin(\theta-\pi)=-\sin\theta$:
\[
y_M(t)=-0.03\sin(5\pi t).
\]
$M$ vibrates in \textbf{opposition of phase} with $O$. This is consistent with the wavelength found below: $x=1.2\ \mathrm{m}=1.5\lambda$, and a separation of an odd multiple of $\lambda/2$ always gives opposition of phase.

\textbf{3. Profile at $t=0.3$ s.} The wavelength is
\[
\lambda=vT=2.0\times0.4=0.8\ \mathrm{m}.
\]
At $t=0.3$ s the wavefront has reached $x=vt=0.6$ m. A point at position $x$ reproduces the motion of $O$ with the delay $x/v$, so
\[
y(x)=0.03\sin\left[5\pi\left(0.3-\tfrac{x}{2.0}\right)\right]=0.03\sin(1.5\pi - 2.5\pi x),\qquad 0\le x\le 0.6\ \mathrm{m},
\]
and $y=0$ for $x>0.6$ m (the medium there is not yet disturbed). At $x=0$: $y=0.03\sin(1.5\pi)=-0.03$ m — the source is momentarily at a trough. The profile is a sinusoid of wavelength $0.8$ m, truncated at the wavefront $x=0.6$ m.
\end{exoresolu}

\begin{popfigure}
\begin{center}
\begin{tikzpicture}[scale=1.0]
\draw[->,thick] (0,0) -- (7.2,0) node[below left]{$x$ (m)};
\draw[->,thick] (0,-1.6) -- (0,1.8) node[left]{$y$ (cm)};
\draw[popblue,very thick,domain=0:0.6,samples=60] plot (\x*10,{3*sin(deg(1.5*pi - 2.5*pi*\x*10))});
\draw[popblue,very thick] (0.6,0) -- (7.0,0);
\draw[dashed,popmuted] (0.6,0.15) -- (0.6,-0.35) node[below]{$0.6$};
\draw[dashed,popmuted] (0,-1.5) node[left]{$-3$} -- (0.08,-1.5);
\fill[poporange] (0,-1.5) circle (0.07) node[above right=1pt,popdark]{source $O$};
\fill[popgreen] (0.6,0) circle (0.07) node[above right=1pt,popdark]{wavefront};
\end{tikzpicture}
\end{center}
\legende{Profile $y(x)$ at $t=0.3$ s: a sinusoid of wavelength $\lambda=0.8$ m truncated at the wavefront $x=vt=0.6$ m; beyond it the medium is still at rest.}
\end{popfigure}

\courssub{Method 3 — Refraction and Diffraction of Waves}

\begin{methode}[Solving refraction and diffraction problems]
\textbf{Refraction} (change of medium):
\begin{enumerate}
\item The \cle{frequency} $f$ is unchanged when a wave crosses an interface; only $v$ and $\lambda$ change: $\lambda=\dfrac{v}{f}$.
\item For light passing from medium 1 to medium 2: $n_1\sin i_1 = n_2\sin i_2$ (Snell--Descartes), and $n=\dfrac{c}{v}$.
\item Ratio form: $\dfrac{\sin i_1}{\sin i_2}=\dfrac{v_1}{v_2}=\dfrac{\lambda_1}{\lambda_2}=\dfrac{n_2}{n_1}$.
\end{enumerate}
\textbf{Diffraction} (single slit):
\begin{enumerate}
\item Diffraction is appreciable only when the slit width $a$ is \emph{comparable to or smaller than} $\lambda$.
\item Angular position of the first minimum: $\sin\theta_1=\dfrac{\lambda}{a}$; for small angles $\theta_1\approx\dfrac{\lambda}{a}$ (rad).
\item Width of the central bright fringe on a screen at distance $D$: $L=\dfrac{2\lambda D}{a}$.
\end{enumerate}
\textbf{Reflexes:} angles in radians for small-angle approximations; a \emph{narrower} slit gives a \emph{wider} central maximum (inverse proportionality).
\end{methode}

\begin{exoresolu}[Single-slit diffraction of light]
Light from a laser of wavelength $\lambda=633$ nm falls normally on a single slit of width $a=0.10$ mm. The diffraction pattern is observed on a screen placed $D=2.0$ m from the slit.
\begin{enumerate}
\item Calculate the angular half-width $\theta_1$ of the central maximum.
\item Calculate the linear width $L$ of the central bright fringe on the screen.
\item The slit is now replaced by one of width $a'=0.050$ mm. Without further calculation, state how $L$ changes, then verify by computation.
\end{enumerate}
\tcblower
\textbf{1. Angular half-width.} First minimum: $\sin\theta_1=\dfrac{\lambda}{a}$.
\[
\sin\theta_1=\frac{633\times10^{-9}}{0.10\times10^{-3}}=6.33\times10^{-3}
\quad\Rightarrow\quad \theta_1 \approx 6.33\times10^{-3}\ \mathrm{rad}\approx 0.36^{\circ}.
\]
The angle is small, so the approximation $\sin\theta\approx\theta$ is fully justified.

\textbf{2. Linear width of the central fringe.} The first minima lie at $y=\pm D\tan\theta_1\approx D\theta_1$:
\[
L=2D\theta_1=\frac{2\lambda D}{a}=\frac{2\times633\times10^{-9}\times2.0}{0.10\times10^{-3}}
=2.53\times10^{-2}\ \mathrm{m}\approx 2.5\ \mathrm{cm}.
\]

\textbf{3. Halving the slit width.} Since $L\propto \dfrac{1}{a}$, dividing $a$ by 2 must \textbf{double} the width of the central fringe:
\[
L'=\frac{2\lambda D}{a'}=\frac{2\times633\times10^{-9}\times2.0}{0.050\times10^{-3}}
=5.06\times10^{-2}\ \mathrm{m}\approx 5.1\ \mathrm{cm}=2L. \ \checkmark
\]
\textbf{Physical comment:} the narrower the slit, the more the light spreads out — the signature of diffraction. This is why diffraction effects are only easily seen when $a$ is of the order of $\lambda$.
\end{exoresolu}

\courssub{Method 4 — Double-Slit Interference and Measurement of Wavelength}

\begin{methode}[Young's double slits: fringes and wavelength]
\begin{enumerate}
\item \textbf{Path difference} at a point $P$ of the screen: $\delta = a\sin\theta \approx \dfrac{ax}{D}$, where $a$ = slit separation, $D$ = slit--screen distance, $x$ = position on the screen from the central fringe.
\item \textbf{Bright fringes} (constructive): $\delta = m\lambda$, i.e. $x_m=\dfrac{m\lambda D}{a}$, $m\in\mathbb{Z}$.
\item \textbf{Dark fringes} (destructive): $\delta = \left(m+\tfrac12\right)\lambda$.
\item \textbf{Fringe spacing} (distance between two consecutive bright or dark fringes):
\[
i=\frac{\lambda D}{a}.
\]
\item \textbf{Measuring $\lambda$:} measure the total width of $N$ fringe spacings (to reduce error), $i=\dfrac{\text{width}}{N}$, then $\lambda=\dfrac{ia}{D}$.
\end{enumerate}
\textbf{Reflexes:} convert everything to metres; check the order of magnitude ($\lambda$ for visible light lies between $400$ and $750$ nm); for $n$ slits the principal maxima stay at $\delta=m\lambda$ but become sharper and brighter.
\end{methode}

\begin{exoresolu}[Measuring a wavelength with Young's slits]
In a Young's double-slit experiment performed in a school laboratory in Bamenda, the two slits are separated by $a=0.50$ mm and the screen is at $D=1.50$ m. A student measures the distance between the centre of the 1st bright fringe on the left and the centre of the 9th bright fringe on the right of the central fringe: she finds $28.8$ mm.
\begin{enumerate}
\item How many fringe spacings does this distance contain?
\item Deduce the fringe spacing $i$ and the wavelength $\lambda$ of the light used.
\item Where is the 3rd dark fringe located on the screen?
\item State one precaution that improves the accuracy of the measurement.
\end{enumerate}
\tcblower
\textbf{1. Number of spacings.} From the 1st bright fringe on the left ($m=-1$) to the 9th on the right ($m=+9$) there are
\[
N = 9-(-1)=10\ \text{fringe spacings}.
\]
\textbf{2. Fringe spacing and wavelength.}
\[
i=\frac{28.8\ \mathrm{mm}}{10}=2.88\ \mathrm{mm}=2.88\times10^{-3}\ \mathrm{m}.
\]
\[
\lambda=\frac{ia}{D}=\frac{2.88\times10^{-3}\times0.50\times10^{-3}}{1.50}
=9.6\times10^{-7}\ \mathrm{m}=960\ \mathrm{nm}.
\]
This value lies in the near infra-red, outside the visible range ($400$--$750$ nm). In a real laboratory session such a result would be a warning sign: the student should re-check the fringe count and the value of $a$. Here we keep the data as given and the \emph{method} is what matters.

\textbf{3. Third dark fringe.} Dark fringes: $x=\left(m+\tfrac12\right)\dfrac{\lambda D}{a}=\left(m+\tfrac12\right)i$. Taking $m=2$:
\[
x_3 = 2.5\,i = 2.5\times 2.88 = 7.2\ \mathrm{mm}
\]
on either side of the central fringe.

\textbf{4. Precaution.} Measure across as many fringes as possible (rather than a single spacing) and divide; this reduces the relative error on $i$. One can also use a travelling microscope or magnifier to locate fringe centres precisely, and ensure the screen is perpendicular to the beam.
\end{exoresolu}

\begin{attention}[Counting fringes]
The distance between the fringe of order $m_1$ and the fringe of order $m_2$ contains $|m_2-m_1|$ spacings, \emph{not} $|m_2-m_1|+1$. From the central fringe ($m=0$) to the 5th bright fringe there are exactly 5 fringe spacings. Miscounting by one is a classic source of error in wavelength measurements.
\end{attention}

\courssub{Method 5 — The Doppler Effect for Sound in Air}

\begin{methode}[Choosing and applying the Doppler formula]
\begin{enumerate}
\item \textbf{Draw a diagram} showing source $S$, observer $O$, and the direction of motion. Take as positive the direction \emph{from source to observer}.
\item \textbf{Key idea:} a moving \emph{source} changes the wavelength in the medium ($\lambda'=\dfrac{v\mp v_s}{f}$); a moving \emph{observer} changes the rate at which wavefronts are received. The two effects are physically different.
\item \textbf{General formula:}
\[
f' = f\,\frac{v \pm v_o}{v \mp v_s},
\]
where $v$ is the speed of sound in air, $v_o$ the speed of the observer, $v_s$ that of the source.
\item \textbf{Sign rule (learn it physically, not by heart):} motion that \emph{decreases} the source--observer distance must give $f'>f$ (higher pitch); motion that \emph{increases} it must give $f'<f$. Choose the signs so that this is respected:
\begin{itemize}
\item source approaching: denominator $v-v_s$; source receding: $v+v_s$;
\item observer approaching: numerator $v+v_o$; observer receding: $v-v_o$.
\end{itemize}
\item \textbf{Always check} the result: approach $\Rightarrow$ $f'$ larger; recession $\Rightarrow$ $f'$ smaller.
\end{enumerate}
\textbf{Reflexes:} convert speeds from km\,h$^{-1}$ to m\,s$^{-1}$ (divide by 3.6); take $v=340\ \mathrm{m\,s^{-1}}$ unless told otherwise.
\end{methode}

\begin{exoresolu}[An ambulance in Yaoundé]
An ambulance travels along a straight road in Yaoundé at $v_s=72$ km\,h$^{-1}$, sounding a siren of frequency $f=800$ Hz. The speed of sound in air is $v=340\ \mathrm{m\,s^{-1}}$.
\begin{enumerate}
\item A pedestrian stands still on the roadside. Calculate the frequency she hears (a) as the ambulance approaches, (b) after it has passed and moves away.
\item The ambulance is stationary with its siren on, while a cyclist rides \emph{towards} it at $v_o=36$ km\,h$^{-1}$. What frequency does the cyclist hear?
\item Compare the results of 1(a) and 2 for equal speeds and comment.
\end{enumerate}
\tcblower
Convert speeds: $v_s=\dfrac{72}{3.6}=20\ \mathrm{m\,s^{-1}}$; $v_o=\dfrac{36}{3.6}=10\ \mathrm{m\,s^{-1}}$.

\textbf{1(a) Approaching source, stationary observer.} The wavelength in front of the source is compressed: $\lambda'=\dfrac{v-v_s}{f}$, so
\[
f'=\frac{v}{\lambda'}=f\,\frac{v}{v-v_s}=800\times\frac{340}{340-20}
=800\times\frac{340}{320}=850\ \mathrm{Hz}.
\]
Check: approach $\Rightarrow$ higher frequency. \checkmark

\textbf{1(b) Receding source.}
\[
f''=f\,\frac{v}{v+v_s}=800\times\frac{340}{360}=755.6\approx 756\ \mathrm{Hz}.
\]
Check: recession $\Rightarrow$ lower frequency. \checkmark

\textbf{2. Moving observer, stationary source.} The wavefronts are spaced by the normal wavelength $\lambda=\dfrac{v}{f}$, but the cyclist meets them at relative speed $v+v_o$:
\[
f'=\frac{v+v_o}{\lambda}=f\,\frac{v+v_o}{v}=800\times\frac{350}{340}=823.5\approx 824\ \mathrm{Hz}.
\]

\textbf{3. Comparison.} For the same relative speed of approach, a moving source ($850$ Hz) produces a \emph{larger} shift than a moving observer ($824$ Hz). The two situations are physically distinct for sound because the medium (air) fixes a privileged frame: a moving source really changes the wavelength in the air, whereas a moving observer only changes his own rate of crossing wavefronts.
\end{exoresolu}

\courssub{Method 6 — Stationary (Standing) Waves on a String and in Pipes}

\begin{methode}[Analysing stationary waves]
\begin{enumerate}
\item \textbf{Formation:} a stationary wave results from the superposition of two identical progressive waves travelling in opposite directions (incident + reflected): $y=2a\sin(kx)\cos(\omega t)$. All points between two \cle{nodes} vibrate \emph{in phase} with different amplitudes; points on either side of a node vibrate in \emph{opposition of phase}.
\item \textbf{Geometry:} distance node--node (or antinode--antinode) $=\dfrac{\lambda}{2}$; node--antinode $=\dfrac{\lambda}{4}$.
\item \textbf{String fixed at both ends} (length $L$): $L=n\dfrac{\lambda_n}{2}$, so
\[
f_n=\frac{nv}{2L},\qquad n=1,2,3,\dots \quad(\text{all harmonics}).
\]
\item \textbf{Pipe open at both ends:} same formula $f_n=\dfrac{nv}{2L}$ (all harmonics).
\item \textbf{Pipe closed at one end:} $L=(2n-1)\dfrac{\lambda}{4}$, so only \emph{odd} harmonics: $f=\dfrac{v}{4L},\ \dfrac{3v}{4L},\ \dfrac{5v}{4L},\dots$
\item \textbf{Speed on a stretched string:} $v=\sqrt{\dfrac{T}{\mu}}$ ($T$ tension, $\mu$ linear density).
\end{enumerate}
\textbf{Reflexes:} identify the boundary conditions first (fixed end = node, open end = antinode); the fundamental corresponds to the \emph{longest} possible wavelength.
\end{methode}

\begin{popfigure}
\begin{center}
\begin{tikzpicture}[scale=1.0]
\draw[thick] (0,0) -- (8,0);
\fill[poppurple] (0,0) circle (0.08); \fill[poppurple] (4,0) circle (0.08); \fill[poppurple] (8,0) circle (0.08);
\draw[popblue,very thick,domain=0:8,samples=100] plot (\x,{1.1*sin(deg(pi*\x/4))});
\draw[poporange,very thick,domain=0:8,samples=100] plot (\x,{-1.1*sin(deg(pi*\x/4))});
\draw[dashed,popmuted] (2,-1.4) -- (2,1.4); \draw[dashed,popmuted] (6,-1.4) -- (6,1.4);
\node[popdark] at (0,-0.45) {node}; \node[popdark] at (4,-0.45) {node}; \node[popdark] at (8,-0.45) {node};
\node[popdark] at (2,1.65) {antinode}; \node[popdark] at (6,1.65) {antinode};
\draw[<->,popgreen] (0,-1.15) -- (4,-1.15) node[midway,below]{$\lambda/2$};
\end{tikzpicture}
\end{center}
\legende{Second harmonic ($n=2$) on a string fixed at both ends: the two extreme positions of the string are shown. Nodes never move; antinodes vibrate with maximum amplitude $2a$.}
\end{popfigure}

\begin{exoresolu}[Harmonics of a guitar string and of a pipe]
\textbf{A.} A guitar string of length $L=0.65$ m and linear density $\mu=1.0\times10^{-3}$ kg\,m$^{-1}$ is stretched by a tension $T=100$ N.
\begin{enumerate}
\item Calculate the speed of transverse waves on the string.
\item Calculate the frequencies of the fundamental and of the second and third harmonics.
\end{enumerate}
\textbf{B.} A pipe open at both ends has the same fundamental frequency as the string. Take $v_{\text{sound}}=340\ \mathrm{m\,s^{-1}}$.
\begin{enumerate}
\item[3.] Find the length of the pipe.
\item[4.] The pipe is now closed at one end, keeping the same length. What is its new fundamental frequency?
\end{enumerate}
\tcblower
\textbf{1. Speed on the string.}
\[
v=\sqrt{\frac{T}{\mu}}=\sqrt{\frac{100}{1.0\times10^{-3}}}=\sqrt{10^{5}}=316\ \mathrm{m\,s^{-1}}.
\]
\textbf{2. Harmonics of the string.} Fixed--fixed string: $f_n=\dfrac{nv}{2L}$.
\[
f_1=\frac{316}{2\times0.65}=243\ \mathrm{Hz},\qquad
f_2=2f_1=486\ \mathrm{Hz},\qquad
f_3=3f_1=729\ \mathrm{Hz}.
\]
\textbf{3. Length of the open pipe.} Open--open pipe: $f_1=\dfrac{v_{\text{sound}}}{2L'}$, hence
\[
L'=\frac{v_{\text{sound}}}{2f_1}=\frac{340}{2\times243}=0.70\ \mathrm{m}.
\]
\textbf{4. Same pipe closed at one end.} Now $L'=\dfrac{\lambda}{4}$ for the fundamental:
\[
f_1'=\frac{v_{\text{sound}}}{4L'}=\frac{340}{4\times0.70}=121\ \mathrm{Hz}.
\]
Closing one end \textbf{halves} the fundamental frequency (the note drops by one octave) and removes all even harmonics: only $121$ Hz, $364$ Hz, $607$ Hz, \dots\ are produced. This is why a stopped organ pipe sounds deeper and has a different timbre from an open one.
\end{exoresolu}

\courssub{Method 7 — Combined Examination Problem (Synoptic Skill)}

\begin{methode}[Tackling a long GCE structured question on waves]
\begin{enumerate}
\item \textbf{Read the whole question first} and identify the sub-topics (progressive wave, interference, Doppler, stationary waves); the parts are usually independent — if you are stuck on one, move on.
\item \textbf{Write down the data} with symbols and SI units before any calculation.
\item \textbf{Quote the relevant law in symbols} before substituting numbers (marks are awarded for the method).
\item \textbf{Keep 3 significant figures} and carry units through every line.
\item \textbf{Exploit links between parts:} a wavelength measured by interference may reappear in a diffraction or Doppler part.
\item \textbf{Final check:} orders of magnitude (visible light $\sim 10^{-7}$ m; sound in air $\sim 340$ m\,s$^{-1}$; fringe spacings $\sim$ mm).
\end{enumerate}
\end{methode}

\begin{exoresolu}[Synoptic problem: light and sound]
\textbf{Part I.} In a double-slit experiment, slits separated by $a=0.20$ mm are illuminated by monochromatic light; the screen is at $D=1.20$ m. The fringe spacing measured is $i=3.54$ mm.
\begin{enumerate}
\item Calculate the wavelength $\lambda$ of the light.
\item The same light passes through a single slit of width $0.050$ mm placed $1.20$ m from a screen. Calculate the width of the central diffraction maximum.
\end{enumerate}
\textbf{Part II.} A stationary car hooter emits a note of frequency $400$ Hz. A motorcyclist rides away from it at $15$ m\,s$^{-1}$. ($v_{\text{sound}}=340$ m\,s$^{-1}$.)
\begin{enumerate}
\item[3.] Calculate the frequency heard by the motorcyclist.
\end{enumerate}
\tcblower
\textbf{1. Wavelength.}
\[
\lambda=\frac{ia}{D}=\frac{3.54\times10^{-3}\times0.20\times10^{-3}}{1.20}
=5.90\times10^{-7}\ \mathrm{m}=590\ \mathrm{nm}.
\]
This is yellow light (sodium-like), consistent with the visible range. \checkmark

\textbf{2. Central diffraction maximum.}
\[
L=\frac{2\lambda D}{a_{\text{slit}}}
=\frac{2\times5.90\times10^{-7}\times1.20}{0.050\times10^{-3}}
=2.83\times10^{-2}\ \mathrm{m}\approx 2.8\ \mathrm{cm}.
\]
Note how the same $\lambda$ governs both phenomena: interference (two slits) and diffraction (one slit).

\textbf{3. Doppler: observer receding from a stationary source.}
\[
f'=f\,\frac{v-v_o}{v}=400\times\frac{340-15}{340}
=400\times\frac{325}{340}=382\ \mathrm{Hz}.
\]
Check: the distance source--observer increases, so the pitch must drop ($382<400$ Hz). \checkmark

\textbf{Examiner's remark:} each part uses a different formula but the \emph{same working discipline} — data in SI units, symbolic law first, numerical substitution, result with unit and a physical check. This discipline is what earns full method marks at the GCE.
\end{exoresolu}

\newpage

\coursec{Exercises}

\courssub{I check my knowledge}

\begin{exercice}[MCQ: nature and classification of waves]
For each question, choose the correct answer and justify briefly.
\begin{enumerate}
\item A wave transports:
\begin{enumerate}
\item matter only \quad b) energy only \quad c) matter and energy \quad d) neither matter nor energy
\end{enumerate}
\item In a transverse wave, the particles of the medium vibrate:
\begin{enumerate}
\item parallel to the direction of propagation \quad b) perpendicular to the direction of propagation \quad c) in circles \quad d) not at all
\end{enumerate}
\item Which of the following waves does \textbf{not} need a material medium?
\begin{enumerate}
\item sound in air \quad b) waves on a rope \quad c) light \quad d) ripples on water
\end{enumerate}
\item The speed of a progressive wave is given by:
\begin{enumerate}
\item $v = \lambda/T$ \quad b) $v = \lambda T$ \quad c) $v = T/\lambda$ \quad d) $v = \lambda T^2$
\end{enumerate}
\end{enumerate}
\end{exercice}

\begin{exercice}[True or false]
State whether each statement is true or false, and justify your answer in one or two sentences.
\begin{enumerate}
\item When a wave passes from one medium to another, its frequency changes but its wavelength stays the same.
\item Diffraction is more pronounced when the aperture is comparable to or smaller than the wavelength.
\item In a stationary wave, all points of the medium vibrate with the same amplitude.
\item The Doppler effect occurs only when the source is moving.
\end{enumerate}
\end{exercice}

\begin{exercice}[Definitions]
Define each of the following terms in one or two clear sentences, giving one example where appropriate:
\begin{enumerate}
\item a progressive wave;
\item the wavelength of a wave;
\item a node and an antinode of a stationary wave;
\item the fringe spacing in a double-slit interference pattern.
\end{enumerate}
\end{exercice}

\begin{exercice}[Direct calculations]
\begin{enumerate}
\item A wave on a rope has wavelength $\lambda = \SI{0.40}{m}$ and frequency $f = \SI{5.0}{Hz}$. Calculate its speed.
\item A sound wave travels in air at $\SI{340}{m.s^{-1}}$ with frequency $\SI{680}{Hz}$. Calculate its wavelength and its period.
\item A stationary wave on a string fixed at both ends has a distance of $\SI{25}{cm}$ between two consecutive nodes. What is the wavelength?
\end{enumerate}
\end{exercice}

\courssub{I practise}

\begin{exercice}[Equation of a progressive wave]
A progressive wave is given, in SI units, by
\[ y(x,t) = 0.04\sin(100\pi t - 2\pi x). \]
\begin{enumerate}
\item Determine the amplitude, the frequency, the wavelength and the speed of the wave.
\item State the direction of propagation and justify your answer.
\item Sketch the graph of $y$ against $x$ at $t = 0$, marking the amplitude and the wavelength.
\end{enumerate}
\end{exercice}

\begin{exercice}[Refraction of water waves]
Plane water waves of wavelength $\SI{3.0}{cm}$ pass from deep water, where their speed is $\SI{24}{cm.s^{-1}}$, into shallow water, where their speed drops to $\SI{18}{cm.s^{-1}}$. The angle of incidence is $30^\circ$.
\begin{enumerate}
\item Explain why the frequency of the wave does not change at the boundary.
\item Calculate the frequency of the wave and its wavelength in the shallow water.
\item Calculate the angle of refraction.
\end{enumerate}
\end{exercice}

\begin{exercice}[Young's double slits]
In a Young's double-slit experiment, the slits are $\SI{0.50}{mm}$ apart and the screen is placed $\SI{2.0}{m}$ away. It is observed that ten fringe spacings occupy a distance of $\SI{2.4}{cm}$ on the screen.
\begin{enumerate}
\item Calculate the fringe spacing.
\item Deduce the wavelength of the light used, expressing your answer in nanometres.
\item State two changes that would increase the fringe spacing.
\end{enumerate}
\end{exercice}

\begin{exercice}[Stationary waves on a string]
A stretched string of length $\SI{0.80}{m}$, fixed at both ends, vibrates in three loops at a frequency of $\SI{300}{Hz}$.
\begin{enumerate}
\item Sketch the vibration pattern, marking nodes and antinodes.
\item Determine the wavelength of the waves on the string.
\item Calculate the speed of transverse waves on the string.
\item Deduce the fundamental frequency of the string.
\end{enumerate}
\end{exercice}

\courssub{I prepare for the GCE}

\begin{exercice}[Graphs of a progressive wave \hfill (12 marks)]
The figure (to be drawn by the candidate) shows, for the same progressive wave travelling along a rope, the graph of displacement $y$ against position $x$ at time $t = 0$, and the graph of $y$ against time $t$ for the particle at $x = 0$. The first graph is sinusoidal with amplitude $\SI{5.0}{cm}$, the distance between two consecutive crests being $\SI{0.80}{m}$; the second graph is sinusoidal with period $\SI{0.40}{s}$.
\begin{enumerate}
\item Define the terms \emph{wavelength} and \emph{period} of a progressive wave. \hfill (2)
\item From the graphs, deduce the wavelength $\lambda$ and the period $T$ of the wave. \hfill (2)
\item Calculate the speed of the wave along the rope. \hfill (2)
\item Write the equation $y(x,t)$ of the wave, taking $y = 0$ and increasing at $x = 0$, $t = 0$. \hfill (3)
\item Calculate the displacement of the particle at $x = \SI{0.20}{m}$ at time $t = \SI{0.10}{s}$. \hfill (3)
\end{enumerate}
\end{exercice}

\begin{exercice}[Diffraction and interference of light \hfill (14 marks)]
\begin{enumerate}
\item A single slit of width $a = \SI{0.10}{mm}$ is illuminated normally by monochromatic light of wavelength $\lambda = \SI{600}{nm}$.
\begin{enumerate}
\item State the condition giving the angular positions of the minima of the single-slit diffraction pattern. \hfill (2)
\item Calculate the angular width of the central diffraction maximum. \hfill (2)
\item The pattern is observed on a screen $\SI{1.5}{m}$ from the slit. Calculate the linear width of the central bright fringe. \hfill (2)
\end{enumerate}
\item The single slit is replaced by two parallel slits $\SI{0.50}{mm}$ apart, illuminated by the same light, the screen remaining at $\SI{1.5}{m}$.
\begin{enumerate}
\item Define the fringe spacing $i$ and give its expression. \hfill (2)
\item Calculate the fringe spacing. \hfill (2)
\item How many bright fringes are visible within the central diffraction maximum of each slit? \hfill (2)
\end{enumerate}
\item The two slits are now replaced by a diffraction grating with 500 lines per mm. Calculate the angle at which the second-order maximum is observed for the same light. \hfill (2)
\end{enumerate}
\end{exercice}

\begin{exercice}[Doppler effect and hearing through a wall \hfill (14 marks)]
Take the speed of sound in air as $v = \SI{340}{m.s^{-1}}$.
\begin{enumerate}
\item An ambulance emitting a siren of frequency $\SI{500}{Hz}$ moves at $\SI{30}{m.s^{-1}}$ along a straight road past a stationary observer.
\begin{enumerate}
\item State the Doppler formula for a moving source and a stationary observer, defining each symbol. \hfill (2)
\item Calculate the frequency heard by the observer as the ambulance approaches. \hfill (2)
\item Calculate the frequency heard as the ambulance moves away. \hfill (2)
\item Explain briefly why the observed frequency changes abruptly as the ambulance passes. \hfill (2)
\end{enumerate}
\item A motorist now drives at $\SI{20}{m.s^{-1}}$ towards a stationary siren emitting $\SI{800}{Hz}$.
\begin{enumerate}
\item Calculate the frequency heard by the motorist. \hfill (2)
\item Calculate the frequency that would be heard if, instead, the siren moved at $\SI{20}{m.s^{-1}}$ towards a stationary motorist. Comment on the difference. \hfill (2)
\end{enumerate}
\item Using labelled diagrams and typical orders of magnitude of the wavelengths of sound and of light, explain why a person standing behind a wall can hear a conversation taking place on the other side but cannot see the speakers. \hfill (2)
\end{enumerate}
\end{exercice}

\newpage

\coursec{Solutions to the exercises}

\begin{corrige}[Exercise 1 --- MCQ: nature and classification of waves]
\begin{enumerate}
\item \textbf{Answer: b) energy only.} A wave is a disturbance that propagates \emph{through a material medium} (or through vacuum for electromagnetic waves) transporting \cle{energy} from one point to another. The particles of the medium oscillate about their mean positions and return to them; there is no net transport of matter. For example, a cork on water bobs up and down as a ripple passes but is not carried along with the wave.
\item \textbf{Answer: b) perpendicular to the direction of propagation.} By definition, in a \cle{transverse} wave the vibration of the particles is perpendicular to the direction of propagation (e.g.\ a wave on a rope, light). Parallel vibration would describe a \emph{longitudinal} wave such as sound in air.
\item \textbf{Answer: c) light.} Light is an \cle{electromagnetic wave}: it consists of oscillating electric and magnetic fields and propagates through vacuum (it reaches us from the Sun through empty space). Sound, rope waves and water ripples are \emph{mechanical} waves and all require a material medium.
\item \textbf{Answer: a) $v = \lambda/T$.} In one period $T$ the wave advances by exactly one wavelength $\lambda$, so $v = \dfrac{\lambda}{T} = \lambda f$. Dimensional check: $[\lambda/T] = \mathrm{m\,s^{-1}}$, which is indeed a speed; options b), c) and d) do not have the dimensions of a speed.
\end{enumerate}
\end{corrige}

\begin{corrige}[Exercise 2 --- True or false]
\begin{enumerate}
\item \textbf{False.} It is the \cle{frequency} that stays the same when a wave crosses a boundary (it is imposed by the source and the wave cannot ``pile up'' at the interface). Since the speed changes, the wavelength must change according to $\lambda = v/f$: the wave slows down in the denser (slower) medium and its wavelength decreases.
\item \textbf{True.} Diffraction --- the spreading of a wave after passing an aperture or meeting an obstacle --- is significant only when the size $a$ of the aperture is of the order of, or smaller than, the wavelength $\lambda$. The angular spread of the central lobe is of order $\theta \approx \lambda/a$, which is large when $a \lesssim \lambda$.
\item \textbf{False.} In a stationary wave the amplitude varies from point to point: it is \emph{zero} at the \cle{nodes} and \emph{maximum} at the \cle{antinodes}. It is the \emph{phase} that is the same for all points between two consecutive nodes, not the amplitude.
\item \textbf{False.} The Doppler effect occurs whenever there is \emph{relative motion} between source and observer: a moving source with a stationary observer, a moving observer with a stationary source, or both moving. (For sound, the two cases are not even equivalent, because the medium --- air --- provides an absolute frame of reference.)
\end{enumerate}
\end{corrige}

\begin{corrige}[Exercise 3 --- Definitions]
\begin{enumerate}
\item \textbf{Progressive wave:} a disturbance that travels through a \emph{material medium} in a given direction, each point of the medium reproducing, with a time delay, the motion of the source, and transporting energy without net transport of matter. \emph{Example:} a pulse sent along a rope, ripples spreading on a pond, sound in air.
\item \textbf{Wavelength:} the wavelength $\lambda$ is the smallest distance, measured in the direction of propagation, between two points of the medium that vibrate \emph{in phase} (e.g.\ two consecutive crests or two consecutive compressions); it is also the distance travelled by the wave during one period $T$, so $\lambda = vT = v/f$.
\item \textbf{Node and antinode:} in a stationary wave, a \cle{node} is a point of the medium that remains permanently at rest (zero amplitude), while an \cle{antinode} is a point that vibrates with maximum amplitude. Consecutive nodes (or consecutive antinodes) are separated by $\lambda/2$.
\item \textbf{Fringe spacing:} in a double-slit interference pattern, the fringe spacing $i$ is the distance, measured on the screen, between the centres of two consecutive bright fringes (or two consecutive dark fringes). For slits separated by $a$ with a screen at distance $D$, $i = \dfrac{\lambda D}{a}$.
\end{enumerate}
\end{corrige}

\begin{corrige}[Exercise 4 --- Direct calculations]
\begin{enumerate}
\item Speed of the wave on the rope:
\[ v = \lambda f = 0.40 \times 5.0 = \SI{2.0}{m.s^{-1}}. \]
\item Wavelength and period of the sound wave:
\[ \lambda = \frac{v}{f} = \frac{340}{680} = \SI{0.50}{m}; \qquad T = \frac{1}{f} = \frac{1}{680} = \SI{1.47e-3}{s} \approx \SI{1.5}{ms}. \]
\item In a stationary wave, two consecutive nodes are separated by half a wavelength:
\[ \frac{\lambda}{2} = \SI{25}{cm} \quad\Longrightarrow\quad \lambda = 2 \times 25 = \SI{50}{cm} = \SI{0.50}{m}. \]
\end{enumerate}
\end{corrige}

\begin{corrige}[Exercise 5 --- Equation of a progressive wave]
The general equation of a sinusoidal progressive wave travelling in the $+x$ direction is
\[ y(x,t) = A\sin(\omega t - kx + \varphi), \qquad \omega = \frac{2\pi}{T} = 2\pi f, \quad k = \frac{2\pi}{\lambda}. \]
Comparing with $y(x,t) = 0.04\sin(100\pi t - 2\pi x)$:
\begin{enumerate}
\item \textbf{Amplitude:} $A = \SI{0.04}{m} = \SI{4}{cm}$.\par
\textbf{Frequency:} $\omega = 100\pi\ \mathrm{rad\,s^{-1}} \Rightarrow f = \dfrac{\omega}{2\pi} = \dfrac{100\pi}{2\pi} = \SI{50}{Hz}$.\par
\textbf{Wavelength:} $k = 2\pi\ \mathrm{m^{-1}} \Rightarrow \lambda = \dfrac{2\pi}{k} = \dfrac{2\pi}{2\pi} = \SI{1.0}{m}$.\par
\textbf{Speed:} $v = \lambda f = 1.0 \times 50 = \SI{50}{m.s^{-1}}$ (equivalently $v = \omega/k = 100\pi/2\pi$).
\item \textbf{Direction of propagation:} the argument has the form $(\omega t - kx)$, i.e.\ the signs of the $t$ and $x$ terms are \emph{opposite}. The wave therefore propagates in the direction of \textbf{increasing $x$} (the $+x$ direction). Indeed, a given phase is maintained when $x$ increases as $t$ increases: $x = (\omega/k)\,t + \text{const}$.
\item At $t = 0$: $y(x,0) = 0.04\sin(-2\pi x) = -0.04\sin(2\pi x)$, a sinusoid of amplitude $\SI{0.04}{m}$ and period (in $x$) $\lambda = \SI{1.0}{m}$, passing through the origin with negative slope.
\begin{popfigure}
\begin{center}
\begin{tikzpicture}[scale=1.0]
\draw[->,thick] (-0.2,0) -- (2.6,0) node[below] {$x$ (m)};
\draw[->,thick] (0,-1.5) -- (0,1.7) node[left] {$y$ (m)};
\draw[popblue,thick,domain=0:2.3,samples=200] plot (\x,{-1.1*sin(360*\x)});
\draw[dashed,popmuted] (-0.1,-1.1) -- (2.4,-1.1) node[right,popink] {$-A$};
\draw[dashed,popmuted] (-0.1,1.1) -- (2.4,1.1) node[right,popink] {$+A$};
\draw[<->,poporange,thick] (0.25,1.35) -- (1.25,1.35) node[midway,above] {$\lambda = 1.0$ m};
\foreach \x in {0,0.5,1,1.5,2} \draw (\x,0.06) -- (\x,-0.06) node[below] {\small \x};
\end{tikzpicture}
\end{center}
\legende{Graph of $y$ against $x$ at $t = 0$: amplitude $A = \SI{0.04}{m}$, wavelength $\lambda = \SI{1.0}{m}$.}
\end{popfigure}
\end{enumerate}
\end{corrige}

\begin{corrige}[Exercise 6 --- Refraction of water waves]
\begin{enumerate}
\item \textbf{Why the frequency is unchanged:} the frequency is fixed by the source (the vibrating dipper). At the boundary, the wave cannot accumulate or disappear: each crest arriving at the boundary in deep water must produce exactly one crest in the shallow water, otherwise crests would pile up indefinitely at the interface. Hence the number of vibrations per second is the same on both sides: $f$ is conserved. What changes is the speed, and therefore the wavelength via $\lambda = v/f$.
\item \textbf{Frequency} (computed in deep water, valid everywhere):
\[ f = \frac{v_1}{\lambda_1} = \frac{24}{3.0} = \SI{8.0}{Hz}. \]
\textbf{Wavelength in shallow water:}
\[ \lambda_2 = \frac{v_2}{f} = \frac{18}{8.0} = \SI{2.25}{cm} \approx \SI{2.3}{cm}. \]
\item \textbf{Angle of refraction.} The law of refraction for waves (Snell--Descartes) gives
\[ \frac{\sin i_1}{v_1} = \frac{\sin i_2}{v_2}
\quad\Longrightarrow\quad
\sin i_2 = \frac{v_2}{v_1}\sin i_1 = \frac{18}{24}\sin 30^\circ = 0.75 \times 0.500 = 0.375. \]
\[ i_2 = \arcsin(0.375) \approx 22^\circ. \]
The wave slows down and bends \emph{towards the normal}, as expected when entering a slower medium.
\end{enumerate}
\end{corrige}

\begin{corrige}[Exercise 7 --- Young's double slits]
\begin{enumerate}
\item \textbf{Fringe spacing:} ten fringe spacings occupy $\SI{2.4}{cm}$, so
\[ i = \frac{2.4}{10} = \SI{0.24}{cm} = \SI{2.4e-3}{m}. \]
\item \textbf{Wavelength:} from $i = \dfrac{\lambda D}{a}$ we obtain
\[ \lambda = \frac{ia}{D} = \frac{2.4\times10^{-3}\times 0.50\times10^{-3}}{2.0}
= 6.0\times10^{-7}\ \mathrm{m} = \SI{600}{nm}. \]
\item \textbf{To increase $i = \lambda D/a$, one can} (any two):
\begin{itemize}
\item increase the slit--screen distance $D$ (move the screen further away);
\item decrease the slit separation $a$ (use slits closer together);
\item use light of longer wavelength $\lambda$ (e.g.\ replace green light by red light).
\end{itemize}
\end{enumerate}
\end{corrige}

\begin{corrige}[Exercise 8 --- Stationary waves on a string]
\begin{enumerate}
\item \textbf{Vibration pattern:} three loops means three half-wavelengths between the fixed ends, with 4 nodes (including the ends) and 3 antinodes.
\begin{popfigure}
\begin{center}
\begin{tikzpicture}[scale=1.0]
\draw[->,thick] (-0.2,0) -- (8.6,0) node[below] {$x$};
\draw[popblue,thick,domain=0:8,samples=300] plot (\x,{1.1*sin(67.5*\x)});
\draw[popblue,dashed,domain=0:8,samples=300] plot (\x,{-1.1*sin(67.5*\x)});
\foreach \x in {0,2.667,5.333,8} {\fill[popdark] (\x,0) circle (0.09);}
\node[below,popdark] at (0,-0.15) {N};
\node[below,popdark] at (2.667,-0.15) {N};
\node[below,popdark] at (5.333,-0.15) {N};
\node[below,popdark] at (8,-0.15) {N};
\node[above,poporange] at (1.333,1.25) {A};
\node[below,poporange] at (4,-1.25) {A};
\node[above,poporange] at (6.667,1.25) {A};
\draw[<->,popmuted] (0,-1.7) -- (8,-1.7) node[midway,below] {$L = \SI{0.80}{m}$};
\end{tikzpicture}
\end{center}
\legende{String vibrating in three loops ($n=3$): nodes (N) at the fixed ends and at $L/3$, $2L/3$; antinodes (A) at the centre of each loop. The dashed curve shows the string half a period later.}
\end{popfigure}
\item \textbf{Wavelength:} three loops fit in the length $L$, and each loop is $\lambda/2$:
\[ L = 3\,\frac{\lambda}{2} \quad\Longrightarrow\quad \lambda = \frac{2L}{3} = \frac{2\times 0.80}{3} = \SI{0.533}{m} \approx \SI{0.53}{m}. \]
\item \textbf{Speed of the waves:}
\[ v = \lambda f = 0.533 \times 300 = \SI{160}{m.s^{-1}}. \]
\item \textbf{Fundamental frequency:} the fundamental mode is a single loop, $\lambda_1 = 2L$, so
\[ f_1 = \frac{v}{2L} = \frac{160}{2\times 0.80} = \SI{100}{Hz}. \]
(Equivalently, the third harmonic is $3f_1 = \SI{300}{Hz}$, giving $f_1 = \SI{100}{Hz}$ directly.)
\end{enumerate}
\end{corrige}

\begin{corrige}[Exercise 9 --- Graphs of a progressive wave \hfill (12 marks)]
\begin{enumerate}
\item \textbf{Definitions.} \emph{Wavelength} $\lambda$: the shortest distance between two points of the medium vibrating in phase; equivalently, the distance travelled by the wave during one period. \hfill (1)\par
\emph{Period} $T$: the duration of one complete vibration of any point of the medium; it is also the time taken by the wave to travel one wavelength. \hfill (1)
\item \textbf{Reading the graphs.} On the $y(x)$ graph, the distance between two consecutive crests is the wavelength:
\[ \lambda = \SI{0.80}{m}. \hfill (1) \]
On the $y(t)$ graph, the duration of one complete oscillation is the period:
\[ T = \SI{0.40}{s}. \hfill (1) \]
\item \textbf{Speed of the wave:}
\[ v = \frac{\lambda}{T} = \frac{0.80}{0.40} = \SI{2.0}{m.s^{-1}}. \hfill (2) \]
(1 mark for the relation $v = \lambda/T$, 1 mark for the numerical value with unit.)
\item \textbf{Equation of the wave.} Amplitude $A = \SI{5.0}{cm} = \SI{0.050}{m}$;
\[ \omega = \frac{2\pi}{T} = \frac{2\pi}{0.40} = 5\pi\ \mathrm{rad\,s^{-1}}, \qquad
k = \frac{2\pi}{\lambda} = \frac{2\pi}{0.80} = 2.5\pi\ \mathrm{m^{-1}}. \hfill (2) \]
The condition $y = 0$ and increasing at $x = 0$, $t = 0$ imposes zero initial phase, so
\[ \boxed{\,y(x,t) = 0.050\sin(5\pi t - 2.5\pi x)\,} \quad (\text{SI units}). \hfill (1) \]
\item \textbf{Displacement at $x = \SI{0.20}{m}$, $t = \SI{0.10}{s}$:}
\[ y = 0.050\sin\!\big(5\pi\times 0.10 - 2.5\pi\times 0.20\big)
= 0.050\sin\!\big(0.5\pi - 0.5\pi\big) = 0.050\sin 0 = \SI{0}{m}. \hfill (3) \]
(1 mark substitution, 1 mark correct phase $= 0$, 1 mark result. Note: $x = 0.20$ m $= \lambda/4$ and $t = 0.10$ s $= T/4$, so the crest that was at $x=0$ at $t=0$ has just arrived --- the particle passes through zero going up.)
\end{enumerate}
\textbf{Examiner's expectations:} definitions must mention ``in phase'' (for $\lambda$) and ``one complete vibration'' (for $T$); the equation must be dimensionally consistent with the phase in radians; all numerical answers must carry a unit.
\end{corrige}

\begin{corrige}[Exercise 10 --- Diffraction and interference of light \hfill (14 marks)]
\begin{enumerate}
\item \textbf{Single slit}, $a = \SI{0.10}{mm} = \SI{1.0e-4}{m}$, $\lambda = \SI{600}{nm} = \SI{6.0e-7}{m}$.
\begin{enumerate}
\item \textbf{Condition for the minima:} destructive interference between secondary sources across the slit occurs when
\[ a\sin\theta_n = n\lambda, \qquad n = 1, 2, 3, \dots \hfill (2) \]
($\theta_n$ measured from the central axis; $n = 0$ is excluded since it gives the central maximum.)
\item \textbf{Angular width of the central maximum:} the first minimum is at
\[ \sin\theta_1 = \frac{\lambda}{a} = \frac{6.0\times10^{-7}}{1.0\times10^{-4}} = 6.0\times10^{-3}
\;\Rightarrow\; \theta_1 \approx 6.0\times10^{-3}\ \mathrm{rad}. \]
The central maximum extends from $-\theta_1$ to $+\theta_1$:
\[ \Delta\theta = 2\theta_1 = \frac{2\lambda}{a} = \SI{1.2e-2}{rad} \approx 0.69^\circ. \hfill (2) \]
\item \textbf{Linear width on the screen} at $D = \SI{1.5}{m}$ (small-angle approximation $\tan\theta \approx \theta$):
\[ L = 2D\tan\theta_1 \approx \frac{2\lambda D}{a}
= \frac{2\times 6.0\times10^{-7}\times 1.5}{1.0\times10^{-4}}
= \SI{1.8e-2}{m} = \SI{18}{mm}. \hfill (2) \]
\end{enumerate}
\item \textbf{Double slits}, slit separation $b = \SI{0.50}{mm} = \SI{5.0e-4}{m}$.
\begin{enumerate}
\item \textbf{Definition and expression:} the fringe spacing $i$ is the distance on the screen between the centres of two consecutive bright (or dark) fringes:
\[ i = \frac{\lambda D}{b}. \hfill (2) \]
\item \textbf{Value:}
\[ i = \frac{6.0\times10^{-7}\times 1.5}{5.0\times10^{-4}} = \SI{1.8e-3}{m} = \SI{1.8}{mm}. \hfill (2) \]
\item \textbf{Bright fringes inside the central diffraction maximum:} the central diffraction lobe has half-width $\SI{9}{mm}$ on the screen (from part a), i.e.\ it spans $\SI{18}{mm}$. Bright fringes lie at $x = ni$ ($n$ integer) and are visible while $|x| < \SI{9}{mm}$:
\[ |n| < \frac{9}{1.8} = 5 \quad\Rightarrow\quad n = 0, \pm1, \pm2, \pm3, \pm4, \]
the fringes $n = \pm5$ falling exactly on the diffraction minima (missing orders). Hence $\boxed{9}$ bright fringes are visible. \hfill (2)
\end{enumerate}
\item \textbf{Diffraction grating:} 500 lines per mm gives a grating spacing
\[ d = \frac{10^{-3}}{500} = \SI{2.0e-6}{m}. \]
Maxima satisfy $d\sin\theta_n = n\lambda$. For the second order ($n = 2$):
\[ \sin\theta_2 = \frac{2\lambda}{d} = \frac{2\times 6.0\times10^{-7}}{2.0\times10^{-6}} = 0.60
\quad\Rightarrow\quad \theta_2 \approx 36.9^\circ. \hfill (2) \]
\end{enumerate}
\textbf{Examiner's expectations:} correct conversion of nm and mm to metres; clear distinction between slit \emph{width} (diffraction) and slit \emph{separation} (interference); in (b)(iii) the candidate must explain the counting, including the central fringe and the missing orders at $n = \pm 5$.
\end{corrige}

\begin{corrige}[Exercise 11 --- Doppler effect and hearing through a wall \hfill (14 marks)]
\begin{enumerate}
\item \textbf{Moving source, stationary observer.}
\begin{enumerate}
\item \textbf{Formula:}
\[ f' = \frac{v}{v \mp v_s}\,f, \]
where $f'$ is the observed frequency, $f$ the source frequency, $v$ the speed of sound in air and $v_s$ the speed of the source; the minus sign applies when the source \emph{approaches} and the plus sign when it \emph{recedes}. \hfill (2)
\item \textbf{Approaching ambulance:}
\[ f' = \frac{340}{340 - 30}\times 500 = \frac{340}{310}\times 500 \approx \SI{548}{Hz}. \hfill (2) \]
\item \textbf{Receding ambulance:}
\[ f'' = \frac{340}{340 + 30}\times 500 = \frac{340}{370}\times 500 \approx \SI{459}{Hz}. \hfill (2) \]
\item \textbf{Why the change is abrupt:} while approaching, the source ``chases'' its own wavefronts, compressing them (shorter wavelength, higher pitch); while receding, it moves away from them, stretching them (longer wavelength, lower pitch). At the instant of passing, the motion relative to the observer switches from approach to recession, so the frequency drops suddenly from $\SI{548}{Hz}$ to $\SI{459}{Hz}$ --- the familiar ``eeee--ooo'' of a passing siren. \hfill (2)
\end{enumerate}
\item \textbf{Moving observer, stationary source.}
\begin{enumerate}
\item \textbf{Motorist approaching the siren:} the observer crosses wavefronts more rapidly:
\[ f' = \frac{v + v_o}{v}\,f = \frac{340 + 20}{340}\times 800 = \frac{360}{340}\times 800 \approx \SI{847}{Hz}. \hfill (2) \]
\item \textbf{Siren moving towards the stationary motorist:}
\[ f'' = \frac{v}{v - v_s}\,f = \frac{340}{340 - 20}\times 800 = \frac{340}{320}\times 800 = \SI{850}{Hz}. \hfill (1) \]
\textbf{Comment:} the two results differ ($847 \neq 850$ Hz). For sound, the air is a privileged reference frame: a moving source changes the \emph{wavelength} in the medium, whereas a moving observer changes the \emph{rate at which wavefronts are met}. The physical situations are therefore not symmetric. \hfill (1)
\end{enumerate}
\item \textbf{Hearing but not seeing through a wall.} Audible sound has wavelengths of order $\lambda_{\text{sound}} \sim \SI{0.2}{m}$ to $\SI{2}{m}$ (e.g.\ $\lambda = 340/500 \approx \SI{0.7}{m}$), comparable to the size of doors, windows and gaps around a wall. Sound is therefore strongly \emph{diffracted} and bends around obstacles and through openings, reaching the listener. Visible light has $\lambda_{\text{light}} \sim \SI{5e-7}{m}$, millions of times smaller than any aperture: diffraction is negligible ($\theta \approx \lambda/a \sim 10^{-6}$ rad), so light travels essentially in straight lines and is blocked by the opaque wall.
\begin{popfigure}
\begin{center}
\begin{tikzpicture}[scale=0.9]
\fill[gray!60] (2.4,-2) rectangle (2.7,-0.4);
\fill[gray!60] (2.4,0.4) rectangle (2.7,2);
\node[gray!40,rotate=90] at (2.55,1.3) {\small wall};
\fill[popdark] (0.5,0) circle (0.08) node[below left=1pt,popink] {\small speaker};
\foreach \r in {0.5,1.0,1.5} \draw[popblue] (2.55,0) ++(180:\r) arc (180:150:\r);
\foreach \r in {0.5,1.0,1.5} \draw[popblue] (2.55,0) ++(180:\r) arc (180:210:\r);
\foreach \r in {0.5,1.0,1.5,2.0} {\draw[popblue] (2.55,0) ++(40:\r) arc (40:-40:\r);}
\fill[popdark] (5.2,1.2) circle (0.08) node[right,popink] {\small listener};
\draw[poporange,thick,->] (0.5,0) -- (2.5,0);
\node[popblue] at (4.2,-1.6) {\small sound diffracts through the gap};
\draw[poporange,thick,dashed] (0.5,0.9) -- (2.4,1.3);
\node[poporange] at (1.2,1.6) {\small light blocked};
\end{tikzpicture}
\end{center}
\legende{Sound (long $\lambda$) diffracts through the doorway and around the wall; light (very short $\lambda$) propagates straight and is stopped.}
\end{popfigure}
\hfill (2)
\end{enumerate}
\textbf{Examiner's expectations:} the Doppler formulas must be stated with the sign convention clearly linked to approach/recession; numerical answers to 3 significant figures with units; in (c) the argument must explicitly compare the wavelength with the size of the aperture/obstacle --- this comparison, not the mere statement ``sound diffracts'', earns the marks.
\end{corrige}

\newpage

\coursec{Summary sheet}

\begin{popfigure}
\begin{center}\tcbox[colback=popink,colframe=popink,coltext=white,arc=9pt,boxsep=4pt,left=6pt,right=6pt]{\bfseries\small Mechanical Waves}\end{center}
\vspace{-2pt}
\begin{tcbraster}[raster columns=3,raster equal height=rows,raster column skip=6pt,raster row skip=6pt,enhanced,blankest]
\begin{tcolorbox}[enhanced,colback=popblue,colframe=popblue,coltext=white,boxrule=0.9pt,arc=3pt,left=4pt,right=4pt,top=3pt,bottom=3pt,boxsep=1pt,halign=left]\bfseries\scriptsize Nature \& Classification\end{tcolorbox}
\begin{tcolorbox}[enhanced,colback=popgreen,colframe=popgreen,coltext=white,boxrule=0.9pt,arc=3pt,left=4pt,right=4pt,top=3pt,bottom=3pt,boxsep=1pt,halign=left]\bfseries\scriptsize Progressive Wave $y(x,t)$\end{tcolorbox}
\begin{tcolorbox}[enhanced,colback=poporange,colframe=poporange,coltext=white,boxrule=0.9pt,arc=3pt,left=4pt,right=4pt,top=3pt,bottom=3pt,boxsep=1pt,halign=left]\bfseries\scriptsize Reflection, Refraction, Diffraction\end{tcolorbox}
\begin{tcolorbox}[enhanced,colback=poppurple,colframe=poppurple,coltext=white,boxrule=0.9pt,arc=3pt,left=4pt,right=4pt,top=3pt,bottom=3pt,boxsep=1pt,halign=left]\bfseries\scriptsize Interference \& $\lambda$ measurement\end{tcolorbox}
\begin{tcolorbox}[enhanced,colback=poppink,colframe=poppink,coltext=white,boxrule=0.9pt,arc=3pt,left=4pt,right=4pt,top=3pt,bottom=3pt,boxsep=1pt,halign=left]\bfseries\scriptsize Doppler Effect\end{tcolorbox}
\begin{tcolorbox}[enhanced,colback=popgold,colframe=popgold,coltext=white,boxrule=0.9pt,arc=3pt,left=4pt,right=4pt,top=3pt,bottom=3pt,boxsep=1pt,halign=left]\bfseries\scriptsize Stationary Waves\end{tcolorbox}
\end{tcbraster}

\legende{Mind map of the chapter: six branches from the central concept of mechanical waves.}
\end{popfigure}

\begin{bilan}
\textbf{Key definitions}
\begin{itemize}
\item \cle{Wave}: propagation of a disturbance (energy) through a medium \emph{without} net transport of matter.
\item \cle{Transverse wave}: vibration $\perp$ propagation (rope, light, water ripples). \cle{Longitudinal wave}: vibration $\parallel$ propagation (sound in air, spring compressions).
\item \cle{Mechanical wave}: needs a material medium (sound, water waves); electromagnetic waves travel in vacuum.
\item Characteristics: period $T$ (s), frequency $f=1/T$ (Hz), wavelength $\lambda$ (m), amplitude $A$, speed $v=\lambda f=\lambda/T$.
\item \cle{Progressive wave}: disturbance travelling away from the source; \cle{stationary wave}: no net propagation, fixed nodes and antinodes.
\end{itemize}

\textbf{Laws and formulas to know}
\begin{itemize}
\item Progressive wave (motion along $+x$): $y(x,t)=A\sin(\omega t-kx+\varphi)$, with $\omega=2\pi f$, $k=2\pi/\lambda$, $v=\omega/k$. Sign $+$ in $(\omega t+kx)$ for motion along $-x$.
\item Phase difference between two points: $\Delta\varphi=\dfrac{2\pi\,\Delta x}{\lambda}$; in time: $\Delta\varphi=\dfrac{2\pi\,\Delta t}{T}$.
\item Refraction (Snell): $n_1\sin i_1=n_2\sin i_2$; $v=\dfrac{c}{n}$; $\lambda$ changes, $f$ stays constant.
\item Single-slit diffraction: first minimum at $\sin\theta=\dfrac{\lambda}{a}$; angular half-width $\theta\approx\dfrac{\lambda}{a}$ (small angles).
\item Double slit (Young): fringe spacing $i=\dfrac{\lambda D}{a}$; bright fringes: path difference $\delta=k\lambda$; dark fringes: $\delta=(2k+1)\dfrac{\lambda}{2}$.
\item Diffraction grating ($N$ lines/mm): $d\sin\theta_k=k\lambda$, $d=1/N$.
\item Doppler effect (sound): moving source, stationary observer: $f'=\dfrac{v}{v\mp v_s}f$ ($-$ approaching, $+$ receding); moving observer, stationary source: $f'=\dfrac{v\pm v_o}{v}f$.
\item Stationary wave on a string fixed at both ends: $L=n\dfrac{\lambda}{2}$, $f_n=\dfrac{nv}{2L}$; distance node--node $=\lambda/2$, node--antinode $=\lambda/4$.
\end{itemize}

\textbf{Methods}
\begin{enumerate}
\item To write a wave equation: identify $A$, compute $\omega=2\pi f$ and $k=2\pi/\lambda$, choose the sign from the direction of travel, fix $\varphi$ from initial conditions.
\item To measure $\lambda$ with Young's slits: measure $D$, $a$, and the width of $m$ fringes ($mi$), then $\lambda=\dfrac{ai}{D}$.
\item For Doppler problems: draw the situation, decide who moves, choose the sign by asking ``does the observed frequency go up (approach) or down (recession)?'', then substitute.
\item For stationary waves: count loops between fixed ends, $L=n\lambda/2$; sketch nodes (zero amplitude) and antinodes (maximum amplitude).
\end{enumerate}

\textbf{Common mistakes}
\begin{itemize}
\item Confusing the two graphs: $y$ vs $t$ (one point, period $T$) and $y$ vs $x$ (snapshot, wavelength $\lambda$).
\item Mixing up $\omega$ and $f$, or $k$ and $\lambda$: always convert first.
\item Thinking frequency changes on refraction --- only $v$ and $\lambda$ change.
\item Writing fringe spacing as $\lambda a/D$ instead of $\lambda D/a$ (check dimensions!).
\item Forgetting that in Doppler formulas $v$ is the speed of \emph{sound in air}, not of the source.
\item Believing particles travel with the wave: only energy propagates.
\item Placing nodes at positions of maximum amplitude in stationary waves.
\end{itemize}

\textbf{GCE tips}
\begin{itemize}
\item Always quote units (Hz, m, rad\,s$^{-1}$, rad\,m$^{-1}$) and give 2--3 significant figures.
\item In structured questions, define every symbol you introduce ($D$, $a$, $i$\ldots) and state the condition (small angles, coherent sources).
\item For interference, state explicitly that sources must be \cle{coherent} (same frequency, constant phase difference).
\item Sketch diagrams: a labelled snapshot of a wave or a standing-wave pattern often earns easy marks.
\item Show the formula \emph{before} substituting numbers --- method marks are awarded even if the arithmetic slips.
\end{itemize}
\end{bilan}

\findecours

\end{document}
